\documentclass[10pt,twocolumn]{extarticle}
% 调整了top和bottom留出页眉页脚的空间，防止页码被截断（收紧边距以控制总页数）
\usepackage[a4paper, left=0.7cm, right=0.7cm, top=1.1cm, bottom=1.1cm, includehead]{geometry}
\usepackage{xeCJK}
\usepackage{titlesec}
\usepackage{enumitem}
\usepackage{xcolor}
\usepackage{listings}
\usepackage{booktabs}
\usepackage{longtable}
\usepackage{array}
\usepackage{fancyhdr}
% 加在 \usepackage{fancyhdr} 之后
\usepackage[hidelinks]{hyperref}
% 字体设置，可根据你的系统替换（如 Microsoft YaHei 或 SimSun）
\setCJKmainfont{SimSun}
\setmainfont{Times New Roman}

% 关闭自动编号，完美解决序号重复问题（只保留文本原有的编号）
\setcounter{secnumdepth}{0}

% 详细页码设置（在页面上方/前面显示）
\pagestyle{fancy}
\fancyhf{}
\fancyhead[C]{\small 第 \thepage\ 页}
\renewcommand{\headrulewidth}{0.4pt}
\renewcommand{\footrulewidth}{0pt}

% 紧凑的标题间距设置（进一步收紧以控制页数）
\titlespacing*{\section}{0pt}{0.9ex plus .1ex}{0.3ex plus .1ex}
\titlespacing*{\subsection}{0pt}{0.7ex plus .1ex}{0.2ex plus .1ex}
\titlespacing*{\subsubsection}{0pt}{0.7ex plus .1ex}{0.2ex plus .1ex}
\setlength{\parindent}{0pt}
\setlength{\parskip}{0.3em}

% 列表紧凑设置
\setlist{nosep, leftmargin=1.2em}

% 代码高亮颜色定义
\definecolor{codegreen}{rgb}{0,0.6,0}
\definecolor{codegray}{rgb}{0.5,0.5,0.5}
\definecolor{codepurple}{rgb}{0.58,0,0.82}
\definecolor{backcolour}{rgb}{0.96,0.96,0.94}

% listings 代码块设置
\lstdefinestyle{mystyle}{
    backgroundcolor=\color{backcolour},
    commentstyle=\color{codegreen},
    keywordstyle=\color{blue},
    numberstyle=\tiny\color{codegray},
    stringstyle=\color{codepurple},
    basicstyle=\ttfamily\small,
    breakatwhitespace=false,
    breaklines=true,
    captionpos=b,
    keepspaces=true,
    showspaces=false,
    showstringspaces=false,
    showtabs=false,
    tabsize=4,
    xleftmargin=0pt,
    aboveskip=0.25em,
    belowskip=0.25em
}
\lstset{style=mystyle}

\begin{document}

% ===== 目录区 =====
\pagenumbering{Roman} % 目录页码设为大写罗马数字 (I, II, III)
\setcounter{page}{1}  % 强制目录从第 I 页开始算，杀掉旧缓存
\setcounter{tocdepth}{2} 

\tableofcontents

\clearpage % 必须先换页！把目录彻底排版完再重置正文

% ===== 正文区 =====
\pagenumbering{arabic} % 正文页码设为阿拉伯数字 (1, 2, 3)
\setcounter{page}{1}   % 强制正文从第 1 页重新开始算

% ↓↓↓ 下面直接紧跟你的 \section{第一章 vector（重点）} ↓↓↓
明白你的需求。你要的是\textbf{上机考试速查手册}，不是教材式讲解，所以我会按照：

\begin{quote}
\textbf{① 名称（函数/知识点） \\
② 用法（参数怎么写 + 示例代码） \\
③ 作用（考试时知道什么时候用）}
\end{quote}

来整理。

由于内容非常庞大（尤其 \texttt{vector} + \texttt{string} + \texttt{algorithm} + \texttt{map/set/queue/stack}），一次回复放不下完整版本。我会分多次连续发送。

整体目录如下：

\section*{C++ STL 上机考试速查大全}

\subsection*{第一部分：vector（极详细）}
\begin{itemize}
    \item vector定义
    \item 初始化方式
    \item 遍历方式
    \item 所有成员函数
    \begin{itemize}
        \item size()
        \item capacity()
        \item empty()
        \item resize()
        \item reserve()
        \item push\_back()
        \item emplace\_back()
        \item pop\_back()
        \item insert()
        \item erase()
        \item clear()
        \item front()
        \item back()
        \item at()
        \item swap()
        \item assign()
        \item data()
        \item begin()
        \item end()
        \item rbegin()
        \item rend()
    \end{itemize}
\end{itemize}

\subsection*{第二部分：string（极详细）}
\begin{itemize}
    \item 创建
    \item 输入输出
    \item 拼接
    \item 查找
    \item 截取
    \item 替换
    \item 删除
    \item 插入
    \item 字符转换
    \item 所有常用函数
\end{itemize}

\subsection*{第三部分：algorithm}
\begin{itemize}
    \item sort
    \item reverse
    \item max/min
    \item swap
    \item find
    \item count
    \item accumulate
    \item fill
    \item replace
    \item unique
    \item lower\_bound
    \item upper\_bound
    \item binary\_search
    \item next\_permutation
    \item prev\_permutation
    \item rotate
    \item partition
\end{itemize}

\subsection*{第四部分：deque}
\subsection*{第五部分：stack}
\subsection*{第六部分：queue}
\subsection*{第七部分：priority\_queue}
\subsection*{第八部分：set / multiset}
\subsection*{第九部分：map / unordered\_map}
\subsection*{第十部分：pair}
\subsection*{第十一部分：迭代器 iterator}
\subsection*{第十二部分：常用万能模板}

\noindent\rule{\linewidth}{0.4pt}

下面开始。

\section{第一章 vector（重点）}

头文件：

\begin{lstlisting}[language=C++]
#include <vector>
using namespace std;
\end{lstlisting}

\noindent\rule{\linewidth}{0.4pt}

\section{1. vector定义}

\subsection{名称}
\begin{lstlisting}[language=C++]
vector
\end{lstlisting}

\subsection{用法}
\begin{lstlisting}[language=C++]
vector<int> v;
\end{lstlisting}
定义一个存放 int 的动态数组。

其他：
\begin{lstlisting}[language=C++]
vector<double> v1;

vector<string> v2;

vector<char> v3;
\end{lstlisting}

二维：
\begin{lstlisting}[language=C++]
vector<vector<int>> matrix;
\end{lstlisting}
相当于二维数组。

\subsection{作用}
vector 是 STL 中最常用的动态数组。

和普通数组区别：\\
普通数组：
\begin{lstlisting}[language=C++]
int a[100];
\end{lstlisting}
大小固定。

vector：
\begin{lstlisting}[language=C++]
vector<int> v;
\end{lstlisting}
可以自动扩容。

\noindent\rule{\linewidth}{0.4pt}

\section{2. vector初始化}

\subsection{名称}
初始化

\subsection{用法}

\subsubsection{空vector}
\begin{lstlisting}[language=C++]
vector<int> v;
\end{lstlisting}
创建：
\begin{lstlisting}
{}
\end{lstlisting}

\noindent\rule{\linewidth}{0.4pt}

\subsubsection{指定数量}
\begin{lstlisting}[language=C++]
vector<int> v(5);
\end{lstlisting}
结果：
\begin{lstlisting}
0 0 0 0 0
\end{lstlisting}

\noindent\rule{\linewidth}{0.4pt}

\subsubsection{指定数量和值}
\begin{lstlisting}[language=C++]
vector<int> v(5,10);
\end{lstlisting}
结果：
\begin{lstlisting}
10 10 10 10 10
\end{lstlisting}

\noindent\rule{\linewidth}{0.4pt}

\subsubsection{数组初始化}
\begin{lstlisting}[language=C++]
int a[]={1,2,3};

vector<int> v(a,a+3);
\end{lstlisting}
结果：
\begin{lstlisting}
1 2 3
\end{lstlisting}

\noindent\rule{\linewidth}{0.4pt}

\subsubsection{直接赋值}
\begin{lstlisting}[language=C++]
vector<int> v={1,2,3,4};
\end{lstlisting}
结果：
\begin{lstlisting}
1 2 3 4
\end{lstlisting}

\subsection{作用}
创建vector。

\noindent\rule{\linewidth}{0.4pt}

\section{3. push\_back()}

\subsection{名称}
\begin{lstlisting}[language=C++]
push_back()
\end{lstlisting}

\subsection{用法}
格式：
\begin{lstlisting}[language=C++]
vector名.push_back(元素);
\end{lstlisting}
例如：
\begin{lstlisting}[language=C++]
vector<int> v;

v.push_back(10);

v.push_back(20);
\end{lstlisting}
结果：
\begin{lstlisting}
10 20
\end{lstlisting}

\subsection{作用}
在vector末尾添加元素。\\
考试最常用。

\noindent\rule{\linewidth}{0.4pt}

\section{4. pop\_back()}

\subsection{名称}
\begin{lstlisting}[language=C++]
pop_back()
\end{lstlisting}

\subsection{用法}
\begin{lstlisting}[language=C++]
v.pop_back();
\end{lstlisting}
例如：
\begin{lstlisting}[language=C++]
vector<int> v={1,2,3};

v.pop_back();
\end{lstlisting}
结果：
\begin{lstlisting}
1 2
\end{lstlisting}

\subsection{作用}
删除最后一个元素。\\
注意：\\
不能指定删除哪个。\\
只能删除尾部。

\noindent\rule{\linewidth}{0.4pt}

\section{5. size()}

\subsection{名称}
\begin{lstlisting}[language=C++]
size()
\end{lstlisting}

\subsection{用法}
\begin{lstlisting}[language=C++]
v.size();
\end{lstlisting}
例如：
\begin{lstlisting}[language=C++]
vector<int> v={1,2,3};

cout<<v.size();
\end{lstlisting}
输出：
\begin{lstlisting}
3
\end{lstlisting}

\subsection{作用}
返回vector当前元素数量。\\
注意：\\
不是容量。\\
例如：
\begin{lstlisting}[language=C++]
vector<int> v(100);
\end{lstlisting}
size:
\begin{lstlisting}
100
\end{lstlisting}

\noindent\rule{\linewidth}{0.4pt}

\section{6. capacity()}

\subsection{名称}
\begin{lstlisting}[language=C++]
capacity()
\end{lstlisting}

\subsection{用法}
\begin{lstlisting}[language=C++]
v.capacity();
\end{lstlisting}
例如：
\begin{lstlisting}[language=C++]
vector<int> v;

v.push_back(1);

cout<<v.capacity();
\end{lstlisting}
可能输出：
\begin{lstlisting}
1
\end{lstlisting}
或者：
\begin{lstlisting}
2
\end{lstlisting}

\subsection{作用}
查看vector当前申请的内存空间。

区别：
\begin{table}[h]
\centering
\begin{tabular}{|l|l|}
\hline
函数 & 作用 \\ \hline
size() & 已有多少元素 \\ \hline
capacity() & 最多能放多少元素 \\ \hline
\end{tabular}
\end{table}

考试一般用size。

\noindent\rule{\linewidth}{0.4pt}

\section{7. empty()}

\subsection{名称}
\begin{lstlisting}[language=C++]
empty()
\end{lstlisting}

\subsection{用法}
\begin{lstlisting}[language=C++]
v.empty();
\end{lstlisting}
例如：
\begin{lstlisting}[language=C++]
vector<int> v;

if(v.empty())
{
    cout<<"空";
}
\end{lstlisting}
返回：
\begin{lstlisting}
true
\end{lstlisting}

\subsection{作用}
判断vector是否为空。\\
等价：
\begin{lstlisting}[language=C++]
v.size()==0
\end{lstlisting}

\noindent\rule{\linewidth}{0.4pt}

\section{8. clear()}

\subsection{名称}
\begin{lstlisting}[language=C++]
clear()
\end{lstlisting}

\subsection{用法}
\begin{lstlisting}[language=C++]
v.clear();
\end{lstlisting}
例如：
\begin{lstlisting}[language=C++]
vector<int> v={1,2,3};

v.clear();
\end{lstlisting}
结果：
\begin{lstlisting}
{}
\end{lstlisting}

\subsection{作用}
删除所有元素。\\
注意：\\
不会释放空间。\\
例如：
\begin{lstlisting}[language=C++]
v.clear();

v.capacity();
\end{lstlisting}
容量可能仍然存在。

\noindent\rule{\linewidth}{0.4pt}

\section{9. front()}

\subsection{名称}
\begin{lstlisting}[language=C++]
front()
\end{lstlisting}

\subsection{用法}
\begin{lstlisting}[language=C++]
v.front();
\end{lstlisting}
例如：
\begin{lstlisting}[language=C++]
vector<int> v={10,20,30};

cout<<v.front();
\end{lstlisting}
输出：
\begin{lstlisting}
10
\end{lstlisting}

\subsection{作用}
获取第一个元素。\\
等价：
\begin{lstlisting}[language=C++]
v[0]
\end{lstlisting}

\noindent\rule{\linewidth}{0.4pt}

\section{10. back()}

\subsection{名称}
\begin{lstlisting}[language=C++]
back()
\end{lstlisting}

\subsection{用法}
\begin{lstlisting}[language=C++]
v.back();
\end{lstlisting}
例如：
\begin{lstlisting}[language=C++]
vector<int> v={10,20,30};

cout<<v.back();
\end{lstlisting}
输出：
\begin{lstlisting}
30
\end{lstlisting}

\subsection{作用}
获取最后一个元素。\\
等价：
\begin{lstlisting}[language=C++]
v[v.size()-1]
\end{lstlisting}

\noindent\rule{\linewidth}{0.4pt}

\section{11. 下标访问 []}

\subsection{名称}
\begin{lstlisting}[language=C++]
operator[]
\end{lstlisting}

\subsection{用法}
\begin{lstlisting}[language=C++]
v[index];
\end{lstlisting}
例如：
\begin{lstlisting}[language=C++]
vector<int> v={5,6,7};

cout<<v[1];
\end{lstlisting}
输出：
\begin{lstlisting}
6
\end{lstlisting}

\subsection{作用}
访问指定位置元素。\\
注意：\\
不检查越界。

\noindent\rule{\linewidth}{0.4pt}

\section{12. at()}

\subsection{名称}
\begin{lstlisting}[language=C++]
at()
\end{lstlisting}

\subsection{用法}
\begin{lstlisting}[language=C++]
v.at(index);
\end{lstlisting}
例如：
\begin{lstlisting}[language=C++]
vector<int> v={1,2,3};

cout<<v.at(1);
\end{lstlisting}
输出：
\begin{lstlisting}
2
\end{lstlisting}

\subsection{作用}
访问元素。

区别：
\begin{lstlisting}[language=C++]
v[10];
\end{lstlisting}
越界可能崩溃。

\begin{lstlisting}[language=C++]
v.at(10);
\end{lstlisting}
会主动报错。

\noindent\rule{\linewidth}{0.4pt}

\section{13. begin()}

\subsection{名称}
\begin{lstlisting}[language=C++]
begin()
\end{lstlisting}

\subsection{用法}
\begin{lstlisting}[language=C++]
v.begin();
\end{lstlisting}
例如：
\begin{lstlisting}[language=C++]
vector<int> v={1,2,3};

cout<<*v.begin();
\end{lstlisting}
输出：
\begin{lstlisting}
1
\end{lstlisting}

\subsection{作用}
返回第一个元素的迭代器。

常用于：
\begin{lstlisting}[language=C++]
sort(v.begin(),v.end());
\end{lstlisting}

\noindent\rule{\linewidth}{0.4pt}

\section{14. end()}

\subsection{名称}
\begin{lstlisting}[language=C++]
end()
\end{lstlisting}

\subsection{用法}
\begin{lstlisting}[language=C++]
v.end();
\end{lstlisting}
例如：
\begin{lstlisting}[language=C++]
sort(v.begin(),v.end());
\end{lstlisting}

\subsection{作用}
返回最后一个元素后面的位置。\\
注意：\\
不是最后一个元素。

最后元素：
\begin{lstlisting}[language=C++]
*(v.end()-1)
\end{lstlisting}

\noindent\rule{\linewidth}{0.4pt}

\section{15. rbegin()}

\subsection{名称}
\begin{lstlisting}[language=C++]
rbegin()
\end{lstlisting}

\subsection{用法}
\begin{lstlisting}[language=C++]
v.rbegin();
\end{lstlisting}
例如：
\begin{lstlisting}[language=C++]
vector<int> v={1,2,3};

cout<<*v.rbegin();
\end{lstlisting}
输出：
\begin{lstlisting}
3
\end{lstlisting}

\subsection{作用}
反向迭代器开始位置。

\noindent\rule{\linewidth}{0.4pt}

\section{16. rend()}

\subsection{名称}
\begin{lstlisting}[language=C++]
rend()
\end{lstlisting}

\subsection{用法}
\begin{lstlisting}[language=C++]
v.rend();
\end{lstlisting}
例如：
\begin{lstlisting}[language=C++]
for(auto i=v.rbegin();i!=v.rend();i++)
{
    cout<<*i;
}
\end{lstlisting}
输出：
\begin{lstlisting}
321
\end{lstlisting}

\subsection{作用}
反向遍历结束位置。

\noindent\rule{\linewidth}{0.4pt}

\section{17. insert()}

\subsection{名称}
\begin{lstlisting}[language=C++]
insert()
\end{lstlisting}

\subsection{用法}
格式：
\begin{lstlisting}[language=C++]
v.insert(位置,元素);
\end{lstlisting}
例如：
\begin{lstlisting}[language=C++]
vector<int> v={1,3};

v.insert(v.begin()+1,2);
\end{lstlisting}
结果：
\begin{lstlisting}
1 2 3
\end{lstlisting}

\noindent\rule{\linewidth}{0.4pt}

插入多个：
\begin{lstlisting}[language=C++]
v.insert(v.begin(),3,10);
\end{lstlisting}
结果：
\begin{lstlisting}
10 10 10 ...
\end{lstlisting}

\subsection{作用}
在指定位置插入元素。

\noindent\rule{\linewidth}{0.4pt}

\section{18. erase()}

\subsection{名称}
\begin{lstlisting}[language=C++]
erase()
\end{lstlisting}

\subsection{用法}
删除一个：
\begin{lstlisting}[language=C++]
v.erase(v.begin()+2);
\end{lstlisting}
例如：
\begin{lstlisting}[language=C++]
vector<int> v={1,2,3};

v.erase(v.begin()+1);
\end{lstlisting}
结果：
\begin{lstlisting}
1 3
\end{lstlisting}

删除区间：
\begin{lstlisting}[language=C++]
v.erase(v.begin(),v.begin()+2);
\end{lstlisting}
删除：
\begin{lstlisting}
前两个元素
\end{lstlisting}

\subsection{作用}
删除指定位置或者范围元素。

\noindent\rule{\linewidth}{0.4pt}

\section{19. resize()}

\subsection{名称}
\begin{lstlisting}[language=C++]
resize()
\end{lstlisting}

\subsection{用法}
改变大小：
\begin{lstlisting}[language=C++]
v.resize(10);
\end{lstlisting}
不足补0。

例如：
\begin{lstlisting}[language=C++]
vector<int> v={1,2};

v.resize(5);
\end{lstlisting}
结果：
\begin{lstlisting}
1 2 0 0 0
\end{lstlisting}

指定值：
\begin{lstlisting}[language=C++]
v.resize(5,100);
\end{lstlisting}
结果：
\begin{lstlisting}
1 2 100 100 100
\end{lstlisting}

\subsection{作用}
改变vector元素数量。

\noindent\rule{\linewidth}{0.4pt}

\section{20. reserve()}

\subsection{名称}
\begin{lstlisting}[language=C++]
reserve()
\end{lstlisting}

\subsection{用法}
\begin{lstlisting}[language=C++]
v.reserve(100);
\end{lstlisting}

\subsection{作用}
提前申请空间。

区别：\\
resize：改变元素数量。\\
reserve：只改变容量。

例如：
\begin{lstlisting}[language=C++]
vector<int> v;

v.reserve(100);

cout<<v.size();
\end{lstlisting}
输出：
\begin{lstlisting}
0
\end{lstlisting}

\noindent\rule{\linewidth}{0.4pt}

\section{21. swap()}

\subsection{名称}
\begin{lstlisting}[language=C++]
swap()
\end{lstlisting}

\subsection{用法}
两个vector交换：
\begin{lstlisting}[language=C++]
v1.swap(v2);
\end{lstlisting}
例如：
\begin{lstlisting}[language=C++]
vector<int> a={1,2};

vector<int> b={3,4};

a.swap(b);
\end{lstlisting}
结果：

a:
\begin{lstlisting}
3 4
\end{lstlisting}
b:
\begin{lstlisting}
1 2
\end{lstlisting}

\subsection{作用}
交换两个vector内容。

\noindent\rule{\linewidth}{0.4pt}

\section{22. assign()}

\subsection{名称}
\begin{lstlisting}[language=C++]
assign()
\end{lstlisting}

\subsection{用法}
数量和值：
\begin{lstlisting}[language=C++]
v.assign(5,10);
\end{lstlisting}
结果：
\begin{lstlisting}
10 10 10 10 10
\end{lstlisting}

区间：
\begin{lstlisting}[language=C++]
v.assign(v2.begin(),v2.end());
\end{lstlisting}

\subsection{作用}
重新给vector赋值。

\noindent\rule{\linewidth}{0.4pt}

\section{23. data()}

\subsection{名称}
\begin{lstlisting}[language=C++]
data()
\end{lstlisting}

\subsection{用法}
\begin{lstlisting}[language=C++]
int *p=v.data();
\end{lstlisting}
例如：
\begin{lstlisting}[language=C++]
vector<int> v={1,2,3};

int *p=v.data();

cout<<p[0];
\end{lstlisting}
输出：
\begin{lstlisting}
1
\end{lstlisting}

\subsection{作用}
返回底层数组指针。\\
用于和C语言数组交互。

\noindent\rule{\linewidth}{0.4pt}

\section{vector遍历（考试必背）}

\subsection{方法1：下标}
\begin{lstlisting}[language=C++]
for(int i=0;i<v.size();i++)
{
    cout<<v[i];
}
\end{lstlisting}

\noindent\rule{\linewidth}{0.4pt}

\subsection{方法2：范围for}
\begin{lstlisting}[language=C++]
for(auto x:v)
{
    cout<<x;
}
\end{lstlisting}

\noindent\rule{\linewidth}{0.4pt}

\subsection{方法3：迭代器}
\begin{lstlisting}[language=C++]
for(auto it=v.begin();it!=v.end();it++)
{
    cout<<*it;
}
\end{lstlisting}

\noindent\rule{\linewidth}{0.4pt}

\subsection{方法4：反向}
\begin{lstlisting}[language=C++]
for(auto it=v.rbegin();it!=v.rend();it++)
{
    cout<<*it;
}
\end{lstlisting}

\noindent\rule{\linewidth}{0.4pt}

\section{第二章 string（重点）}

头文件：
\begin{lstlisting}[language=C++]
#include <string>
using namespace std;
\end{lstlisting}

\texttt{string} 是 STL 中处理字符串的容器。\\
相比 C 语言的传统字符数组：
\begin{lstlisting}[language=C++]
char s[100];
\end{lstlisting}
使用 \texttt{string}：
\begin{lstlisting}[language=C++]
string s;
\end{lstlisting}
优点是**不用考虑长度，不需要手动管理结束标记 \texttt{'\textbackslash 0'}**。

\noindent\rule{\linewidth}{0.4pt}

\section{1. string定义}

\subsection{名称}
\begin{lstlisting}[language=C++]
string
\end{lstlisting}

\subsection{用法}
普通定义：
\begin{lstlisting}[language=C++]
string s;
\end{lstlisting}
初始化：
\begin{lstlisting}[language=C++]
string s="hello";
\end{lstlisting}
定义多个：
\begin{lstlisting}[language=C++]
string a="abc";
string b="def";
\end{lstlisting}

\subsection{作用}
创建字符串对象。

\noindent\rule{\linewidth}{0.4pt}

\section{2. string输入}

\subsection{方法1：cin（输入一个单词）}
\subsubsection{名称}
\begin{lstlisting}[language=C++]
cin
\end{lstlisting}
\subsubsection{用法}
\begin{lstlisting}[language=C++]
string s;
cin>>s;
\end{lstlisting}
若输入：
\begin{lstlisting}
hello
\end{lstlisting}
结果：
\begin{lstlisting}[language=C++]
s="hello"
\end{lstlisting}
\textbf{注意：遇到空格、换行即停止读取。}

\subsection{方法2：getline（输入整行）}
\subsubsection{名称}
\begin{lstlisting}[language=C++]
getline()
\end{lstlisting}
\subsubsection{用法}
\begin{lstlisting}[language=C++]
getline(cin,s);
\end{lstlisting}
示例：
\begin{lstlisting}[language=C++]
string s;
getline(cin,s);
cout<<s;
\end{lstlisting}
若输入：
\begin{lstlisting}
hello world
\end{lstlisting}
输出：
\begin{lstlisting}
hello world
\end{lstlisting}

\subsection{作用}
读取包含空格的一整行字符串，**上机考试非常常用**。

\noindent\rule{\linewidth}{0.4pt}

\section{3. size()}

\subsection{名称}
\begin{lstlisting}[language=C++]
size()
\end{lstlisting}

\subsection{用法}
\begin{lstlisting}[language=C++]
s.size();
\end{lstlisting}
示例：
\begin{lstlisting}[language=C++]
string s="hello";
cout<<s.size();
\end{lstlisting}
输出结果：
\begin{lstlisting}
5
\end{lstlisting}

\subsection{作用}
获取字符串当前的实际长度。

\noindent\rule{\linewidth}{0.4pt}

\section{4. length()}

\subsection{名称}
\begin{lstlisting}[language=C++]
length()
\end{lstlisting}

\subsection{用法}
\begin{lstlisting}[language=C++]
s.length();
\end{lstlisting}
示例：
\begin{lstlisting}[language=C++]
string s="abc";
cout<<s.length();
\end{lstlisting}
输出结果：
\begin{lstlisting}
3
\end{lstlisting}

\subsection{作用}
获取字符串长度。\\
\textbf{注意：功能与 \texttt{size()} 完全一样，机考用哪个都可以。}

\noindent\rule{\linewidth}{0.4pt}

\section{5. empty()}

\subsection{名称}
\begin{lstlisting}[language=C++]
empty()
\end{lstlisting}

\subsection{用法}
\begin{lstlisting}[language=C++]
s.empty();
\end{lstlisting}
示例：
\begin{lstlisting}[language=C++]
string s="";
if(s.empty())
{
    cout<<"空";
}
\end{lstlisting}

\subsection{作用}
判断字符串是否为空（长度是否为0）。

\noindent\rule{\linewidth}{0.4pt}

\section{6. 字符访问 []}

\subsection{名称}
\begin{lstlisting}[language=C++]
operator[]
\end{lstlisting}

\subsection{用法}
格式：
\begin{lstlisting}[language=C++]
s[index];
\end{lstlisting}
例如：
\begin{lstlisting}[language=C++]
string s="hello";

cout<<s[0];
\end{lstlisting}
输出：
\begin{lstlisting}
h
\end{lstlisting}
修改：
\begin{lstlisting}[language=C++]
s[0]='H';
\end{lstlisting}
结果：
\begin{lstlisting}
Hello
\end{lstlisting}

\subsection{作用}
访问或者修改指定位置字符。\\
注意：\\
下标从 \textbf{0开始}。

\noindent\rule{\linewidth}{0.4pt}

\section{7. at()}

\subsection{名称}
\begin{lstlisting}[language=C++]
at()
\end{lstlisting}

\subsection{用法}
\begin{lstlisting}[language=C++]
s.at(index);
\end{lstlisting}
例如：
\begin{lstlisting}[language=C++]
string s="abc";

cout<<s.at(1);
\end{lstlisting}
输出：
\begin{lstlisting}
b
\end{lstlisting}
修改：
\begin{lstlisting}[language=C++]
s.at(1)='B';
\end{lstlisting}
结果：
\begin{lstlisting}
aBc
\end{lstlisting}

\subsection{作用}
访问指定位置字符。

区别：
\begin{lstlisting}[language=C++]
s[100];
\end{lstlisting}
越界不检查。

\begin{lstlisting}[language=C++]
s.at(100);
\end{lstlisting}
会抛异常。

\noindent\rule{\linewidth}{0.4pt}

\section{8. front()}

\subsection{名称}
\begin{lstlisting}[language=C++]
front()
\end{lstlisting}

\subsection{用法}
\begin{lstlisting}[language=C++]
s.front();
\end{lstlisting}
例如：
\begin{lstlisting}[language=C++]
string s="hello";

cout<<s.front();
\end{lstlisting}
输出：
\begin{lstlisting}
h
\end{lstlisting}

\subsection{作用}
获取第一个字符。\\
等价：
\begin{lstlisting}[language=C++]
s[0]
\end{lstlisting}

\noindent\rule{\linewidth}{0.4pt}

\section{9. back()}

\subsection{名称}
\begin{lstlisting}[language=C++]
back()
\end{lstlisting}

\subsection{用法}
\begin{lstlisting}[language=C++]
s.back();
\end{lstlisting}
例如：
\begin{lstlisting}[language=C++]
string s="hello";

cout<<s.back();
\end{lstlisting}
输出：
\begin{lstlisting}
o
\end{lstlisting}

\subsection{作用}
获取最后一个字符。\\
等价：
\begin{lstlisting}[language=C++]
s[s.size()-1]
\end{lstlisting}

\noindent\rule{\linewidth}{0.4pt}

\section{10. push\_back()}

\subsection{名称}
\begin{lstlisting}[language=C++]
push_back()
\end{lstlisting}

\subsection{用法}
格式：
\begin{lstlisting}[language=C++]
s.push_back(字符);
\end{lstlisting}
例如：
\begin{lstlisting}[language=C++]
string s="abc";

s.push_back('d');
\end{lstlisting}
结果：
\begin{lstlisting}
abcd
\end{lstlisting}
注意：\\
只能添加一个字符。

错误：
\begin{lstlisting}[language=C++]
s.push_back("hello");
\end{lstlisting}
正确：
\begin{lstlisting}[language=C++]
s+="hello";
\end{lstlisting}

\subsection{作用}
在字符串末尾添加字符。

\noindent\rule{\linewidth}{0.4pt}

\section{11. pop\_back()}

\subsection{名称}
\begin{lstlisting}[language=C++]
pop_back()
\end{lstlisting}

\subsection{用法}
\begin{lstlisting}[language=C++]
s.pop_back();
\end{lstlisting}
例如：
\begin{lstlisting}[language=C++]
string s="hello";

s.pop_back();
\end{lstlisting}
结果：
\begin{lstlisting}
hell
\end{lstlisting}

\subsection{作用}
删除最后一个字符。

\noindent\rule{\linewidth}{0.4pt}

\section{12. append()}

\subsection{名称}
\begin{lstlisting}[language=C++]
append()
\end{lstlisting}

\subsection{用法}
添加字符串：
\begin{lstlisting}[language=C++]
s.append("abc");
\end{lstlisting}
例如：
\begin{lstlisting}[language=C++]
string s="hello";

s.append(" world");
\end{lstlisting}
结果：
\begin{lstlisting}
hello world
\end{lstlisting}
也可以：
\begin{lstlisting}[language=C++]
s.append(t);
\end{lstlisting}
其中：
\begin{lstlisting}[language=C++]
string t="abc";
\end{lstlisting}

\subsection{作用}
在末尾追加字符串。\\
等价：
\begin{lstlisting}[language=C++]
s+="abc";
\end{lstlisting}

\noindent\rule{\linewidth}{0.4pt}

\section{13. 字符串连接 +}

\subsection{名称}
字符串拼接

\subsection{用法}
\begin{lstlisting}[language=C++]
s1+s2;
\end{lstlisting}
例如：
\begin{lstlisting}[language=C++]
string a="hello";

string b="world";

string c=a+b;
\end{lstlisting}
结果：
\begin{lstlisting}
helloworld
\end{lstlisting}
添加字符：
\begin{lstlisting}[language=C++]
s+='!';
\end{lstlisting}
结果：
\begin{lstlisting}
hello!
\end{lstlisting}

\subsection{作用}
连接两个字符串。

\noindent\rule{\linewidth}{0.4pt}

\section{14. insert()}

\subsection{名称}
\begin{lstlisting}[language=C++]
insert()
\end{lstlisting}

\subsection{用法}
格式：
\begin{lstlisting}[language=C++]
s.insert(位置,字符串);
\end{lstlisting}
例如：
\begin{lstlisting}[language=C++]
string s="abcd";

s.insert(2,"XX");
\end{lstlisting}
过程：
\begin{lstlisting}
ab + XX + cd
\end{lstlisting}
结果：
\begin{lstlisting}
abXXcd
\end{lstlisting}

\noindent\rule{\linewidth}{0.4pt}

插入一个字符：
\begin{lstlisting}[language=C++]
s.insert(0,1,'A');
\end{lstlisting}
结果：
\begin{lstlisting}
Aabcd
\end{lstlisting}

\subsection{作用}
在指定位置插入内容。

\noindent\rule{\linewidth}{0.4pt}

\section{15. erase()}

\subsection{名称}
\begin{lstlisting}[language=C++]
erase()
\end{lstlisting}

\subsection{用法1：删除指定位置}
格式：
\begin{lstlisting}[language=C++]
s.erase(pos);
\end{lstlisting}
例如：
\begin{lstlisting}[language=C++]
string s="abcdef";

s.erase(2);
\end{lstlisting}
结果：
\begin{lstlisting}
ab
\end{lstlisting}
从位置2开始全部删除。

\noindent\rule{\linewidth}{0.4pt}

\subsection{用法2：删除指定长度}
格式：
\begin{lstlisting}[language=C++]
s.erase(pos,len);
\end{lstlisting}
例如：
\begin{lstlisting}[language=C++]
string s="abcdef";

s.erase(2,3);
\end{lstlisting}
删除：
\begin{lstlisting}
cde
\end{lstlisting}
结果：
\begin{lstlisting}
abf
\end{lstlisting}

\subsection{作用}
删除字符串中的字符。

\noindent\rule{\linewidth}{0.4pt}

\section{16. clear()}

\subsection{名称}
\begin{lstlisting}[language=C++]
clear()
\end{lstlisting}

\subsection{用法}
\begin{lstlisting}[language=C++]
s.clear();
\end{lstlisting}
例如：
\begin{lstlisting}[language=C++]
string s="hello";

s.clear();
\end{lstlisting}
结果：
\begin{lstlisting}
""
\end{lstlisting}

\subsection{作用}
清空字符串。

\noindent\rule{\linewidth}{0.4pt}

\section{17. substr()}

\subsection{名称}
\begin{lstlisting}[language=C++]
substr()
\end{lstlisting}

\subsection{用法}
格式：
\begin{lstlisting}[language=C++]
s.substr(开始位置,长度);
\end{lstlisting}
例如：
\begin{lstlisting}[language=C++]
string s="abcdef";

string t=s.substr(2,3);
\end{lstlisting}
从：
\begin{lstlisting}
c开始
取3个
\end{lstlisting}
结果：
\begin{lstlisting}
cde
\end{lstlisting}

\noindent\rule{\linewidth}{0.4pt}

省略长度：
\begin{lstlisting}[language=C++]
s.substr(2);
\end{lstlisting}
结果：
\begin{lstlisting}
cdef
\end{lstlisting}

\subsection{作用}
截取字符串。\\
考试非常常用。

\noindent\rule{\linewidth}{0.4pt}

\section{18. find()}

\subsection{名称}
\begin{lstlisting}[language=C++]
find()
\end{lstlisting}

\subsection{用法}
格式：
\begin{lstlisting}[language=C++]
s.find(查找内容);
\end{lstlisting}
例如：
\begin{lstlisting}[language=C++]
string s="hello";

cout<<s.find("ll");
\end{lstlisting}
输出：
\begin{lstlisting}
2
\end{lstlisting}
因为：
\begin{lstlisting}
h e l l o
0 1 2 3 4
\end{lstlisting}

\noindent\rule{\linewidth}{0.4pt}

查找字符：
\begin{lstlisting}[language=C++]
s.find('o');
\end{lstlisting}
输出：
\begin{lstlisting}
4
\end{lstlisting}

\noindent\rule{\linewidth}{0.4pt}

不存在：
\begin{lstlisting}[language=C++]
s.find("xx");
\end{lstlisting}
返回：
\begin{lstlisting}[language=C++]
string::npos
\end{lstlisting}
例如：
\begin{lstlisting}[language=C++]
if(s.find("abc")!=string::npos)
{
    cout<<"存在";
}
\end{lstlisting}

\subsection{作用}
查找子串或者字符第一次出现的位置。

\noindent\rule{\linewidth}{0.4pt}

\section{19. rfind()}

\subsection{名称}
\begin{lstlisting}[language=C++]
rfind()
\end{lstlisting}

\subsection{用法}
\begin{lstlisting}[language=C++]
s.rfind("abc");
\end{lstlisting}
例如：
\begin{lstlisting}[language=C++]
string s="abcabc";

cout<<s.rfind("abc");
\end{lstlisting}
输出：
\begin{lstlisting}
3
\end{lstlisting}

\subsection{作用}
从后往前查找。\\
返回最后一次出现位置。

\noindent\rule{\linewidth}{0.4pt}

\section{20. replace()}

\subsection{名称}
\begin{lstlisting}[language=C++]
replace()
\end{lstlisting}

\subsection{用法}
格式：
\begin{lstlisting}[language=C++]
s.replace(位置,长度,替换内容);
\end{lstlisting}
例如：
\begin{lstlisting}[language=C++]
string s="abcdef";

s.replace(2,3,"XYZ");
\end{lstlisting}
原：
\begin{lstlisting}
ab cde f
\end{lstlisting}
替换：
\begin{lstlisting}
ab XYZ f
\end{lstlisting}
结果：
\begin{lstlisting}
abXYZf
\end{lstlisting}

\subsection{作用}
替换字符串中的部分内容。

\noindent\rule{\linewidth}{0.4pt}

\section{21. compare()}

\subsection{名称}
\begin{lstlisting}[language=C++]
compare()
\end{lstlisting}

\subsection{用法}
\begin{lstlisting}[language=C++]
s1.compare(s2);
\end{lstlisting}
例如：
\begin{lstlisting}[language=C++]
string a="abc";

string b="abc";

cout<<a.compare(b);
\end{lstlisting}
输出：
\begin{lstlisting}
0
\end{lstlisting}
结果：
\begin{table}[h]
\centering
\begin{tabular}{|l|l|}
\hline
返回值 & 意义 \\ \hline
0 & 相等 \\ \hline
<0 & a小于b \\ \hline
>0 & a大于b \\ \hline
\end{tabular}
\end{table}

\subsection{作用}
比较字符串。

\noindent\rule{\linewidth}{0.4pt}

\section{22. begin()}

\subsection{名称}
\begin{lstlisting}[language=C++]
begin()
\end{lstlisting}

\subsection{用法}
\begin{lstlisting}[language=C++]
s.begin();
\end{lstlisting}
例如：
\begin{lstlisting}[language=C++]
cout<<*s.begin();
\end{lstlisting}
输出：
\begin{lstlisting}
第一个字符
\end{lstlisting}

\subsection{作用}
返回字符串第一个字符迭代器。

常用于：
\begin{lstlisting}[language=C++]
sort(s.begin(),s.end());
\end{lstlisting}

\noindent\rule{\linewidth}{0.4pt}

\section{23. end()}

\subsection{名称}
\begin{lstlisting}[language=C++]
end()
\end{lstlisting}

\subsection{用法}
\begin{lstlisting}[language=C++]
s.end();
\end{lstlisting}
例如：
\begin{lstlisting}[language=C++]
sort(s.begin(),s.end());
\end{lstlisting}

\subsection{作用}
返回最后一个字符后的位置。

\noindent\rule{\linewidth}{0.4pt}

\section{24. rbegin()}

\subsection{名称}
\begin{lstlisting}[language=C++]
rbegin()
\end{lstlisting}

\subsection{用法}
\begin{lstlisting}[language=C++]
s.rbegin();
\end{lstlisting}
例如：
\begin{lstlisting}[language=C++]
cout<<*s.rbegin();
\end{lstlisting}
输出：
\begin{lstlisting}
最后一个字符
\end{lstlisting}

\subsection{作用}
反向遍历开始。

\noindent\rule{\linewidth}{0.4pt}

\section{25. rend()}

\subsection{名称}
\begin{lstlisting}[language=C++]
rend()
\end{lstlisting}

\subsection{用法}
\begin{lstlisting}[language=C++]
s.rend();
\end{lstlisting}
例如：
\begin{lstlisting}[language=C++]
for(auto it=s.rbegin();it!=s.rend();it++)
{
    cout<<*it;
}
\end{lstlisting}
字符串：
\begin{lstlisting}
abc
\end{lstlisting}
输出：
\begin{lstlisting}
cba
\end{lstlisting}

\subsection{作用}
反向遍历结束位置。

\noindent\rule{\linewidth}{0.4pt}

\section{26. stoi()}

\subsection{名称}
\begin{lstlisting}[language=C++]
stoi()
\end{lstlisting}

\subsection{用法}
\begin{lstlisting}[language=C++]
stoi(字符串);
\end{lstlisting}
例如：
\begin{lstlisting}[language=C++]
string s="123";

int x=stoi(s);
\end{lstlisting}
结果：
\begin{lstlisting}
x=123
\end{lstlisting}

\subsection{作用}
字符串转整数。

\noindent\rule{\linewidth}{0.4pt}

\section{27. stol()}

\subsection{名称}
\begin{lstlisting}[language=C++]
stol()
\end{lstlisting}

\subsection{用法}
\begin{lstlisting}[language=C++]
long long x=stoll(s);
\end{lstlisting}
例如：
\begin{lstlisting}[language=C++]
string s="123456789";

long long x=stoll(s);
\end{lstlisting}

\subsection{作用}
字符串转 long long。

\noindent\rule{\linewidth}{0.4pt}

\section{28. stod()}

\subsection{名称}
\begin{lstlisting}[language=C++]
stod()
\end{lstlisting}

\subsection{用法}
\begin{lstlisting}[language=C++]
double x=stod(s);
\end{lstlisting}
例如：
\begin{lstlisting}[language=C++]
string s="3.14";

double x=stod(s);
\end{lstlisting}

\subsection{作用}
字符串转 double。

\noindent\rule{\linewidth}{0.4pt}

\section{29. to\_string()}

\subsection{名称}
\begin{lstlisting}[language=C++]
to_string()
\end{lstlisting}

\subsection{用法}
\begin{lstlisting}[language=C++]
to_string(数字);
\end{lstlisting}
例如：
\begin{lstlisting}[language=C++]
int x=123;

string s=to_string(x);
\end{lstlisting}
结果：
\begin{lstlisting}
"123"
\end{lstlisting}

\subsection{作用}
数字转字符串。

\noindent\rule{\linewidth}{0.4pt}

\section{30. sort字符串}

\subsection{名称}
\begin{lstlisting}[language=C++]
sort()
\end{lstlisting}

\subsection{用法}
\begin{lstlisting}[language=C++]
sort(s.begin(),s.end());
\end{lstlisting}
例如：
\begin{lstlisting}[language=C++]
string s="dcba";

sort(s.begin(),s.end());
\end{lstlisting}
结果：
\begin{lstlisting}
abcd
\end{lstlisting}

\subsection{作用}
按照ASCII排序字符。

\noindent\rule{\linewidth}{0.4pt}

\section{31. reverse字符串}

\subsection{名称}
\begin{lstlisting}[language=C++]
reverse()
\end{lstlisting}

\subsection{用法}
\begin{lstlisting}[language=C++]
reverse(s.begin(),s.end());
\end{lstlisting}
例如：
\begin{lstlisting}[language=C++]
string s="hello";

reverse(s.begin(),s.end());
\end{lstlisting}
结果：
\begin{lstlisting}
olleh
\end{lstlisting}

\subsection{作用}
反转字符串。

\noindent\rule{\linewidth}{0.4pt}

\section{string考试高频组合}

\subsection{1. 判断回文字符串}
\begin{lstlisting}[language=C++]
string s;

string t=s;

reverse(t.begin(),t.end());

if(s==t)
    cout<<"yes";
\end{lstlisting}
作用：\\
判断正反是否一样。

\noindent\rule{\linewidth}{0.4pt}

\subsection{2. 统计字符数量}
\begin{lstlisting}[language=C++]
int cnt[26]={0};

for(char c:s)
{
    cnt[c-'a']++;
}
\end{lstlisting}
作用：\\
统计每个字母出现次数。

\noindent\rule{\linewidth}{0.4pt}

\subsection{3. 删除空格}
\begin{lstlisting}[language=C++]
s.erase(remove(s.begin(),s.end(),' '),s.end());
\end{lstlisting}
作用：\\
删除所有空格。\\
（这里涉及 algorithm 的 remove，后面详细讲）

\noindent\rule{\linewidth}{0.4pt}
\section{第三章 algorithm（STL算法库，考试核心）}

头文件：
\begin{lstlisting}[language=C++]
#include <algorithm>
\end{lstlisting}
部分数学算法需要：
\begin{lstlisting}[language=C++]
#include <numeric>
\end{lstlisting}

\noindent\rule{\linewidth}{0.4pt}

\section{1. sort()}

\subsection{名称}
\begin{lstlisting}[language=C++]
sort()
\end{lstlisting}

\noindent\rule{\linewidth}{0.4pt}

\subsection{用法}
基本格式：
\begin{lstlisting}[language=C++]
sort(开始位置,结束位置);
\end{lstlisting}
例如 vector：
\begin{lstlisting}[language=C++]
vector<int> v={3,1,5,2};

sort(v.begin(),v.end());
\end{lstlisting}
结果：
\begin{lstlisting}
1 2 3 5
\end{lstlisting}

\noindent\rule{\linewidth}{0.4pt}

数组：
\begin{lstlisting}[language=C++]
int a[]={3,1,5,2};

sort(a,a+4);
\end{lstlisting}
结果：
\begin{lstlisting}
1 2 3 5
\end{lstlisting}

\noindent\rule{\linewidth}{0.4pt}

降序：\\
格式：
\begin{lstlisting}[language=C++]
sort(begin,end,greater<类型>());
\end{lstlisting}
例如：
\begin{lstlisting}[language=C++]
sort(v.begin(),v.end(),greater<int>());
\end{lstlisting}
结果：
\begin{lstlisting}
5 3 2 1
\end{lstlisting}

\noindent\rule{\linewidth}{0.4pt}

自定义排序：\\
格式：
\begin{lstlisting}[language=C++]
sort(begin,end,比较函数);
\end{lstlisting}
例如：
\begin{lstlisting}[language=C++]
bool cmp(int a,int b)
{
    return a>b;
}

sort(v.begin(),v.end(),cmp);
\end{lstlisting}

\noindent\rule{\linewidth}{0.4pt}

\subsection{作用}
排序。\\
考试最常用函数之一。

默认：
\begin{lstlisting}
升序
\end{lstlisting}

\noindent\rule{\linewidth}{0.4pt}

\section{2. stable\_sort()}

\subsection{名称}
\begin{lstlisting}[language=C++]
stable_sort()
\end{lstlisting}

\noindent\rule{\linewidth}{0.4pt}

\subsection{用法}
和sort一样：
\begin{lstlisting}[language=C++]
stable_sort(v.begin(),v.end());
\end{lstlisting}

\noindent\rule{\linewidth}{0.4pt}

\subsection{作用}
稳定排序。

区别：\\
sort：\\
相同元素顺序可能改变。\\
stable\_sort：\\
保证相同元素原来的顺序。

考试：\\
普通排序直接sort。

\noindent\rule{\linewidth}{0.4pt}

\section{3. reverse()}

\subsection{名称}
\begin{lstlisting}[language=C++]
reverse()
\end{lstlisting}

\noindent\rule{\linewidth}{0.4pt}

\subsection{用法}
格式：
\begin{lstlisting}[language=C++]
reverse(开始,结束);
\end{lstlisting}
例如：
\begin{lstlisting}[language=C++]
vector<int> v={1,2,3,4};

reverse(v.begin(),v.end());
\end{lstlisting}
结果：
\begin{lstlisting}
4 3 2 1
\end{lstlisting}

字符串：
\begin{lstlisting}[language=C++]
string s="hello";

reverse(s.begin(),s.end());
\end{lstlisting}
结果：
\begin{lstlisting}
olleh
\end{lstlisting}

\noindent\rule{\linewidth}{0.4pt}

\subsection{作用}
反转顺序。

\noindent\rule{\linewidth}{0.4pt}

\section{4. rotate()}

\subsection{名称}
\begin{lstlisting}[language=C++]
rotate()
\end{lstlisting}

\noindent\rule{\linewidth}{0.4pt}

\subsection{用法}
格式：
\begin{lstlisting}[language=C++]
rotate(first,middle,last);
\end{lstlisting}
例如：
\begin{lstlisting}[language=C++]
vector<int> v={1,2,3,4,5};

rotate(v.begin(),v.begin()+2,v.end());
\end{lstlisting}
过程：\\
以第3个元素为开头：
\begin{lstlisting}
3 4 5 1 2
\end{lstlisting}

\noindent\rule{\linewidth}{0.4pt}

\subsection{作用}
旋转数组。

例如：\\
左旋：
\begin{lstlisting}
1 2 3 4 5

左移2位

3 4 5 1 2
\end{lstlisting}

\noindent\rule{\linewidth}{0.4pt}

\section{5. swap()}

\subsection{名称}
\begin{lstlisting}[language=C++]
swap()
\end{lstlisting}

\noindent\rule{\linewidth}{0.4pt}

\subsection{用法}
\begin{lstlisting}[language=C++]
swap(a,b);
\end{lstlisting}
例如：
\begin{lstlisting}[language=C++]
int a=10;
int b=20;

swap(a,b);
\end{lstlisting}
结果：
\begin{lstlisting}
a=20
b=10
\end{lstlisting}

\noindent\rule{\linewidth}{0.4pt}

字符串：
\begin{lstlisting}[language=C++]
swap(s1,s2);
\end{lstlisting}

\noindent\rule{\linewidth}{0.4pt}

\subsection{作用}
交换两个变量。

\noindent\rule{\linewidth}{0.4pt}

\section{6. max()}

\subsection{名称}
\begin{lstlisting}[language=C++]
max()
\end{lstlisting}

\noindent\rule{\linewidth}{0.4pt}

\subsection{用法}
两个数：
\begin{lstlisting}[language=C++]
max(a,b);
\end{lstlisting}
例如：
\begin{lstlisting}[language=C++]
cout<<max(3,5);
\end{lstlisting}
输出：
\begin{lstlisting}
5
\end{lstlisting}

\noindent\rule{\linewidth}{0.4pt}

多个：
\begin{lstlisting}[language=C++]
max({1,5,3});
\end{lstlisting}
输出：
\begin{lstlisting}
5
\end{lstlisting}

\noindent\rule{\linewidth}{0.4pt}

\subsection{作用}
返回最大值。

\noindent\rule{\linewidth}{0.4pt}

\section{7. min()}

\subsection{名称}
\begin{lstlisting}[language=C++]
min()
\end{lstlisting}

\noindent\rule{\linewidth}{0.4pt}

\subsection{用法}
\begin{lstlisting}[language=C++]
min(a,b);
\end{lstlisting}
例如：
\begin{lstlisting}[language=C++]
cout<<min(3,5);
\end{lstlisting}
输出：
\begin{lstlisting}
3
\end{lstlisting}

\noindent\rule{\linewidth}{0.4pt}

多个：
\begin{lstlisting}[language=C++]
min({3,1,5});
\end{lstlisting}
输出：
\begin{lstlisting}
1
\end{lstlisting}

\noindent\rule{\linewidth}{0.4pt}

\subsection{作用}
返回最小值。

\noindent\rule{\linewidth}{0.4pt}

\section{8. max\_element()}

\subsection{名称}
\begin{lstlisting}[language=C++]
max_element()
\end{lstlisting}

\noindent\rule{\linewidth}{0.4pt}

\subsection{用法}
格式：
\begin{lstlisting}[language=C++]
max_element(begin,end);
\end{lstlisting}
注意：\\
返回的是迭代器。

例如：
\begin{lstlisting}[language=C++]
vector<int> v={3,7,2};

auto it=max_element(v.begin(),v.end());

cout<<*it;
\end{lstlisting}
输出：
\begin{lstlisting}
7
\end{lstlisting}

\noindent\rule{\linewidth}{0.4pt}

\subsection{作用}
找容器最大元素位置。

如果想得到下标：
\begin{lstlisting}[language=C++]
int pos=max_element(v.begin(),v.end())-v.begin();
\end{lstlisting}

\noindent\rule{\linewidth}{0.4pt}

\section{9. min\_element()}

\subsection{名称}
\begin{lstlisting}[language=C++]
min_element()
\end{lstlisting}

\noindent\rule{\linewidth}{0.4pt}

\subsection{用法}
\begin{lstlisting}[language=C++]
auto it=min_element(v.begin(),v.end());
\end{lstlisting}
例如：
\begin{lstlisting}[language=C++]
cout<<*min_element(v.begin(),v.end());
\end{lstlisting}
输出：
\begin{lstlisting}
2
\end{lstlisting}

\noindent\rule{\linewidth}{0.4pt}

\subsection{作用}
找最小元素。

\noindent\rule{\linewidth}{0.4pt}

\section{10. find()}

\subsection{名称}
\begin{lstlisting}[language=C++]
find()
\end{lstlisting}

\noindent\rule{\linewidth}{0.4pt}

\subsection{用法}
格式：
\begin{lstlisting}[language=C++]
find(begin,end,value);
\end{lstlisting}
例如：
\begin{lstlisting}[language=C++]
vector<int> v={1,2,3,4};

auto it=find(v.begin(),v.end(),3);
\end{lstlisting}
找到：
\begin{lstlisting}
3
\end{lstlisting}
输出：
\begin{lstlisting}[language=C++]
cout<<*it;
\end{lstlisting}

\noindent\rule{\linewidth}{0.4pt}

不存在：
\begin{lstlisting}[language=C++]
if(it==v.end())
{
    cout<<"不存在";
}
\end{lstlisting}

\noindent\rule{\linewidth}{0.4pt}

\subsection{作用}
查找元素。

注意：\\
返回迭代器。

\noindent\rule{\linewidth}{0.4pt}

\section{11. count()}

\subsection{名称}
\begin{lstlisting}[language=C++]
count()
\end{lstlisting}

\noindent\rule{\linewidth}{0.4pt}

\subsection{用法}
格式：
\begin{lstlisting}[language=C++]
count(begin,end,value);
\end{lstlisting}
例如：
\begin{lstlisting}[language=C++]
vector<int> v={1,2,2,3};

int x=count(v.begin(),v.end(),2);
\end{lstlisting}
结果：
\begin{lstlisting}
2
\end{lstlisting}

\noindent\rule{\linewidth}{0.4pt}

字符串：
\begin{lstlisting}[language=C++]
string s="hello";

cout<<count(s.begin(),s.end(),'l');
\end{lstlisting}
输出：
\begin{lstlisting}
2
\end{lstlisting}

\noindent\rule{\linewidth}{0.4pt}

\subsection{作用}
统计元素出现次数。

\noindent\rule{\linewidth}{0.4pt}

\section{12. accumulate()}

头文件：
\begin{lstlisting}[language=C++]
#include <numeric>
\end{lstlisting}

\subsection{名称}
\begin{lstlisting}[language=C++]
accumulate()
\end{lstlisting}

\noindent\rule{\linewidth}{0.4pt}

\subsection{用法}
格式：
\begin{lstlisting}[language=C++]
accumulate(begin,end,初始值);
\end{lstlisting}
三个参数：
\begin{enumerate}
    \item 开始位置
    \item 结束位置
    \item 初始值
\end{enumerate}

例如：
\begin{lstlisting}[language=C++]
vector<int> v={1,2,3,4};

int sum=accumulate(v.begin(),v.end(),0);
\end{lstlisting}
计算：
\begin{lstlisting}
0+1+2+3+4
\end{lstlisting}
结果：
\begin{lstlisting}
10
\end{lstlisting}

\noindent\rule{\linewidth}{0.4pt}

注意：\\
第三个参数类型决定返回类型。

例如：
\begin{lstlisting}[language=C++]
vector<double> v={1.2,2.3};

double sum=accumulate(v.begin(),v.end(),0.0);
\end{lstlisting}
必须：
\begin{lstlisting}
0.0
\end{lstlisting}
不能：
\begin{lstlisting}
0
\end{lstlisting}
否则可能整数计算。

\noindent\rule{\linewidth}{0.4pt}

\subsection{作用}
求和。

\noindent\rule{\linewidth}{0.4pt}

\section{13. fill()}

\subsection{名称}
\begin{lstlisting}[language=C++]
fill()
\end{lstlisting}

\noindent\rule{\linewidth}{0.4pt}

\subsection{用法}
格式：
\begin{lstlisting}[language=C++]
fill(begin,end,value);
\end{lstlisting}
例如：\\
数组：
\begin{lstlisting}[language=C++]
int a[5];

fill(a,a+5,10);
\end{lstlisting}
结果：
\begin{lstlisting}
10 10 10 10 10
\end{lstlisting}
vector：
\begin{lstlisting}[language=C++]
fill(v.begin(),v.end(),0);
\end{lstlisting}

\noindent\rule{\linewidth}{0.4pt}

\subsection{作用}
全部赋值。

\noindent\rule{\linewidth}{0.4pt}

\section{14. replace()}

\subsection{名称}
\begin{lstlisting}[language=C++]
replace()
\end{lstlisting}

\noindent\rule{\linewidth}{0.4pt}

\subsection{用法}
格式：
\begin{lstlisting}[language=C++]
replace(begin,end,旧值,新值);
\end{lstlisting}
例如：
\begin{lstlisting}[language=C++]
vector<int> v={1,2,2,3};

replace(v.begin(),v.end(),2,99);
\end{lstlisting}
结果：
\begin{lstlisting}
1 99 99 3
\end{lstlisting}

\noindent\rule{\linewidth}{0.4pt}

字符串：
\begin{lstlisting}[language=C++]
string s="hello";

replace(s.begin(),s.end(),'l','x');
\end{lstlisting}
结果：
\begin{lstlisting}
hexxo
\end{lstlisting}

\noindent\rule{\linewidth}{0.4pt}

\subsection{作用}
替换所有指定元素。

\noindent\rule{\linewidth}{0.4pt}

\section{15. unique()}

\subsection{名称}
\begin{lstlisting}[language=C++]
unique()
\end{lstlisting}

\noindent\rule{\linewidth}{0.4pt}

\subsection{用法}
常和sort一起：
\begin{lstlisting}[language=C++]
sort(v.begin(),v.end());

v.erase(unique(v.begin(),v.end()),v.end());
\end{lstlisting}
例如：\\
原：
\begin{lstlisting}
1 2 2 3 3 3
\end{lstlisting}
unique：
\begin{lstlisting}
1 2 3 ? ? ?
\end{lstlisting}
erase删除后：
\begin{lstlisting}
1 2 3
\end{lstlisting}

\noindent\rule{\linewidth}{0.4pt}

\subsection{作用}
去除连续重复元素。\\
注意：\\
单独unique不会真正删除。

\noindent\rule{\linewidth}{0.4pt}

\section{16. lower\_bound()}

\subsection{名称}
\begin{lstlisting}[language=C++]
lower_bound()
\end{lstlisting}

\noindent\rule{\linewidth}{0.4pt}

\subsection{用法}
要求：\\
数组必须有序。

格式：
\begin{lstlisting}[language=C++]
lower_bound(begin,end,value);
\end{lstlisting}
例如：
\begin{lstlisting}[language=C++]
vector<int> v={1,3,5,7};

auto it=lower_bound(v.begin(),v.end(),5);
\end{lstlisting}
结果：
\begin{lstlisting}
5的位置
\end{lstlisting}
下标：
\begin{lstlisting}[language=C++]
int pos=it-v.begin();
\end{lstlisting}
输出：
\begin{lstlisting}
2
\end{lstlisting}

\noindent\rule{\linewidth}{0.4pt}

\subsection{作用}
寻找第一个：
\begin{lstlisting}
>= value
\end{lstlisting}
的位置。

例如：\\
数组：
\begin{lstlisting}
1 3 5 7
\end{lstlisting}
找4：\\
返回：
\begin{lstlisting}
5
\end{lstlisting}

\noindent\rule{\linewidth}{0.4pt}

\section{17. upper\_bound()}

\subsection{名称}
\begin{lstlisting}[language=C++]
upper_bound()
\end{lstlisting}

\noindent\rule{\linewidth}{0.4pt}

\subsection{用法}
\begin{lstlisting}[language=C++]
upper_bound(v.begin(),v.end(),value);
\end{lstlisting}
例如：
\begin{lstlisting}[language=C++]
vector<int> v={1,3,5,7};

auto it=upper_bound(v.begin(),v.end(),5);
\end{lstlisting}
返回：
\begin{lstlisting}
7
\end{lstlisting}

\noindent\rule{\linewidth}{0.4pt}

\subsection{作用}
寻找第一个：
\begin{lstlisting}
> value
\end{lstlisting}
的位置。

\noindent\rule{\linewidth}{0.4pt}

\section{18. binary\_search()}

\subsection{名称}
\begin{lstlisting}[language=C++]
binary_search()
\end{lstlisting}

\noindent\rule{\linewidth}{0.4pt}

\subsection{用法}
要求：\\
必须排序。

格式：
\begin{lstlisting}[language=C++]
binary_search(begin,end,value);
\end{lstlisting}
例如：
\begin{lstlisting}[language=C++]
vector<int> v={1,3,5,7};

bool flag=binary_search(v.begin(),v.end(),5);
\end{lstlisting}
结果：
\begin{lstlisting}
true
\end{lstlisting}

\noindent\rule{\linewidth}{0.4pt}

\subsection{作用}
二分查找是否存在。

返回：
\begin{lstlisting}
true/false
\end{lstlisting}

\noindent\rule{\linewidth}{0.4pt}

\section{19. next\_permutation()}

\subsection{名称}
\begin{lstlisting}[language=C++]
next_permutation()
\end{lstlisting}

\noindent\rule{\linewidth}{0.4pt}

\subsection{用法}
\begin{lstlisting}[language=C++]
next_permutation(begin,end);
\end{lstlisting}
例如：
\begin{lstlisting}[language=C++]
string s="abc";

next_permutation(s.begin(),s.end());
\end{lstlisting}
结果：
\begin{lstlisting}
acb
\end{lstlisting}
连续：
\begin{lstlisting}[language=C++]
while(next_permutation(s.begin(),s.end()))
{
    cout<<s;
}
\end{lstlisting}
输出所有排列。

\noindent\rule{\linewidth}{0.4pt}

\subsection{作用}
生成下一个字典序排列。

\noindent\rule{\linewidth}{0.4pt}

\section{20. prev\_permutation()}

\subsection{名称}
\begin{lstlisting}[language=C++]
prev_permutation()
\end{lstlisting}

\noindent\rule{\linewidth}{0.4pt}

\subsection{用法}
\begin{lstlisting}[language=C++]
prev_permutation(s.begin(),s.end());
\end{lstlisting}
例如：
\begin{lstlisting}
cba
\end{lstlisting}
变：
\begin{lstlisting}
cab
\end{lstlisting}

\noindent\rule{\linewidth}{0.4pt}

\subsection{作用}
生成上一个排列。

\noindent\rule{\linewidth}{0.4pt}

\section{algorithm 高频考试组合}

\subsection{1. vector排序去重}
\begin{lstlisting}[language=C++]
sort(v.begin(),v.end());

v.erase(unique(v.begin(),v.end()),v.end());
\end{lstlisting}
作用：\\
去除重复元素。

\noindent\rule{\linewidth}{0.4pt}

\subsection{2. 求最大最小值}
\begin{lstlisting}[language=C++]
int mx=*max_element(v.begin(),v.end());

int mn=*min_element(v.begin(),v.end());
\end{lstlisting}

\noindent\rule{\linewidth}{0.4pt}

\subsection{3. 求和}
\begin{lstlisting}[language=C++]
int sum=accumulate(v.begin(),v.end(),0);
\end{lstlisting}

\noindent\rule{\linewidth}{0.4pt}

\subsection{4. 判断存在}
\begin{lstlisting}[language=C++]
if(find(v.begin(),v.end(),x)!=v.end())
{
    cout<<"存在";
}
\end{lstlisting}

\noindent\rule{\linewidth}{0.4pt}

\section{第四章 deque / stack / queue / priority\_queue}

\noindent\rule{\linewidth}{0.4pt}

\section{一、deque（双端队列）}

头文件：
\begin{lstlisting}[language=C++]
#include <deque>
\end{lstlisting}

\noindent\rule{\linewidth}{0.4pt}

\subsection{1. deque定义}

\subsection{名称}
\begin{lstlisting}[language=C++]
deque
\end{lstlisting}

\subsection{用法}
\begin{lstlisting}[language=C++]
deque<int> dq;
\end{lstlisting}
初始化：
\begin{lstlisting}[language=C++]
deque<int> dq={1,2,3};
\end{lstlisting}

\subsection{作用}
双端队列。

特点：
\begin{itemize}
    \item 两端都可以插入
    \item 两端都可以删除
    \item 支持随机访问
\end{itemize}

类似：
\begin{lstlisting}
vector + queue
\end{lstlisting}

\noindent\rule{\linewidth}{0.4pt}

\section{2. push\_back()}

\subsection{名称}
\begin{lstlisting}[language=C++]
push_back()
\end{lstlisting}

\subsection{用法}
\begin{lstlisting}[language=C++]
dq.push_back(x);
\end{lstlisting}
例如：
\begin{lstlisting}[language=C++]
deque<int> dq;

dq.push_back(10);
dq.push_back(20);
\end{lstlisting}
结果：
\begin{lstlisting}
10 20
\end{lstlisting}

\subsection{作用}
尾部添加元素。

\noindent\rule{\linewidth}{0.4pt}

\section{3. push\_front()}

\subsection{名称}
\begin{lstlisting}[language=C++]
push_front()
\end{lstlisting}

\subsection{用法}
\begin{lstlisting}[language=C++]
dq.push_front(x);
\end{lstlisting}
例如：
\begin{lstlisting}[language=C++]
deque<int> dq={2,3};

dq.push_front(1);
\end{lstlisting}
结果：
\begin{lstlisting}
1 2 3
\end{lstlisting}

\subsection{作用}
头部插入元素。

\noindent\rule{\linewidth}{0.4pt}

\section{4. pop\_back()}

\subsection{名称}
\begin{lstlisting}[language=C++]
pop_back()
\end{lstlisting}

\subsection{用法}
\begin{lstlisting}[language=C++]
dq.pop_back();
\end{lstlisting}
例如：
\begin{lstlisting}[language=C++]
deque<int> dq={1,2,3};

dq.pop_back();
\end{lstlisting}
结果：
\begin{lstlisting}
1 2
\end{lstlisting}

\subsection{作用}
删除最后元素。

\noindent\rule{\linewidth}{0.4pt}

\section{5. pop\_front()}

\subsection{名称}
\begin{lstlisting}[language=C++]
pop_front()
\end{lstlisting}

\subsection{用法}
\begin{lstlisting}[language=C++]
dq.pop_front();
\end{lstlisting}
例如：
\begin{lstlisting}[language=C++]
deque<int> dq={1,2,3};

dq.pop_front();
\end{lstlisting}
结果：
\begin{lstlisting}
2 3
\end{lstlisting}

\subsection{作用}
删除第一个元素。

\noindent\rule{\linewidth}{0.4pt}

\section{6. front()}

\subsection{名称}
\begin{lstlisting}[language=C++]
front()
\end{lstlisting}

\subsection{用法}
\begin{lstlisting}[language=C++]
dq.front();
\end{lstlisting}
例如：
\begin{lstlisting}[language=C++]
cout<<dq.front();
\end{lstlisting}

\subsection{作用}
访问第一个元素。

\noindent\rule{\linewidth}{0.4pt}

\section{7. back()}

\subsection{名称}
\begin{lstlisting}[language=C++]
back()
\end{lstlisting}

\subsection{用法}
\begin{lstlisting}[language=C++]
dq.back();
\end{lstlisting}

\subsection{作用}
访问最后元素。

\noindent\rule{\linewidth}{0.4pt}

\section{8. size()}

\subsection{名称}
\begin{lstlisting}[language=C++]
size()
\end{lstlisting}

\subsection{用法}
\begin{lstlisting}[language=C++]
dq.size();
\end{lstlisting}

\subsection{作用}
获取元素数量。

\noindent\rule{\linewidth}{0.4pt}

\section{9. empty()}

\subsection{名称}
\begin{lstlisting}[language=C++]
empty()
\end{lstlisting}

\subsection{用法}
\begin{lstlisting}[language=C++]
dq.empty();
\end{lstlisting}

\subsection{作用}
判断是否为空。

\noindent\rule{\linewidth}{0.4pt}

\section{10. clear()}

\subsection{名称}
\begin{lstlisting}[language=C++]
clear()
\end{lstlisting}

\subsection{用法}
\begin{lstlisting}[language=C++]
dq.clear();
\end{lstlisting}

\subsection{作用}
清空deque。

\noindent\rule{\linewidth}{0.4pt}

\section{二、stack（栈）}

头文件：
\begin{lstlisting}[language=C++]
#include <stack>
\end{lstlisting}

\noindent\rule{\linewidth}{0.4pt}

\subsection{1. stack定义}

\subsection{名称}
\begin{lstlisting}[language=C++]
stack
\end{lstlisting}

\subsection{用法}
\begin{lstlisting}[language=C++]
stack<int> st;
\end{lstlisting}

\subsection{作用}
先进后出：
\begin{lstlisting}
LIFO
\end{lstlisting}
类似：\\
一摞盘子。\\
最后放的最先拿。

\noindent\rule{\linewidth}{0.4pt}

\section{2. push()}

\subsection{名称}
\begin{lstlisting}[language=C++]
push()
\end{lstlisting}

\subsection{用法}
\begin{lstlisting}[language=C++]
st.push(x);
\end{lstlisting}
例如：
\begin{lstlisting}[language=C++]
stack<int> st;

st.push(1);
st.push(2);
st.push(3);
\end{lstlisting}
栈：
\begin{lstlisting}
3
2
1
\end{lstlisting}

\subsection{作用}
入栈。

\noindent\rule{\linewidth}{0.4pt}

\section{3. pop()}

\subsection{名称}
\begin{lstlisting}[language=C++]
pop()
\end{lstlisting}

\subsection{用法}
\begin{lstlisting}[language=C++]
st.pop();
\end{lstlisting}
例如：
\begin{lstlisting}[language=C++]
st.pop();
\end{lstlisting}
删除：
\begin{lstlisting}
3
\end{lstlisting}
剩：
\begin{lstlisting}
2
1
\end{lstlisting}

\subsection{作用}
删除栈顶元素。

注意：\\
pop没有返回值。

错误：
\begin{lstlisting}[language=C++]
int x=st.pop();
\end{lstlisting}
错误。

\noindent\rule{\linewidth}{0.4pt}

\section{4. top()}

\subsection{名称}
\begin{lstlisting}[language=C++]
top()
\end{lstlisting}

\subsection{用法}
\begin{lstlisting}[language=C++]
st.top();
\end{lstlisting}
例如：
\begin{lstlisting}[language=C++]
cout<<st.top();
\end{lstlisting}

\subsection{作用}
查看栈顶元素。

\noindent\rule{\linewidth}{0.4pt}

\section{5. size()}

\subsection{名称}
\begin{lstlisting}[language=C++]
size()
\end{lstlisting}

\subsection{用法}
\begin{lstlisting}[language=C++]
st.size();
\end{lstlisting}

\subsection{作用}
获取元素数量。

\noindent\rule{\linewidth}{0.4pt}

\section{6. empty()}

\subsection{名称}
\begin{lstlisting}[language=C++]
empty()
\end{lstlisting}

\subsection{用法}
\begin{lstlisting}[language=C++]
st.empty();
\end{lstlisting}

\subsection{作用}
判断栈是否为空。

\noindent\rule{\linewidth}{0.4pt}

\subsection{stack经典模板}

\subsubsection{括号匹配}
\begin{lstlisting}[language=C++]
string s;

stack<char> st;

for(char c:s)
{
    if(c=='(')
        st.push(c);

    else if(c==')')
    {
        if(st.empty())
            cout<<"NO";

        st.pop();
    }
}
\end{lstlisting}
作用：\\
判断括号是否合法。

\noindent\rule{\linewidth}{0.4pt}

\section{三、queue（普通队列）}

头文件：
\begin{lstlisting}[language=C++]
#include <queue>
\end{lstlisting}

\noindent\rule{\linewidth}{0.4pt}

\section{1. queue定义}

\subsection{名称}
\begin{lstlisting}[language=C++]
queue
\end{lstlisting}

\subsection{用法}
\begin{lstlisting}[language=C++]
queue<int> q;
\end{lstlisting}

\subsection{作用}
先进先出：
\begin{lstlisting}
FIFO
\end{lstlisting}
类似排队：\\
先来先走。

\noindent\rule{\linewidth}{0.4pt}

\section{2. push()}

\subsection{名称}
\begin{lstlisting}[language=C++]
push()
\end{lstlisting}

\subsection{用法}
\begin{lstlisting}[language=C++]
q.push(x);
\end{lstlisting}
例如：
\begin{lstlisting}[language=C++]
q.push(1);
q.push(2);
q.push(3);
\end{lstlisting}
队列：
\begin{lstlisting}
1 2 3
\end{lstlisting}

\subsection{作用}
队尾加入元素。

\noindent\rule{\linewidth}{0.4pt}

\section{3. pop()}

\subsection{名称}
\begin{lstlisting}[language=C++]
pop()
\end{lstlisting}

\subsection{用法}
\begin{lstlisting}[language=C++]
q.pop();
\end{lstlisting}
例如：\\
原：
\begin{lstlisting}
1 2 3
\end{lstlisting}
执行：
\begin{lstlisting}[language=C++]
q.pop();
\end{lstlisting}
结果：
\begin{lstlisting}
2 3
\end{lstlisting}

\subsection{作用}
删除队头元素。

注意：\\
没有返回值。

\noindent\rule{\linewidth}{0.4pt}

\section{4. front()}

\subsection{名称}
\begin{lstlisting}[language=C++]
front()
\end{lstlisting}

\subsection{用法}
\begin{lstlisting}[language=C++]
q.front();
\end{lstlisting}
例如：
\begin{lstlisting}[language=C++]
cout<<q.front();
\end{lstlisting}
输出：\\
队头元素。

\subsection{作用}
访问队头。

\noindent\rule{\linewidth}{0.4pt}

\section{5. back()}

\subsection{名称}
\begin{lstlisting}[language=C++]
back()
\end{lstlisting}

\subsection{用法}
\begin{lstlisting}[language=C++]
q.back();
\end{lstlisting}

\subsection{作用}
访问队尾。

\noindent\rule{\linewidth}{0.4pt}

\section{6. size()}

\subsection{名称}
\begin{lstlisting}[language=C++]
size()
\end{lstlisting}

\subsection{用法}
\begin{lstlisting}[language=C++]
q.size();
\end{lstlisting}

\subsection{作用}
元素数量。

\noindent\rule{\linewidth}{0.4pt}

\section{7. empty()}

\subsection{名称}
\begin{lstlisting}[language=C++]
empty()
\end{lstlisting}

\subsection{用法}
\begin{lstlisting}[language=C++]
q.empty();
\end{lstlisting}

\subsection{作用}
判断队列为空。

\noindent\rule{\linewidth}{0.4pt}

\section{四、priority\_queue（优先队列）}

头文件：
\begin{lstlisting}[language=C++]
#include <queue>
\end{lstlisting}

\noindent\rule{\linewidth}{0.4pt}

\section{1. priority\_queue定义}

\subsection{名称}
\begin{lstlisting}[language=C++]
priority_queue
\end{lstlisting}

\subsection{用法}
默认：
\begin{lstlisting}[language=C++]
priority_queue<int> pq;
\end{lstlisting}

\subsection{作用}
自动排序的队列。

默认：\\
大根堆：
\begin{lstlisting}
最大值永远在顶部
\end{lstlisting}

\noindent\rule{\linewidth}{0.4pt}

\section{2. push()}

\subsection{名称}
\begin{lstlisting}[language=C++]
push()
\end{lstlisting}

\subsection{用法}
\begin{lstlisting}[language=C++]
pq.push(x);
\end{lstlisting}
例如：
\begin{lstlisting}[language=C++]
priority_queue<int> pq;

pq.push(3);
pq.push(10);
pq.push(5);
\end{lstlisting}
内部：
\begin{lstlisting}
10
5
3
\end{lstlisting}

\noindent\rule{\linewidth}{0.4pt}

\section{3. top()}

\subsection{名称}
\begin{lstlisting}[language=C++]
top()
\end{lstlisting}

\subsection{用法}
\begin{lstlisting}[language=C++]
pq.top();
\end{lstlisting}
例如：
\begin{lstlisting}[language=C++]
cout<<pq.top();
\end{lstlisting}
输出：
\begin{lstlisting}
10
\end{lstlisting}

\subsection{作用}
查看最大元素。

\noindent\rule{\linewidth}{0.4pt}

\section{4. pop()}

\subsection{名称}
\begin{lstlisting}[language=C++]
pop()
\end{lstlisting}

\subsection{用法}
\begin{lstlisting}[language=C++]
pq.pop();
\end{lstlisting}
删除：\\
当前最大值。

\noindent\rule{\linewidth}{0.4pt}

\section{5. empty()}

\subsection{名称}
\begin{lstlisting}[language=C++]
empty()
\end{lstlisting}

\subsection{用法}
\begin{lstlisting}[language=C++]
pq.empty();
\end{lstlisting}

\subsection{作用}
判断为空。

\noindent\rule{\linewidth}{0.4pt}

\section{6. size()}

\subsection{名称}
\begin{lstlisting}[language=C++]
size()
\end{lstlisting}

\subsection{用法}
\begin{lstlisting}[language=C++]
pq.size();
\end{lstlisting}

\subsection{作用}
元素数量。

\noindent\rule{\linewidth}{0.4pt}

\section{7. 小根堆}

\subsection{名称}
小根priority\_queue

\subsection{用法}
\begin{lstlisting}[language=C++]
priority_queue<int,vector<int>,greater<int>> pq;
\end{lstlisting}
例如：
\begin{lstlisting}[language=C++]
pq.push(5);
pq.push(2);
pq.push(8);
\end{lstlisting}
top：
\begin{lstlisting}
2
\end{lstlisting}

\noindent\rule{\linewidth}{0.4pt}

\subsection{作用}
让最小值在顶部。\\
考试非常重要。

\noindent\rule{\linewidth}{0.4pt}

\section{priority\_queue自定义排序}

例如：\\
结构体：
\begin{lstlisting}[language=C++]
struct Node
{
    int x;
    int y;
};
\end{lstlisting}
比较：
\begin{lstlisting}[language=C++]
struct cmp
{
    bool operator()(Node a,Node b)
    {
        return a.x>b.x;
    }
};
\end{lstlisting}
定义：
\begin{lstlisting}[language=C++]
priority_queue<Node,vector<Node>,cmp> pq;
\end{lstlisting}
作用：\\
复杂排序。

\noindent\rule{\linewidth}{0.4pt}

\section{五、四种容器总结（考试记忆）}

\begin{table}[h]
\centering
\begin{tabular}{|l|l|l|}
\hline
容器 & 特点 & 常用函数 \\ \hline
vector & 动态数组 & push\_back erase sort \\ \hline
deque & 两端操作 & push\_front \\ \hline
stack & 后进先出 & push pop top \\ \hline
queue & 先进先出 & push pop front \\ \hline
priority\_queue & 自动排序 & push top pop \\ \hline
\end{tabular}
\end{table}

\noindent\rule{\linewidth}{0.4pt}

\section{第五章 set / multiset / map / unordered\_map（STL核心）}

这一章非常重要。

考试中常见：
\begin{itemize}
    \item 去重
    \item 查找某个元素是否存在
    \item 统计次数
    \item 建立映射关系
    \item 字符频率统计
\end{itemize}

\noindent\rule{\linewidth}{0.4pt}

\section{一、set（集合）}

头文件：
\begin{lstlisting}[language=C++]
#include <set>
\end{lstlisting}

\noindent\rule{\linewidth}{0.4pt}

\section{1. set定义}

\subsection{名称}
\begin{lstlisting}[language=C++]
set
\end{lstlisting}

\subsection{用法}
\begin{lstlisting}[language=C++]
set<int> s;
\end{lstlisting}
初始化：
\begin{lstlisting}[language=C++]
set<int> s={3,1,2,2};
\end{lstlisting}
结果：
\begin{lstlisting}
1 2 3
\end{lstlisting}

特点：
\begin{enumerate}
    \item 自动排序
    \item 自动去重
\end{enumerate}

\noindent\rule{\linewidth}{0.4pt}

\subsection{作用}
保存唯一元素。\\
类似数学集合。

\noindent\rule{\linewidth}{0.4pt}

\section{2. insert()}

\subsection{名称}
\begin{lstlisting}[language=C++]
insert()
\end{lstlisting}

\subsection{用法}
\begin{lstlisting}[language=C++]
s.insert(value);
\end{lstlisting}
例如：
\begin{lstlisting}[language=C++]
set<int> s;

s.insert(3);
s.insert(1);
s.insert(3);
\end{lstlisting}
结果：
\begin{lstlisting}
1 3
\end{lstlisting}

\subsection{作用}
插入元素。\\
注意：\\
重复元素不会加入。

\noindent\rule{\linewidth}{0.4pt}

\section{3. erase()}

\subsection{名称}
\begin{lstlisting}[language=C++]
erase()
\end{lstlisting}

\subsection{用法1：删除指定元素}
\begin{lstlisting}[language=C++]
s.erase(value);
\end{lstlisting}
例如：
\begin{lstlisting}[language=C++]
set<int> s={1,2,3};

s.erase(2);
\end{lstlisting}
结果：
\begin{lstlisting}
1 3
\end{lstlisting}

\noindent\rule{\linewidth}{0.4pt}

\subsection{用法2：删除迭代器位置}
\begin{lstlisting}[language=C++]
auto it=s.begin();

s.erase(it);
\end{lstlisting}
删除第一个元素。

\noindent\rule{\linewidth}{0.4pt}

\subsection{用法3：删除区间}
\begin{lstlisting}[language=C++]
s.erase(s.begin(),s.end());
\end{lstlisting}
清空。

\noindent\rule{\linewidth}{0.4pt}

\subsection{作用}
删除元素。

\noindent\rule{\linewidth}{0.4pt}

\section{4. find()}

\subsection{名称}
\begin{lstlisting}[language=C++]
find()
\end{lstlisting}

\subsection{用法}
\begin{lstlisting}[language=C++]
s.find(value);
\end{lstlisting}
例如：
\begin{lstlisting}[language=C++]
set<int> s={1,2,3};

auto it=s.find(2);
\end{lstlisting}
如果存在：\\
返回位置。

不存在：
\begin{lstlisting}[language=C++]
s.end()
\end{lstlisting}
判断：
\begin{lstlisting}[language=C++]
if(s.find(2)!=s.end())
{
    cout<<"存在";
}
\end{lstlisting}

\subsection{作用}
查找元素。

\noindent\rule{\linewidth}{0.4pt}

\section{5. count()}

\subsection{名称}
\begin{lstlisting}[language=C++]
count()
\end{lstlisting}

\subsection{用法}
\begin{lstlisting}[language=C++]
s.count(value);
\end{lstlisting}
例如：
\begin{lstlisting}[language=C++]
set<int> s={1,2,3};

cout<<s.count(2);
\end{lstlisting}
输出：
\begin{lstlisting}
1
\end{lstlisting}
不存在：
\begin{lstlisting}
0
\end{lstlisting}

\subsection{作用}
判断元素是否存在。\\
因为set没有重复：\\
返回只能：
\begin{lstlisting}
0 或 1
\end{lstlisting}

\noindent\rule{\linewidth}{0.4pt}

\section{6. size()}

\subsection{名称}
\begin{lstlisting}[language=C++]
size()
\end{lstlisting}

\subsection{用法}
\begin{lstlisting}[language=C++]
s.size();
\end{lstlisting}

\subsection{作用}
元素数量。

\noindent\rule{\linewidth}{0.4pt}

\section{7. empty()}

\subsection{名称}
\begin{lstlisting}[language=C++]
empty()
\end{lstlisting}

\subsection{用法}
\begin{lstlisting}[language=C++]
s.empty();
\end{lstlisting}

\subsection{作用}
判断为空。

\noindent\rule{\linewidth}{0.4pt}

\section{8. begin()}

\subsection{名称}
\begin{lstlisting}[language=C++]
begin()
\end{lstlisting}

\subsection{用法}
\begin{lstlisting}[language=C++]
auto it=s.begin();
\end{lstlisting}
例如：
\begin{lstlisting}[language=C++]
cout<<*s.begin();
\end{lstlisting}
输出：\\
最小元素。

\noindent\rule{\linewidth}{0.4pt}

\section{9. end()}

\subsection{名称}
\begin{lstlisting}[language=C++]
end()
\end{lstlisting}

\subsection{用法}
\begin{lstlisting}[language=C++]
s.end();
\end{lstlisting}
作用：\\
最后一个元素之后。

\noindent\rule{\linewidth}{0.4pt}

\section{10. lower\_bound()}

\subsection{名称}
\begin{lstlisting}[language=C++]
lower_bound()
\end{lstlisting}

\subsection{用法}
\begin{lstlisting}[language=C++]
auto it=s.lower_bound(x);
\end{lstlisting}
例如：
\begin{lstlisting}[language=C++]
set<int> s={1,3,5,7};

auto it=s.lower_bound(4);
\end{lstlisting}
返回：
\begin{lstlisting}
5
\end{lstlisting}

\subsection{作用}
找第一个：
\begin{lstlisting}
>=x
\end{lstlisting}

\noindent\rule{\linewidth}{0.4pt}

\section{11. upper\_bound()}

\subsection{名称}
\begin{lstlisting}[language=C++]
upper_bound()
\end{lstlisting}

\subsection{用法}
\begin{lstlisting}[language=C++]
auto it=s.upper_bound(x);
\end{lstlisting}
例如：
\begin{lstlisting}[language=C++]
auto it=s.upper_bound(5);
\end{lstlisting}
返回：
\begin{lstlisting}
7
\end{lstlisting}

\subsection{作用}
找第一个：
\begin{lstlisting}
>x
\end{lstlisting}

\noindent\rule{\linewidth}{0.4pt}

\section{set遍历}

\subsection{方法1：范围for}
\begin{lstlisting}[language=C++]
for(auto x:s)
{
    cout<<x;
}
\end{lstlisting}
输出：
\begin{lstlisting}
1 2 3
\end{lstlisting}

\noindent\rule{\linewidth}{0.4pt}

\subsection{方法2：迭代器}
\begin{lstlisting}[language=C++]
for(auto it=s.begin();it!=s.end();it++)
{
    cout<<*it;
}
\end{lstlisting}

\noindent\rule{\linewidth}{0.4pt}

\section{set常用场景}

\subsection{去重排序}
原：
\begin{lstlisting}[language=C++]
vector<int> v={3,1,2,2};
\end{lstlisting}
方法：
\begin{lstlisting}[language=C++]
set<int> s(v.begin(),v.end());
\end{lstlisting}
结果：
\begin{lstlisting}
1 2 3
\end{lstlisting}

\noindent\rule{\linewidth}{0.4pt}
\noindent\rule{\linewidth}{0.4pt}

\section{二、multiset（可重复集合）}

头文件：
\begin{lstlisting}[language=C++]
#include <set>
\end{lstlisting}

\noindent\rule{\linewidth}{0.4pt}

\section{1. multiset定义}

\subsection{名称}
\begin{lstlisting}[language=C++]
multiset
\end{lstlisting}

\subsection{用法}
\begin{lstlisting}[language=C++]
multiset<int> ms;
\end{lstlisting}
例如：
\begin{lstlisting}[language=C++]
multiset<int> ms={1,2,2,3};
\end{lstlisting}
结果：
\begin{lstlisting}
1 2 2 3
\end{lstlisting}

\noindent\rule{\linewidth}{0.4pt}

\subsection{作用}
允许重复元素的set。

区别：
\begin{table}[h]
\centering
\begin{tabular}{|l|l|l|}
\hline
 & set & multiset \\ \hline
重复 & 不允许 & 允许 \\ \hline
排序 & 自动 & 自动 \\ \hline
\end{tabular}
\end{table}

\noindent\rule{\linewidth}{0.4pt}

\section{2. insert()}

\subsection{用法}
\begin{lstlisting}[language=C++]
ms.insert(x);
\end{lstlisting}
例如：
\begin{lstlisting}[language=C++]
ms.insert(2);
ms.insert(2);
\end{lstlisting}
结果：
\begin{lstlisting}
2 2
\end{lstlisting}

\noindent\rule{\linewidth}{0.4pt}

\section{3. erase()}

重点！

\subsection{删除一个}
\begin{lstlisting}[language=C++]
auto it=ms.find(2);

ms.erase(it);
\end{lstlisting}
删除一个2。

\noindent\rule{\linewidth}{0.4pt}

\subsection{删除全部}
\begin{lstlisting}[language=C++]
ms.erase(2);
\end{lstlisting}
所有2都会删除。

\noindent\rule{\linewidth}{0.4pt}

\subsection{作用}
删除元素。

\noindent\rule{\linewidth}{0.4pt}

\section{4. count()}

\subsection{用法}
\begin{lstlisting}[language=C++]
ms.count(x);
\end{lstlisting}
例如：
\begin{lstlisting}[language=C++]
multiset<int> ms={1,2,2,2,3};

cout<<ms.count(2);
\end{lstlisting}
输出：
\begin{lstlisting}
3
\end{lstlisting}

\subsection{作用}
统计出现次数。

\noindent\rule{\linewidth}{0.4pt}

\section{三、map（映射）}

头文件：
\begin{lstlisting}[language=C++]
#include <map>
\end{lstlisting}

\noindent\rule{\linewidth}{0.4pt}

\section{1. map定义}

\subsection{名称}
\begin{lstlisting}[language=C++]
map
\end{lstlisting}

\subsection{用法}
\begin{lstlisting}[language=C++]
map<string,int> mp;
\end{lstlisting}
表示：
\begin{lstlisting}
字符串 -> 整数
\end{lstlisting}
例如：
\begin{lstlisting}[language=C++]
mp["apple"]=5;
\end{lstlisting}
形成：
\begin{lstlisting}
apple 5
\end{lstlisting}

\noindent\rule{\linewidth}{0.4pt}

\subsection{作用}
键值映射。

类似：\\
字典：
\begin{lstlisting}
单词 → 解释
姓名 → 成绩
\end{lstlisting}

\noindent\rule{\linewidth}{0.4pt}

\section{2. 插入元素}

\subsection{方法1：下标}
\begin{lstlisting}[language=C++]
mp[key]=value;
\end{lstlisting}
例如：
\begin{lstlisting}[language=C++]
map<string,int> mp;

mp["Tom"]=90;
\end{lstlisting}

\noindent\rule{\linewidth}{0.4pt}

\subsection{方法2：insert}
\begin{lstlisting}[language=C++]
mp.insert(pair<string,int>("Tom",90));
\end{lstlisting}
或者：
\begin{lstlisting}[language=C++]
mp.insert({"Tom",90});
\end{lstlisting}

\noindent\rule{\linewidth}{0.4pt}

\subsection{作用}
添加键值对。

\noindent\rule{\linewidth}{0.4pt}

\section{3. 访问元素}

\subsection{名称}
\begin{lstlisting}[language=C++]
[]
\end{lstlisting}

\subsection{用法}
\begin{lstlisting}[language=C++]
mp[key];
\end{lstlisting}
例如：
\begin{lstlisting}[language=C++]
cout<<mp["Tom"];
\end{lstlisting}
输出：
\begin{lstlisting}
90
\end{lstlisting}

\noindent\rule{\linewidth}{0.4pt}

注意：\\
如果key不存在：
\begin{lstlisting}[language=C++]
mp["Bob"];
\end{lstlisting}
会自动创建：
\begin{lstlisting}
Bob 0
\end{lstlisting}

\noindent\rule{\linewidth}{0.4pt}

\section{4. erase()}

\subsection{用法}
删除key：
\begin{lstlisting}[language=C++]
mp.erase(key);
\end{lstlisting}
例如：
\begin{lstlisting}[language=C++]
mp.erase("Tom");
\end{lstlisting}
删除：
\begin{lstlisting}
Tom
\end{lstlisting}

\noindent\rule{\linewidth}{0.4pt}

\subsection{作用}
删除映射。

\noindent\rule{\linewidth}{0.4pt}

\section{5. find()}

\subsection{用法}
\begin{lstlisting}[language=C++]
mp.find(key);
\end{lstlisting}
例如：
\begin{lstlisting}[language=C++]
if(mp.find("Tom")!=mp.end())
{
    cout<<"存在";
}
\end{lstlisting}

\subsection{作用}
查找key。

\noindent\rule{\linewidth}{0.4pt}

\section{6. count()}

\subsection{用法}
\begin{lstlisting}[language=C++]
mp.count(key);
\end{lstlisting}
例如：
\begin{lstlisting}[language=C++]
mp.count("Tom");
\end{lstlisting}
返回：
\begin{lstlisting}
1
\end{lstlisting}
不存在：
\begin{lstlisting}
0
\end{lstlisting}

\noindent\rule{\linewidth}{0.4pt}

\section{7. size()}

\subsection{用法}
\begin{lstlisting}[language=C++]
mp.size();
\end{lstlisting}
作用：\\
键值对数量。

\noindent\rule{\linewidth}{0.4pt}

\section{8. 遍历map（重点）}

\subsection{名称}
\begin{lstlisting}[language=C++]
pair
\end{lstlisting}

\subsection{用法}
\begin{lstlisting}[language=C++]
for(auto p:mp)
{
    cout<<p.first<<" "<<p.second;
}
\end{lstlisting}
例如：\\
map：
\begin{lstlisting}
Tom 90
Bob 80
\end{lstlisting}
输出：
\begin{lstlisting}
Tom 90
Bob 80
\end{lstlisting}

\noindent\rule{\linewidth}{0.4pt}

其中：
\begin{lstlisting}[language=C++]
p.first
\end{lstlisting}
表示：\\
key

\begin{lstlisting}[language=C++]
p.second
\end{lstlisting}
表示：\\
value

\noindent\rule{\linewidth}{0.4pt}

\section{9. map自动排序}

默认：\\
按照key升序。

例如：
\begin{lstlisting}[language=C++]
map<int,string> mp;

mp[3]="c";
mp[1]="a";
mp[2]="b";
\end{lstlisting}
遍历：
\begin{lstlisting}
1 a
2 b
3 c
\end{lstlisting}

\noindent\rule{\linewidth}{0.4pt}

\section{map经典用途}

\subsection{统计次数}
例如：\\
统计字符串字符：
\begin{lstlisting}[language=C++]
string s="hello";

map<char,int> mp;


for(char c:s)
{
    mp[c]++;
}
\end{lstlisting}
结果：
\begin{lstlisting}
h 1
e 1
l 2
o 1
\end{lstlisting}

\noindent\rule{\linewidth}{0.4pt}

\section{四、unordered\_map（哈希表）}

头文件：
\begin{lstlisting}[language=C++]
#include <unordered_map>
\end{lstlisting}

\noindent\rule{\linewidth}{0.4pt}

\section{1. unordered\_map定义}

\subsection{名称}
\begin{lstlisting}[language=C++]
unordered_map
\end{lstlisting}

\subsection{用法}
\begin{lstlisting}[language=C++]
unordered_map<string,int> mp;
\end{lstlisting}

\noindent\rule{\linewidth}{0.4pt}

\subsection{作用}
类似map。

区别：
\begin{table}[h]
\centering
\begin{tabular}{|l|l|l|}
\hline
 & map & unordered\_map \\ \hline
底层 & 红黑树 & 哈希表 \\ \hline
排序 & 自动排序 & 无序 \\ \hline
速度 & O(logn) & 平均O(1) \\ \hline
\end{tabular}
\end{table}

\noindent\rule{\linewidth}{0.4pt}

\section{2. 插入}
\begin{lstlisting}[language=C++]
mp[key]=value;
\end{lstlisting}
例如：
\begin{lstlisting}[language=C++]
mp["apple"]=10;
\end{lstlisting}

\noindent\rule{\linewidth}{0.4pt}

\section{3. 查找}
\begin{lstlisting}[language=C++]
mp.find(key);
\end{lstlisting}
例如：
\begin{lstlisting}[language=C++]
if(mp.find("apple")!=mp.end())
{
    cout<<"yes";
}
\end{lstlisting}

\noindent\rule{\linewidth}{0.4pt}

\section{4. 统计次数}
最常用：
\begin{lstlisting}[language=C++]
unordered_map<char,int> cnt;


for(char c:s)
{
    cnt[c]++;
}
\end{lstlisting}

\noindent\rule{\linewidth}{0.4pt}

\section{map和unordered\_map选择}

考试：\\
需要排序：\\
↓
\begin{lstlisting}[language=C++]
map
\end{lstlisting}
只需要快速查找：\\
↓
\begin{lstlisting}[language=C++]
unordered_map
\end{lstlisting}
统计次数：\\
↓
\begin{lstlisting}[language=C++]
unordered_map
\end{lstlisting}

\noindent\rule{\linewidth}{0.4pt}

\section{第六章 pair / tuple / iterator（迭代器）}

这一章属于 STL 基础连接部分。

很多 STL 函数：
\begin{lstlisting}[language=C++]
sort()
find()
lower_bound()
map遍历
\end{lstlisting}
都会用到。

\noindent\rule{\linewidth}{0.4pt}

\section{一、pair（存两个数据）}

头文件：
\begin{lstlisting}[language=C++]
#include <utility>
\end{lstlisting}
通常：
\begin{lstlisting}[language=C++]
#include <bits/stdc++.h>
\end{lstlisting}
已经包含。

\noindent\rule{\linewidth}{0.4pt}

\section{1. pair定义}

\subsection{名称}
\begin{lstlisting}[language=C++]
pair
\end{lstlisting}

\subsection{用法}
格式：
\begin{lstlisting}[language=C++]
pair<类型1,类型2> 名称;
\end{lstlisting}
例如：
\begin{lstlisting}[language=C++]
pair<int,string> p;
\end{lstlisting}
表示：\\
存：
\begin{lstlisting}
整数 + 字符串
\end{lstlisting}

\noindent\rule{\linewidth}{0.4pt}

初始化：
\begin{lstlisting}[language=C++]
pair<int,string> p={1,"Tom"};
\end{lstlisting}
或者：
\begin{lstlisting}[language=C++]
pair<int,string> p;

p.first=1;
p.second="Tom";
\end{lstlisting}

\noindent\rule{\linewidth}{0.4pt}

\subsection{作用}
把两个变量绑定在一起。

例如：\\
学生：
\begin{lstlisting}
学号 + 姓名
\end{lstlisting}
坐标：
\begin{lstlisting}
x + y
\end{lstlisting}

\noindent\rule{\linewidth}{0.4pt}

\section{2. first}

\subsection{名称}
\begin{lstlisting}[language=C++]
first
\end{lstlisting}

\subsection{用法}
\begin{lstlisting}[language=C++]
p.first;
\end{lstlisting}
例如：
\begin{lstlisting}[language=C++]
pair<int,string> p={100,"Tom"};

cout<<p.first;
\end{lstlisting}
输出：
\begin{lstlisting}
100
\end{lstlisting}

\subsection{作用}
访问pair第一个元素。

\noindent\rule{\linewidth}{0.4pt}

\section{3. second}

\subsection{名称}
\begin{lstlisting}[language=C++]
second
\end{lstlisting}

\subsection{用法}
\begin{lstlisting}[language=C++]
p.second;
\end{lstlisting}
例如：
\begin{lstlisting}[language=C++]
cout<<p.second;
\end{lstlisting}
输出：
\begin{lstlisting}
Tom
\end{lstlisting}

\subsection{作用}
访问第二个元素。

\noindent\rule{\linewidth}{0.4pt}

\section{4. make\_pair()}

\subsection{名称}
\begin{lstlisting}[language=C++]
make_pair()
\end{lstlisting}

\subsection{用法}
格式：
\begin{lstlisting}[language=C++]
make_pair(数据1,数据2);
\end{lstlisting}
例如：
\begin{lstlisting}[language=C++]
auto p=make_pair(10,"hello");
\end{lstlisting}
等价：
\begin{lstlisting}[language=C++]
pair<int,string> p={10,"hello"};
\end{lstlisting}

\noindent\rule{\linewidth}{0.4pt}

\subsection{作用}
快速创建pair。

\noindent\rule{\linewidth}{0.4pt}

\section{5. pair比较}

pair可以直接比较。

例如：
\begin{lstlisting}[language=C++]
pair<int,int> a={1,2};

pair<int,int> b={1,3};


cout<<(a<b);
\end{lstlisting}
比较规则：\\
先比较：
\begin{lstlisting}
first
\end{lstlisting}
如果相同：\\
比较：
\begin{lstlisting}
second
\end{lstlisting}
例如：
\begin{lstlisting}
(1,2)
(1,3)
\end{lstlisting}
因为：
\begin{lstlisting}
2 < 3
\end{lstlisting}
所以：
\begin{lstlisting}
a < b
\end{lstlisting}

\noindent\rule{\linewidth}{0.4pt}

\subsection{作用}
常用于排序。

\noindent\rule{\linewidth}{0.4pt}

\section{6. pair排序}

例如：
\begin{lstlisting}[language=C++]
vector<pair<int,int>> v;

v.push_back({2,5});
v.push_back({1,3});
v.push_back({2,1});


sort(v.begin(),v.end());
\end{lstlisting}
结果：
\begin{lstlisting}
(1,3)
(2,1)
(2,5)
\end{lstlisting}
规则：\\
第一关键字升序。\\
第二关键字升序。

\noindent\rule{\linewidth}{0.4pt}

\section{7. 自定义pair排序}

例如：\\
按照second排序：
\begin{lstlisting}[language=C++]
bool cmp(pair<int,int>a,pair<int,int>b)
{
    return a.second<b.second;
}


sort(v.begin(),v.end(),cmp);
\end{lstlisting}
作用：\\
改变排序规则。

\noindent\rule{\linewidth}{0.4pt}

\section{二、tuple（三个及以上数据）}

头文件：
\begin{lstlisting}[language=C++]
#include <tuple>
\end{lstlisting}

\noindent\rule{\linewidth}{0.4pt}

\section{1. tuple定义}

\subsection{名称}
\begin{lstlisting}[language=C++]
tuple
\end{lstlisting}

\subsection{用法}
\begin{lstlisting}[language=C++]
tuple<int,string,double> t;
\end{lstlisting}
表示：
\begin{lstlisting}
int
string
double
\end{lstlisting}
三个数据。

\noindent\rule{\linewidth}{0.4pt}

初始化：
\begin{lstlisting}[language=C++]
tuple<int,string,double> t
=
{
1,
"Tom",
99.5
};
\end{lstlisting}

\noindent\rule{\linewidth}{0.4pt}

\subsection{作用}
存储多个不同类型的数据。

pair只能两个。\\
tuple可以多个。

\noindent\rule{\linewidth}{0.4pt}

\section{2. get()}

\subsection{名称}
\begin{lstlisting}[language=C++]
get<>
\end{lstlisting}

\subsection{用法}
格式：
\begin{lstlisting}[language=C++]
get<编号>(tuple);
\end{lstlisting}
例如：
\begin{lstlisting}[language=C++]
tuple<int,string,double> t={1,"Tom",99.5};


cout<<get<0>(t);
\end{lstlisting}
输出：
\begin{lstlisting}
1
\end{lstlisting}

\begin{lstlisting}[language=C++]
cout<<get<1>(t);
\end{lstlisting}
输出：
\begin{lstlisting}
Tom
\end{lstlisting}

\noindent\rule{\linewidth}{0.4pt}

注意：\\
编号从0开始。
\begin{lstlisting}
get<0>
第一个

get<1>
第二个

get<2>
第三个
\end{lstlisting}

\noindent\rule{\linewidth}{0.4pt}

\subsection{作用}
访问tuple元素。

\noindent\rule{\linewidth}{0.4pt}

\section{三、iterator（迭代器）}

iterator 是 STL 的核心。

简单理解：\\
迭代器 = 指向容器元素的指针。

类似：\\
C语言：
\begin{lstlisting}[language=C++]
int *p;
\end{lstlisting}
STL：
\begin{lstlisting}[language=C++]
vector<int>::iterator it;
\end{lstlisting}

\noindent\rule{\linewidth}{0.4pt}

\section{1. iterator定义}

\subsection{名称}
\begin{lstlisting}[language=C++]
iterator
\end{lstlisting}

\subsection{用法}
vector：
\begin{lstlisting}[language=C++]
vector<int>::iterator it;
\end{lstlisting}
例如：
\begin{lstlisting}[language=C++]
vector<int> v={1,2,3};

vector<int>::iterator it;
\end{lstlisting}

\noindent\rule{\linewidth}{0.4pt}

现在推荐：
\begin{lstlisting}[language=C++]
auto it;
\end{lstlisting}
例如：
\begin{lstlisting}[language=C++]
auto it=v.begin();
\end{lstlisting}

\noindent\rule{\linewidth}{0.4pt}

\subsection{作用}
遍历容器。

\noindent\rule{\linewidth}{0.4pt}

\section{2. begin()}

\subsection{名称}
\begin{lstlisting}[language=C++]
begin()
\end{lstlisting}

\subsection{用法}
\begin{lstlisting}[language=C++]
auto it=v.begin();
\end{lstlisting}
例如：
\begin{lstlisting}[language=C++]
vector<int> v={1,2,3};

auto it=v.begin();

cout<<*it;
\end{lstlisting}
输出：
\begin{lstlisting}
1
\end{lstlisting}

\noindent\rule{\linewidth}{0.4pt}

\subsection{作用}
返回第一个元素位置。

\noindent\rule{\linewidth}{0.4pt}

\section{3. end()}

\subsection{名称}
\begin{lstlisting}[language=C++]
end()
\end{lstlisting}

\subsection{用法}
\begin{lstlisting}[language=C++]
auto it=v.end();
\end{lstlisting}
注意：\\
end不是最后元素。

它指向：
\begin{lstlisting}
最后一个元素后面
\end{lstlisting}
所以：\\
错误：
\begin{lstlisting}[language=C++]
cout<<*v.end();
\end{lstlisting}
正确：
\begin{lstlisting}[language=C++]
cout<<*(v.end()-1);
\end{lstlisting}

\noindent\rule{\linewidth}{0.4pt}

\subsection{作用}
表示遍历结束。

\noindent\rule{\linewidth}{0.4pt}

\section{4. 解引用 *}

\subsection{名称}
\begin{lstlisting}[language=C++]
*
\end{lstlisting}

\subsection{用法}
\begin{lstlisting}[language=C++]
*iterator
\end{lstlisting}
例如：
\begin{lstlisting}[language=C++]
auto it=v.begin();

cout<<*it;
\end{lstlisting}
输出：
\begin{lstlisting}
第一个元素
\end{lstlisting}

\noindent\rule{\linewidth}{0.4pt}

\subsection{作用}
获得迭代器指向的数据。

\noindent\rule{\linewidth}{0.4pt}

\section{5. iterator移动}

\subsection{名称}
++

--

\subsection{用法}
向后：
\begin{lstlisting}[language=C++]
it++;
\end{lstlisting}
向前：
\begin{lstlisting}[language=C++]
it--;
\end{lstlisting}
例如：
\begin{lstlisting}[language=C++]
auto it=v.begin();

it++;

cout<<*it;
\end{lstlisting}
输出：\\
第二个元素。

\noindent\rule{\linewidth}{0.4pt}

\section{6. iterator遍历vector}

\subsection{方法}
\begin{lstlisting}[language=C++]
for(auto it=v.begin();
    it!=v.end();
    it++)
{
    cout<<*it;
}
\end{lstlisting}
输出：
\begin{lstlisting}
123
\end{lstlisting}

\noindent\rule{\linewidth}{0.4pt}

\section{7. iterator遍历set}

set不能：
\begin{lstlisting}[language=C++]
s[0]
\end{lstlisting}
只能：
\begin{lstlisting}[language=C++]
for(auto it=s.begin();
    it!=s.end();
    it++)
{
    cout<<*it;
}
\end{lstlisting}

\noindent\rule{\linewidth}{0.4pt}

\section{8. iterator遍历map（重点）}

map里面：\\
元素类型：
\begin{lstlisting}[language=C++]
pair
\end{lstlisting}
例如：
\begin{lstlisting}[language=C++]
map<string,int> mp;


mp["Tom"]=90;
mp["Bob"]=80;


for(auto it=mp.begin();
    it!=mp.end();
    it++)
{
    cout<<it->first;
    cout<<it->second;
}
\end{lstlisting}
注意：\\
这里：
\begin{lstlisting}[language=C++]
it->first
\end{lstlisting}
等价：
\begin{lstlisting}[language=C++]
(*it).first
\end{lstlisting}

\noindent\rule{\linewidth}{0.4pt}

\section{9. 箭头符号 ->}

\subsection{名称}
成员访问

\subsection{用法}
如果：
\begin{lstlisting}[language=C++]
it
\end{lstlisting}
是指针：\\
访问成员：
\begin{lstlisting}[language=C++]
it->成员
\end{lstlisting}
例如：\\
map：
\begin{lstlisting}[language=C++]
it->first
it->second
\end{lstlisting}

\noindent\rule{\linewidth}{0.4pt}

\subsection{作用}
访问指针对象成员。

\noindent\rule{\linewidth}{0.4pt}

\section{10. const\_iterator}

\subsection{名称}
\begin{lstlisting}[language=C++]
const_iterator
\end{lstlisting}

\subsection{用法}
\begin{lstlisting}[language=C++]
vector<int>::const_iterator it;
\end{lstlisting}
或者：
\begin{lstlisting}[language=C++]
auto it=v.cbegin();
\end{lstlisting}

\noindent\rule{\linewidth}{0.4pt}

\subsection{作用}
只读迭代器。

不能：
\begin{lstlisting}[language=C++]
*it=10;
\end{lstlisting}

\noindent\rule{\linewidth}{0.4pt}

\section{11. reverse\_iterator}

\subsection{名称}
反向迭代器

\subsection{用法}
\begin{lstlisting}[language=C++]
auto it=v.rbegin();
\end{lstlisting}
遍历：
\begin{lstlisting}[language=C++]
for(auto it=v.rbegin();
    it!=v.rend();
    it++)
{
    cout<<*it;
}
\end{lstlisting}
例如：\\
vector:
\begin{lstlisting}
1 2 3
\end{lstlisting}
输出：
\begin{lstlisting}
321
\end{lstlisting}

\noindent\rule{\linewidth}{0.4pt}

\subsection{作用}
反向遍历。

\noindent\rule{\linewidth}{0.4pt}

\section{四、auto关键字（STL必用）}

\subsection{名称}
\begin{lstlisting}[language=C++]
auto
\end{lstlisting}

\subsection{用法}
不用写复杂类型：

传统：
\begin{lstlisting}[language=C++]
vector<int>::iterator it;
\end{lstlisting}
auto：
\begin{lstlisting}[language=C++]
auto it=v.begin();
\end{lstlisting}

\noindent\rule{\linewidth}{0.4pt}

map：

传统：
\begin{lstlisting}[language=C++]
map<string,int>::iterator it;
\end{lstlisting}
auto：
\begin{lstlisting}[language=C++]
auto it=mp.begin();
\end{lstlisting}

\noindent\rule{\linewidth}{0.4pt}

\subsection{作用}
让编译器自动推断类型。\\
考试非常推荐。

\noindent\rule{\linewidth}{0.4pt}

\section{五、范围for循环（C++11）}

\subsection{名称}
range-based for

\subsection{用法}
vector：
\begin{lstlisting}[language=C++]
for(auto x:v)
{
    cout<<x;
}
\end{lstlisting}
string：
\begin{lstlisting}[language=C++]
for(char c:s)
{
    cout<<c;
}
\end{lstlisting}
map：
\begin{lstlisting}[language=C++]
for(auto p:mp)
{
    cout<<p.first;
    cout<<p.second;
}
\end{lstlisting}

\noindent\rule{\linewidth}{0.4pt}

\subsection{作用}
快速遍历。

\noindent\rule{\linewidth}{0.4pt}

\section{六、STL常用头文件总表（考试）}

\begin{table}[h]
\centering
\begin{tabular}{|l|l|}
\hline
功能 & 头文件 \\ \hline
vector & string \\ \hline
stack & \texttt{<stack>} \\ \hline
queue & \texttt{<queue>} \\ \hline
priority\_queue & \texttt{<queue>} \\ \hline
set/map & \texttt{<set>} \texttt{<map>} \\ \hline
算法 & \texttt{<algorithm>} \\ \hline
累加 & \texttt{<numeric>} \\ \hline
pair & \texttt{<utility>} \\ \hline
tuple & \texttt{<tuple>} \\ \hline
\end{tabular}
\end{table}

\noindent\rule{\linewidth}{0.4pt}

\section{第七章 C++ STL上机考试万能模板}

这一部分不是介绍函数，而是整理\textbf{考试直接能套用的代码模板}。

\noindent\rule{\linewidth}{0.4pt}

\section{1. 万能头文件模板}

\subsection{名称}
\begin{lstlisting}[language=C++]
#include <bits/stdc++.h>
\end{lstlisting}

\subsection{用法}
\begin{lstlisting}[language=C++]
#include <bits/stdc++.h>
using namespace std;

int main()
{


    return 0;
}
\end{lstlisting}

\subsection{作用}
一次包含几乎所有 STL：

包括：
\begin{lstlisting}[language=C++]
vector
string
map
set
queue
stack
algorithm
numeric
\end{lstlisting}
等。

考试最推荐。

\noindent\rule{\linewidth}{0.4pt}

\section{2. 快速输入输出}

\subsection{名称}
\begin{lstlisting}[language=C++]
ios::sync_with_stdio(false);
cin.tie(nullptr);
\end{lstlisting}

\subsection{用法}
\begin{lstlisting}[language=C++]
int main()
{
    ios::sync_with_stdio(false);
    cin.tie(nullptr);


}
\end{lstlisting}

\subsection{作用}
关闭 C 和 C++ 输入同步，提高速度。\\
普通题可以不写。\\
数据量大建议写。

\noindent\rule{\linewidth}{0.4pt}

\section{3. vector输入模板}

\subsection{场景}
输入n个数字。

\subsection{写法}
\begin{lstlisting}[language=C++]
int n;

cin>>n;


vector<int> v(n);


for(int i=0;i<n;i++)
{
    cin>>v[i];
}
\end{lstlisting}

\subsection{作用}
快速创建数组。

\noindent\rule{\linewidth}{0.4pt}

另一种：
\begin{lstlisting}[language=C++]
vector<int> v;


for(int i=0;i<n;i++)
{
    int x;

    cin>>x;

    v.push_back(x);
}
\end{lstlisting}

\noindent\rule{\linewidth}{0.4pt}

区别：
\begin{table}[h]
\centering
\begin{tabular}{|l|l|}
\hline
方法 & 特点 \\ \hline
vector(n) & 提前知道数量 \\ \hline
push\_back & 不知道数量 \\ \hline
\end{tabular}
\end{table}

\noindent\rule{\linewidth}{0.4pt}

\section{4. vector遍历模板}

\subsection{方法1：下标}
\begin{lstlisting}[language=C++]
for(int i=0;i<v.size();i++)
{
    cout<<v[i]<<" ";
}
\end{lstlisting}
适合：\\
需要下标。

\noindent\rule{\linewidth}{0.4pt}

\subsection{方法2：范围for}
\begin{lstlisting}[language=C++]
for(auto x:v)
{
    cout<<x<<" ";
}
\end{lstlisting}
适合：\\
只读取。

\noindent\rule{\linewidth}{0.4pt}

\subsection{方法3：引用}
\begin{lstlisting}[language=C++]
for(auto &x:v)
{
    cin>>x;
}
\end{lstlisting}
作用：\\
直接修改原元素。

例如：
\begin{lstlisting}[language=C++]
for(auto &x:v)
{
    x*=2;
}
\end{lstlisting}

\noindent\rule{\linewidth}{0.4pt}

\section{5. 数组排序模板}

\subsection{升序}
\begin{lstlisting}[language=C++]
sort(a,a+n);
\end{lstlisting}
例如：
\begin{lstlisting}[language=C++]
int a[5]={3,1,5,2,4};

sort(a,a+5);
\end{lstlisting}
结果：
\begin{lstlisting}
1 2 3 4 5
\end{lstlisting}

\noindent\rule{\linewidth}{0.4pt}

\subsection{vector排序}
\begin{lstlisting}[language=C++]
sort(v.begin(),v.end());
\end{lstlisting}

\noindent\rule{\linewidth}{0.4pt}

\subsection{降序}
数组：
\begin{lstlisting}[language=C++]
sort(a,a+n,greater<int>());
\end{lstlisting}
vector：
\begin{lstlisting}[language=C++]
sort(v.begin(),v.end(),greater<int>());
\end{lstlisting}

\noindent\rule{\linewidth}{0.4pt}

\section{6. 字符串排序模板}

\subsection{名称}
字符排序

\subsection{用法}
\begin{lstlisting}[language=C++]
string s="dcba";


sort(s.begin(),s.end());
\end{lstlisting}
结果：
\begin{lstlisting}
abcd
\end{lstlisting}

\noindent\rule{\linewidth}{0.4pt}

降序：
\begin{lstlisting}[language=C++]
sort(s.begin(),s.end(),greater<char>());
\end{lstlisting}
结果：
\begin{lstlisting}
dcba
\end{lstlisting}

\noindent\rule{\linewidth}{0.4pt}

\section{7. vector求最大最小值}

\subsection{最大值}
\begin{lstlisting}[language=C++]
int mx=*max_element(v.begin(),v.end());
\end{lstlisting}

\noindent\rule{\linewidth}{0.4pt}

\subsection{最小值}
\begin{lstlisting}[language=C++]
int mn=*min_element(v.begin(),v.end());
\end{lstlisting}

\noindent\rule{\linewidth}{0.4pt}

\subsection{最大值位置}
\begin{lstlisting}[language=C++]
int pos=max_element(v.begin(),v.end())-v.begin();
\end{lstlisting}
例如：
\begin{lstlisting}[language=C++]
v={3,7,2}
\end{lstlisting}
返回：
\begin{lstlisting}
1
\end{lstlisting}

\noindent\rule{\linewidth}{0.4pt}

\section{8. vector求和模板}

\subsection{名称}
accumulate

\subsection{用法}
\begin{lstlisting}[language=C++]
int sum=accumulate(
    v.begin(),
    v.end(),
    0
);
\end{lstlisting}
例如：
\begin{lstlisting}[language=C++]
v={1,2,3}
\end{lstlisting}
结果：
\begin{lstlisting}
6
\end{lstlisting}

\noindent\rule{\linewidth}{0.4pt}

\section{9. 去重模板}

\subsection{场景}
数组：
\begin{lstlisting}
1 2 2 3 3
\end{lstlisting}
变：
\begin{lstlisting}
1 2 3
\end{lstlisting}

\subsection{写法}
\begin{lstlisting}[language=C++]
sort(v.begin(),v.end());


v.erase(
    unique(v.begin(),v.end()),
    v.end()
);
\end{lstlisting}
步骤：\\
第一步：\\
排序：
\begin{lstlisting}
1 2 2 3 3
\end{lstlisting}
第二步：\\
unique：
\begin{lstlisting}
1 2 3 ? ?
\end{lstlisting}
第三步：\\
erase删除。

\noindent\rule{\linewidth}{0.4pt}

\section{10. 查找元素模板}

\subsection{方法1：find}
\begin{lstlisting}[language=C++]
auto it=find(
    v.begin(),
    v.end(),
    x
);


if(it!=v.end())
{
    cout<<"找到";
}
\end{lstlisting}

\noindent\rule{\linewidth}{0.4pt}

\subsection{方法2：count}
\begin{lstlisting}[language=C++]
if(count(v.begin(),v.end(),x))
{
    cout<<"存在";
}
\end{lstlisting}

\noindent\rule{\linewidth}{0.4pt}

\section{11. set去重模板}

\subsection{场景}
自动排序去重。

代码：
\begin{lstlisting}[language=C++]
set<int> s;


for(int x:v)
{
    s.insert(x);
}
\end{lstlisting}
遍历：
\begin{lstlisting}[language=C++]
for(auto x:s)
{
    cout<<x<<" ";
}
\end{lstlisting}

\noindent\rule{\linewidth}{0.4pt}

\section{12. map统计次数模板（超级高频）}

\subsection{场景}
统计：
\begin{itemize}
    \item 字符出现次数
    \item 数字出现次数
    \item 单词次数
\end{itemize}

代码：
\begin{lstlisting}[language=C++]
map<int,int> mp;


for(int x:v)
{
    mp[x]++;
}
\end{lstlisting}
例如：\\
输入：
\begin{lstlisting}
1 2 2 3
\end{lstlisting}
结果：
\begin{lstlisting}
1 -> 1次
2 -> 2次
3 -> 1次
\end{lstlisting}

\noindent\rule{\linewidth}{0.4pt}

字符串：
\begin{lstlisting}[language=C++]
string s="hello";


map<char,int> mp;


for(char c:s)
{
    mp[c]++;
}
\end{lstlisting}
结果：
\begin{lstlisting}
h 1
e 1
l 2
o 1
\end{lstlisting}

\noindent\rule{\linewidth}{0.4pt}

\section{13. unordered\_map统计模板}

推荐：
\begin{lstlisting}[language=C++]
unordered_map<int,int> mp;


for(int x:v)
{
    mp[x]++;
}
\end{lstlisting}
速度更快。

\noindent\rule{\linewidth}{0.4pt}

\section{14. map遍历模板}

\subsection{方法1}
\begin{lstlisting}[language=C++]
for(auto p:mp)
{
    cout<<p.first<<" ";

    cout<<p.second;
}
\end{lstlisting}

\noindent\rule{\linewidth}{0.4pt}

\subsection{方法2（修改）}
\begin{lstlisting}[language=C++]
for(auto &p:mp)
{
    p.second++;
}
\end{lstlisting}

\noindent\rule{\linewidth}{0.4pt}

\section{15. stack括号匹配模板}

\subsection{场景}
判断：
\begin{lstlisting}
(())
\end{lstlisting}
是否合法。

代码：
\begin{lstlisting}[language=C++]
string s;


cin>>s;


stack<char> st;


for(char c:s)
{
    if(c=='(')
    {
        st.push(c);
    }


    else
    {
        if(st.empty())
        {
            cout<<"NO";

            return 0;
        }


        st.pop();
    }
}


if(st.empty())
    cout<<"YES";

else
    cout<<"NO";
\end{lstlisting}

\noindent\rule{\linewidth}{0.4pt}

\section{16. queue BFS模板}

\subsection{场景}
图搜索：
\begin{itemize}
    \item 迷宫
    \item 最短路径
    \item 层序遍历
\end{itemize}

模板：
\begin{lstlisting}[language=C++]
queue<int> q;


q.push(start);


while(!q.empty())
{
    int x=q.front();

    q.pop();


    //处理x


}
\end{lstlisting}

\noindent\rule{\linewidth}{0.4pt}

\section{17. priority\_queue最大值模板}

\subsection{默认大根堆}
\begin{lstlisting}[language=C++]
priority_queue<int> pq;


pq.push(5);

pq.push(10);

pq.push(3);


cout<<pq.top();
\end{lstlisting}
输出：
\begin{lstlisting}
10
\end{lstlisting}

\noindent\rule{\linewidth}{0.4pt}

\section{18. priority\_queue小根堆}
\begin{lstlisting}[language=C++]
priority_queue<
    int,
    vector<int>,
    greater<int>
> pq;
\end{lstlisting}
例如：
\begin{lstlisting}[language=C++]
pq.push(5);
pq.push(2);
pq.push(8);
\end{lstlisting}
top：
\begin{lstlisting}
2
\end{lstlisting}

\noindent\rule{\linewidth}{0.4pt}

\section{19. 二分查找模板}

前提：\\
必须排序。

\noindent\rule{\linewidth}{0.4pt}

\subsection{lower\_bound}
\begin{lstlisting}[language=C++]
auto it=
lower_bound(
    v.begin(),
    v.end(),
    x
);


int pos=it-v.begin();
\end{lstlisting}
寻找：
\begin{lstlisting}
>=x
\end{lstlisting}

\noindent\rule{\linewidth}{0.4pt}

\subsection{upper\_bound}
\begin{lstlisting}[language=C++]
auto it=
upper_bound(
    v.begin(),
    v.end(),
    x
);
\end{lstlisting}
寻找：
\begin{lstlisting}
>x
\end{lstlisting}

\noindent\rule{\linewidth}{0.4pt}

\section{20. 全排列模板}

\subsection{next\_permutation}
例如：
\begin{lstlisting}[language=C++]
string s="abc";


do
{
    cout<<s<<endl;

}
while(next_permutation(
    s.begin(),
    s.end()
));
\end{lstlisting}
输出：
\begin{lstlisting}
abc
acb
bac
bca
cab
cba
\end{lstlisting}

\noindent\rule{\linewidth}{0.4pt}

\section{21. 二维vector模板}

\subsection{定义}
\begin{lstlisting}[language=C++]
vector<vector<int>> v(
    n,
    vector<int>(m)
);
\end{lstlisting}
创建：\\
n行m列。

例如：
\begin{lstlisting}[language=C++]
vector<vector<int>> a(
    3,
    vector<int>(4,0)
);
\end{lstlisting}
得到：\\
3×4 全0矩阵。

\noindent\rule{\linewidth}{0.4pt}

\section{22. 二维vector遍历}
\begin{lstlisting}[language=C++]
for(int i=0;i<v.size();i++)
{
    for(int j=0;j<v[i].size();j++)
    {
        cout<<v[i][j];
    }
}
\end{lstlisting}

\noindent\rule{\linewidth}{0.4pt}

\section{23. pair排序模板}

\subsection{默认}
\begin{lstlisting}[language=C++]
vector<pair<int,int>> v;


sort(v.begin(),v.end());
\end{lstlisting}
排序：\\
第一关键字\\
↓\\
第二关键字。

\noindent\rule{\linewidth}{0.4pt}

\subsection{按第二关键字}
\begin{lstlisting}[language=C++]
sort(
v.begin(),
v.end(),
[](pair<int,int>a,pair<int,int>b)
{
    return a.second<b.second;
});
\end{lstlisting}

\noindent\rule{\linewidth}{0.4pt}

\section{24. 常用字符串处理组合}

\subsection{删除某字符}
例如删除空格：
\begin{lstlisting}[language=C++]
s.erase(
remove(
s.begin(),
s.end(),
' '
),
s.end()
);
\end{lstlisting}

\noindent\rule{\linewidth}{0.4pt}

\subsection{字符串反转}
\begin{lstlisting}[language=C++]
reverse(
s.begin(),
s.end()
);
\end{lstlisting}

\noindent\rule{\linewidth}{0.4pt}

\subsection{字符串转数字}
\begin{lstlisting}[language=C++]
int x=stoi(s);
\end{lstlisting}

\noindent\rule{\linewidth}{0.4pt}

\subsection{数字转字符串}
\begin{lstlisting}[language=C++]
string s=to_string(x);
\end{lstlisting}

\noindent\rule{\linewidth}{0.4pt}

\section{25. STL复杂度速记}

\begin{table}[h]
\centering
\begin{tabular}{|l|l|}
\hline
操作 & 复杂度 \\ \hline
vector尾插push\_back & O(1) \\ \hline
vector中间插入 & O(n) \\ \hline
sort & O(nlogn) \\ \hline
find(vector) & O(n) \\ \hline
binary\_search & O(logn) \\ \hline
set查找 & O(logn) \\ \hline
map查找 & O(logn) \\ \hline
unordered\_map查找 & 平均O(1) \\ \hline
priority\_queue插入 & O(logn) \\ \hline
\end{tabular}
\end{table}

\noindent\rule{\linewidth}{0.4pt}

\section{26. 考试最常用TOP20函数（建议背）}
\begin{lstlisting}[language=C++]
vector:
push_back()
pop_back()
size()
erase()
insert()
clear()
begin()
end()


string:
getline()
size()
substr()
find()
erase()
replace()
append()


algorithm:
sort()
reverse()
find()
count()
accumulate()
max_element()
min_element()
lower_bound()
upper_bound()


map:
insert
[]
find
count


stack:
push
pop
top


queue:
push
pop
front


priority_queue:
push
pop
top
\end{lstlisting}

\noindent\rule{\linewidth}{0.4pt}
\section{第八章 STL补充知识（容易遗漏但考试可能出现）}

这一章主要补充：
\begin{itemize}
    \item array
    \item unordered\_set
    \item bitset
    \item numeric更多函数
    \item functional
    \item lambda表达式
    \item 自定义比较器
    \item STL常见坑
\end{itemize}

\noindent\rule{\linewidth}{0.4pt}

\section{一、array（固定大小数组）}

头文件：
\begin{lstlisting}[language=C++]
#include <array>
\end{lstlisting}

\noindent\rule{\linewidth}{0.4pt}

\section{1. array定义}

\subsection{名称}
\begin{lstlisting}[language=C++]
array
\end{lstlisting}

\subsection{用法}
格式：
\begin{lstlisting}[language=C++]
array<类型,大小> 名称;
\end{lstlisting}
例如：
\begin{lstlisting}[language=C++]
array<int,5> a;
\end{lstlisting}
表示：
\begin{lstlisting}
长度固定为5的int数组
\end{lstlisting}
初始化：
\begin{lstlisting}[language=C++]
array<int,5> a={1,2,3,4,5};
\end{lstlisting}

\noindent\rule{\linewidth}{0.4pt}

\subsection{作用}
STL版本的普通数组。

区别：\\
普通数组：
\begin{lstlisting}[language=C++]
int a[5];
\end{lstlisting}
array：
\begin{lstlisting}[language=C++]
array<int,5> a;
\end{lstlisting}
优点：
\begin{itemize}
    \item 支持STL函数
    \item 有size()
    \item 有迭代器
\end{itemize}

\noindent\rule{\linewidth}{0.4pt}

\section{2. size()}

\subsection{名称}
\begin{lstlisting}[language=C++]
size()
\end{lstlisting}

\subsection{用法}
\begin{lstlisting}[language=C++]
a.size();
\end{lstlisting}
例如：
\begin{lstlisting}[language=C++]
array<int,5> a;

cout<<a.size();
\end{lstlisting}
输出：
\begin{lstlisting}
5
\end{lstlisting}

\subsection{作用}
获取数组大小。

\noindent\rule{\linewidth}{0.4pt}

\section{3. front()}

\subsection{用法}
\begin{lstlisting}[language=C++]
a.front();
\end{lstlisting}
例如：
\begin{lstlisting}[language=C++]
cout<<a.front();
\end{lstlisting}
作用：\\
访问第一个元素。

\noindent\rule{\linewidth}{0.4pt}

\section{4. back()}

\subsection{用法}
\begin{lstlisting}[language=C++]
a.back();
\end{lstlisting}
作用：\\
访问最后元素。

\noindent\rule{\linewidth}{0.4pt}

\section{5. fill()}

\subsection{用法}
\begin{lstlisting}[language=C++]
a.fill(0);
\end{lstlisting}
例如：
\begin{lstlisting}[language=C++]
array<int,5> a;

a.fill(10);
\end{lstlisting}
结果：
\begin{lstlisting}
10 10 10 10 10
\end{lstlisting}

\subsection{作用}
全部赋值。

\noindent\rule{\linewidth}{0.4pt}

\section{6. begin()}

\subsection{用法}
\begin{lstlisting}[language=C++]
a.begin();
\end{lstlisting}
作用：\\
返回首迭代器。

\noindent\rule{\linewidth}{0.4pt}

\section{7. end()}

\subsection{用法}
\begin{lstlisting}[language=C++]
a.end();
\end{lstlisting}
作用：\\
返回尾后迭代器。

\noindent\rule{\linewidth}{0.4pt}

\section{二、unordered\_set（哈希集合）}

头文件：
\begin{lstlisting}[language=C++]
#include <unordered_set>
\end{lstlisting}

\noindent\rule{\linewidth}{0.4pt}

\section{1. unordered\_set定义}

\subsection{名称}
\begin{lstlisting}[language=C++]
unordered_set
\end{lstlisting}

\subsection{用法}
\begin{lstlisting}[language=C++]
unordered_set<int> s;
\end{lstlisting}

\noindent\rule{\linewidth}{0.4pt}

\subsection{作用}
类似set。

区别：
\begin{table}[h]
\centering
\begin{tabular}{|l|l|l|}
\hline
 & set & unordered\_set \\ \hline
底层 & 红黑树 & 哈希表 \\ \hline
排序 & 自动排序 & 无序 \\ \hline
速度 & O(logn) & 平均O(1) \\ \hline
\end{tabular}
\end{table}

\noindent\rule{\linewidth}{0.4pt}

\section{2. insert()}

\subsection{用法}
\begin{lstlisting}[language=C++]
s.insert(x);
\end{lstlisting}
例如：
\begin{lstlisting}[language=C++]
unordered_set<int> s;

s.insert(3);
s.insert(1);
\end{lstlisting}
输出顺序不固定。

\noindent\rule{\linewidth}{0.4pt}

\section{3. find()}

\subsection{用法}
\begin{lstlisting}[language=C++]
auto it=s.find(x);
\end{lstlisting}
判断：
\begin{lstlisting}[language=C++]
if(it!=s.end())
{
    cout<<"存在";
}
\end{lstlisting}

\noindent\rule{\linewidth}{0.4pt}

\section{4. count()}

\subsection{用法}
\begin{lstlisting}[language=C++]
s.count(x);
\end{lstlisting}
返回：\\
存在：
\begin{lstlisting}
1
\end{lstlisting}
不存在：
\begin{lstlisting}
0
\end{lstlisting}

\noindent\rule{\linewidth}{0.4pt}

\section{5. erase()}

\subsection{用法}
\begin{lstlisting}[language=C++]
s.erase(x);
\end{lstlisting}
作用：\\
删除元素。

\noindent\rule{\linewidth}{0.4pt}

\section{6. size()}

\subsection{用法}
\begin{lstlisting}[language=C++]
s.size();
\end{lstlisting}
作用：\\
元素数量。

\noindent\rule{\linewidth}{0.4pt}

\section{三、bitset（二进制位操作）}

头文件：
\begin{lstlisting}[language=C++]
#include <bitset>
\end{lstlisting}

\noindent\rule{\linewidth}{0.4pt}

\section{1. bitset定义}

\subsection{名称}
\begin{lstlisting}[language=C++]
bitset
\end{lstlisting}

\subsection{用法}
\begin{lstlisting}[language=C++]
bitset<8> b;
\end{lstlisting}
表示：\\
8位二进制。

初始化：
\begin{lstlisting}[language=C++]
bitset<8> b("1010");
\end{lstlisting}
结果：
\begin{lstlisting}
00001010
\end{lstlisting}

\noindent\rule{\linewidth}{0.4pt}

\subsection{作用}
处理二进制非常方便。

例如：
\begin{itemize}
    \item 位运算
    \item 状态压缩
\end{itemize}

\noindent\rule{\linewidth}{0.4pt}

\section{2. count()}

\subsection{用法}
\begin{lstlisting}[language=C++]
b.count();
\end{lstlisting}
例如：
\begin{lstlisting}[language=C++]
bitset<8> b("1011");

cout<<b.count();
\end{lstlisting}
输出：
\begin{lstlisting}
3
\end{lstlisting}

\subsection{作用}
统计1的数量。

\noindent\rule{\linewidth}{0.4pt}

\section{3. size()}

\subsection{用法}
\begin{lstlisting}[language=C++]
b.size();
\end{lstlisting}
作用：\\
位数。

\noindent\rule{\linewidth}{0.4pt}

\section{4. test()}

\subsection{名称}
\begin{lstlisting}[language=C++]
test()
\end{lstlisting}

\subsection{用法}
\begin{lstlisting}[language=C++]
b.test(pos);
\end{lstlisting}
例如：
\begin{lstlisting}[language=C++]
bitset<8> b("1010");

cout<<b.test(1);
\end{lstlisting}
作用：\\
判断第pos位是否为1。

\noindent\rule{\linewidth}{0.4pt}

\section{5. set()}

\subsection{用法}
全部置1：
\begin{lstlisting}[language=C++]
b.set();
\end{lstlisting}
指定位置：
\begin{lstlisting}[language=C++]
b.set(2);
\end{lstlisting}
作用：\\
把某位变成1。

\noindent\rule{\linewidth}{0.4pt}

\section{6. reset()}

\subsection{用法}
全部置0：
\begin{lstlisting}[language=C++]
b.reset();
\end{lstlisting}
指定位置：
\begin{lstlisting}[language=C++]
b.reset(2);
\end{lstlisting}
作用：\\
把某位变成0。

\noindent\rule{\linewidth}{0.4pt}

\section{7. flip()}

\subsection{用法}
全部取反：
\begin{lstlisting}[language=C++]
b.flip();
\end{lstlisting}
指定：
\begin{lstlisting}[language=C++]
b.flip(2);
\end{lstlisting}
例如：\\
0→1\\
1→0

\noindent\rule{\linewidth}{0.4pt}

\section{8. to\_string()}

\subsection{用法}
\begin{lstlisting}[language=C++]
string s=b.to_string();
\end{lstlisting}
作用：\\
转换成字符串。

\noindent\rule{\linewidth}{0.4pt}

\section{9. to\_ulong()}

\subsection{用法}
\begin{lstlisting}[language=C++]
unsigned long x=b.to_ulong();
\end{lstlisting}
作用：\\
转换成整数。

\noindent\rule{\linewidth}{0.4pt}

\section{四、numeric头文件其他函数}

头文件：
\begin{lstlisting}[language=C++]
#include <numeric>
\end{lstlisting}

\noindent\rule{\linewidth}{0.4pt}

\section{1. accumulate()}

前面讲过：
\begin{lstlisting}[language=C++]
accumulate(
begin,
end,
初始值
);
\end{lstlisting}
作用：\\
求和。

\noindent\rule{\linewidth}{0.4pt}

\section{2. iota()}

\subsection{名称}
\begin{lstlisting}[language=C++]
iota()
\end{lstlisting}

\subsection{用法}
\begin{lstlisting}[language=C++]
iota(
begin,
end,
初始值
);
\end{lstlisting}
例如：
\begin{lstlisting}[language=C++]
vector<int> v(5);


iota(
v.begin(),
v.end(),
1
);
\end{lstlisting}
结果：
\begin{lstlisting}
1 2 3 4 5
\end{lstlisting}

\noindent\rule{\linewidth}{0.4pt}

\subsection{作用}
连续赋值。

类似：
\begin{lstlisting}[language=C++]
for循环赋值
\end{lstlisting}

\noindent\rule{\linewidth}{0.4pt}

\section{3. inner\_product()}

\subsection{名称}
\begin{lstlisting}[language=C++]
inner_product()
\end{lstlisting}

\subsection{用法}
\begin{lstlisting}[language=C++]
inner_product(
a.begin(),
a.end(),
b.begin(),
0
);
\end{lstlisting}
例如：\\
a:
\begin{lstlisting}
1 2 3
\end{lstlisting}
b:
\begin{lstlisting}
4 5 6
\end{lstlisting}
计算：
\begin{lstlisting}
1*4+2*5+3*6
\end{lstlisting}
结果：
\begin{lstlisting}
32
\end{lstlisting}

\noindent\rule{\linewidth}{0.4pt}

\subsection{作用}
计算向量点积。

\noindent\rule{\linewidth}{0.4pt}

\section{4. partial\_sum()}

\subsection{名称}
\begin{lstlisting}[language=C++]
partial_sum()
\end{lstlisting}

\subsection{用法}
\begin{lstlisting}[language=C++]
partial_sum(
begin,
end,
输出位置
);
\end{lstlisting}
例如：
\begin{lstlisting}[language=C++]
vector<int> v={1,2,3,4};

vector<int> ans(4);


partial_sum(
v.begin(),
v.end(),
ans.begin()
);
\end{lstlisting}
结果：
\begin{lstlisting}
1 3 6 10
\end{lstlisting}

\noindent\rule{\linewidth}{0.4pt}

\subsection{作用}
前缀和。

\noindent\rule{\linewidth}{0.4pt}

\section{5. adjacent\_difference()}

\subsection{名称}
\begin{lstlisting}[language=C++]
adjacent_difference()
\end{lstlisting}

\subsection{用法}
\begin{lstlisting}[language=C++]
adjacent_difference(
begin,
end,
输出位置
);
\end{lstlisting}
例如：\\
输入：
\begin{lstlisting}
1 3 6 10
\end{lstlisting}
输出：
\begin{lstlisting}
1 2 3 4
\end{lstlisting}

\noindent\rule{\linewidth}{0.4pt}

\subsection{作用}
求相邻差。

\noindent\rule{\linewidth}{0.4pt}

\section{五、functional（函数对象）}

头文件：
\begin{lstlisting}[language=C++]
#include <functional>
\end{lstlisting}
常用于：\\
sort第三个参数。

\noindent\rule{\linewidth}{0.4pt}

\section{1. greater()}

\subsection{名称}
\begin{lstlisting}[language=C++]
greater<>
\end{lstlisting}

\subsection{用法}
降序：
\begin{lstlisting}[language=C++]
sort(
v.begin(),
v.end(),
greater<int>()
);
\end{lstlisting}
作用：\\
从大到小排序。

\noindent\rule{\linewidth}{0.4pt}

\section{2. less()}

默认：
\begin{lstlisting}[language=C++]
less<int>()
\end{lstlisting}
用法：
\begin{lstlisting}[language=C++]
sort(
v.begin(),
v.end(),
less<int>()
);
\end{lstlisting}
作用：\\
升序排序。

\noindent\rule{\linewidth}{0.4pt}

\section{3. plus()}

\subsection{用法}
\begin{lstlisting}[language=C++]
plus<int>()
\end{lstlisting}
作用：\\
加法。

\noindent\rule{\linewidth}{0.4pt}

\section{4. minus()}
\begin{lstlisting}[language=C++]
minus<int>()
\end{lstlisting}
作用：\\
减法。

\noindent\rule{\linewidth}{0.4pt}

\section{5. multiplies()}
\begin{lstlisting}[language=C++]
multiplies<int>()
\end{lstlisting}
作用：\\
乘法。

\noindent\rule{\linewidth}{0.4pt}

\section{6. divides()}
\begin{lstlisting}[language=C++]
divides<int>()
\end{lstlisting}
作用：\\
除法。

\noindent\rule{\linewidth}{0.4pt}

\section{六、lambda表达式（非常重要）}

C++11以后：\\
可以直接写匿名函数。

\noindent\rule{\linewidth}{0.4pt}

\section{1. 基本格式}
\begin{lstlisting}[language=C++]
[捕获](参数)
{
    函数体
}
\end{lstlisting}

\noindent\rule{\linewidth}{0.4pt}

\section{2. sort中使用（最常见）}
例如：\\
按照升序：
\begin{lstlisting}[language=C++]
sort(
v.begin(),
v.end(),
[](int a,int b)
{
    return a<b;
}
);
\end{lstlisting}
等价：
\begin{lstlisting}[language=C++]
sort(v.begin(),v.end());
\end{lstlisting}

\noindent\rule{\linewidth}{0.4pt}

\section{3. 降序}
\begin{lstlisting}[language=C++]
sort(
v.begin(),
v.end(),
[](int a,int b)
{
    return a>b;
}
);
\end{lstlisting}

\noindent\rule{\linewidth}{0.4pt}

\section{4. pair排序}
按照second：
\begin{lstlisting}[language=C++]
sort(
v.begin(),
v.end(),
[](pair<int,int>a,pair<int,int>b)
{
    return a.second<b.second;
}
);
\end{lstlisting}

\noindent\rule{\linewidth}{0.4pt}

\section{5. 自定义规则}
例如：\\
奇数在前，偶数在后：
\begin{lstlisting}[language=C++]
sort(
v.begin(),
v.end(),
[](int a,int b)
{
    return a%2>b%2;
}
);
\end{lstlisting}

\noindent\rule{\linewidth}{0.4pt}

\section{七、STL常见坑}

\noindent\rule{\linewidth}{0.4pt}

\section{1. vector删除元素}
错误：
\begin{lstlisting}[language=C++]
for(auto x:v)
{
    if(x==3)
        v.erase(...);
}
\end{lstlisting}
原因：\\
删除导致迭代器失效。

\noindent\rule{\linewidth}{0.4pt}

正确：
\begin{lstlisting}[language=C++]
for(auto it=v.begin();
    it!=v.end();)
{
    if(*it==3)
        it=v.erase(it);

    else
        it++;
}
\end{lstlisting}

\noindent\rule{\linewidth}{0.4pt}

\section{2. vector访问越界}
错误：
\begin{lstlisting}[language=C++]
vector<int> v;

v[0]=10;
\end{lstlisting}
因为：\\
没有空间。

正确：
\begin{lstlisting}[language=C++]
v.push_back(10);
\end{lstlisting}
或者：
\begin{lstlisting}[language=C++]
vector<int> v(10);
\end{lstlisting}

\noindent\rule{\linewidth}{0.4pt}

\section{3. map中[]自动创建}
例如：
\begin{lstlisting}[language=C++]
map<string,int> mp;


cout<<mp["Tom"];
\end{lstlisting}
即使不存在：\\
也会创建：
\begin{lstlisting}
Tom -> 0
\end{lstlisting}
如果只是查询：\\
用：
\begin{lstlisting}[language=C++]
find()
\end{lstlisting}

\noindent\rule{\linewidth}{0.4pt}

\section{4. sort区间}
错误：
\begin{lstlisting}[language=C++]
sort(v.begin(),v.end()-1);
\end{lstlisting}
少排序最后一个。

正确：
\begin{lstlisting}[language=C++]
sort(v.begin(),v.end());
\end{lstlisting}

\noindent\rule{\linewidth}{0.4pt}

\section{5. unique必须配erase}
错误：
\begin{lstlisting}[language=C++]
unique(v.begin(),v.end());
\end{lstlisting}
不会删除。

正确：
\begin{lstlisting}[language=C++]
v.erase(
unique(v.begin(),v.end()),
v.end()
);
\end{lstlisting}

\noindent\rule{\linewidth}{0.4pt}

\section{6. lower\_bound必须有序}
错误：
\begin{lstlisting}[language=C++]
vector<int> v={5,1,3};

lower_bound(v.begin(),v.end(),3);
\end{lstlisting}
结果错误。

必须：
\begin{lstlisting}[language=C++]
sort(v.begin(),v.end());
\end{lstlisting}

\noindent\rule{\linewidth}{0.4pt}

\section{7. priority\_queue不能遍历}
错误：
\begin{lstlisting}[language=C++]
for(auto x:pq)
\end{lstlisting}
priority\_queue没有迭代器。

只能：
\begin{lstlisting}[language=C++]
while(!pq.empty())
{
    cout<<pq.top();

    pq.pop();
}
\end{lstlisting}

\noindent\rule{\linewidth}{0.4pt}

\section{第九章 C++ STL函数速查表（机考打印版）}

这一章按照\textbf{考试查阅习惯}整理。

格式：
\begin{quote}
\textbf{名称 → 完整写法 → 参数怎么填 → 作用 → 示例}
\end{quote}

\noindent\rule{\linewidth}{0.4pt}

\section{一、vector速查表（重点）}

头文件：
\begin{lstlisting}[language=C++]
#include <vector>
\end{lstlisting}
定义：
\begin{lstlisting}[language=C++]
vector<int> v;
\end{lstlisting}

\noindent\rule{\linewidth}{0.4pt}

\section{1. push\_back()}

\subsection{完整写法}
\begin{lstlisting}[language=C++]
v.push_back(x);
\end{lstlisting}

\subsection{参数}
一个元素：
\begin{lstlisting}[language=C++]
x
\end{lstlisting}
例如：
\begin{lstlisting}[language=C++]
v.push_back(10);
\end{lstlisting}

\subsection{作用}
尾部添加元素。

\noindent\rule{\linewidth}{0.4pt}

\section{2. pop\_back()}

\subsection{完整写法}
\begin{lstlisting}[language=C++]
v.pop_back();
\end{lstlisting}

\subsection{参数}
无。

\subsection{作用}
删除最后一个元素。

\noindent\rule{\linewidth}{0.4pt}

\section{3. size()}

\subsection{完整写法}
\begin{lstlisting}[language=C++]
v.size();
\end{lstlisting}

\subsection{参数}
无。

\subsection{返回}
元素个数。

例如：
\begin{lstlisting}[language=C++]
cout<<v.size();
\end{lstlisting}

\noindent\rule{\linewidth}{0.4pt}

\section{4. empty()}

\subsection{完整写法}
\begin{lstlisting}[language=C++]
v.empty();
\end{lstlisting}

\subsection{参数}
无。

\subsection{返回}
\begin{lstlisting}[language=C++]
true / false
\end{lstlisting}
作用：\\
判断是否为空。

\noindent\rule{\linewidth}{0.4pt}

\section{5. clear()}

\subsection{完整写法}
\begin{lstlisting}[language=C++]
v.clear();
\end{lstlisting}
参数：\\
无。

作用：\\
清空vector。

\noindent\rule{\linewidth}{0.4pt}

\section{6. front()}

\subsection{完整写法}
\begin{lstlisting}[language=C++]
v.front();
\end{lstlisting}
参数：\\
无。

作用：\\
访问第一个元素。

\noindent\rule{\linewidth}{0.4pt}

\section{7. back()}

\subsection{完整写法}
\begin{lstlisting}[language=C++]
v.back();
\end{lstlisting}
参数：\\
无。

作用：\\
访问最后元素。

\noindent\rule{\linewidth}{0.4pt}

\section{8. begin()}

\subsection{完整写法}
\begin{lstlisting}[language=C++]
v.begin();
\end{lstlisting}
参数：\\
无。

作用：\\
返回首迭代器。

\noindent\rule{\linewidth}{0.4pt}

\section{9. end()}

\subsection{完整写法}
\begin{lstlisting}[language=C++]
v.end();
\end{lstlisting}
参数：\\
无。

作用：\\
返回尾后迭代器。

\noindent\rule{\linewidth}{0.4pt}

\section{10. insert()}

\subsection{完整写法}
\begin{lstlisting}[language=C++]
v.insert(pos,value);
\end{lstlisting}
参数：\\
两个：\\
① 插入位置\\
② 插入值

例如：
\begin{lstlisting}[language=C++]
v.insert(v.begin()+2,100);
\end{lstlisting}
作用：\\
指定位置插入。

\noindent\rule{\linewidth}{0.4pt}

多个元素：
\begin{lstlisting}[language=C++]
v.insert(pos,n,value);
\end{lstlisting}
例如：
\begin{lstlisting}[language=C++]
v.insert(v.begin(),3,5);
\end{lstlisting}
结果：
\begin{lstlisting}
5 5 5 ...
\end{lstlisting}
作用：\\
插入n个相同元素。

\noindent\rule{\linewidth}{0.4pt}

\section{11. erase()}

\subsection{完整写法1}
删除一个：
\begin{lstlisting}[language=C++]
v.erase(pos);
\end{lstlisting}
参数：\\
迭代器。

例如：
\begin{lstlisting}[language=C++]
v.erase(v.begin()+2);
\end{lstlisting}
作用：\\
删除指定位置。

\noindent\rule{\linewidth}{0.4pt}

\subsection{完整写法2}
删除区间：
\begin{lstlisting}[language=C++]
v.erase(first,last);
\end{lstlisting}
例如：
\begin{lstlisting}[language=C++]
v.erase(v.begin()+1,v.begin()+3);
\end{lstlisting}
删除：\\
下标1到2。

\noindent\rule{\linewidth}{0.4pt}

\section{12. resize()}

\subsection{完整写法}
\begin{lstlisting}[language=C++]
v.resize(n);
\end{lstlisting}
参数：\\
新大小。

例如：
\begin{lstlisting}[language=C++]
v.resize(10);
\end{lstlisting}
作用：\\
改变vector长度。

\noindent\rule{\linewidth}{0.4pt}

填充值：
\begin{lstlisting}[language=C++]
v.resize(10,0);
\end{lstlisting}
不足：\\
补0。

\noindent\rule{\linewidth}{0.4pt}

\section{13. capacity()}

\subsection{完整写法}
\begin{lstlisting}[language=C++]
v.capacity();
\end{lstlisting}
参数：\\
无。

作用：\\
返回当前容量。

注意：\\
不是元素数量。

区别：
\begin{lstlisting}[language=C++]
size()
\end{lstlisting}
实际元素数量。
\begin{lstlisting}[language=C++]
capacity()
\end{lstlisting}
申请空间大小。

\noindent\rule{\linewidth}{0.4pt}

\section{14. reserve()}

\subsection{完整写法}
\begin{lstlisting}[language=C++]
v.reserve(n);
\end{lstlisting}
参数：\\
预留容量。

例如：
\begin{lstlisting}[language=C++]
v.reserve(100);
\end{lstlisting}
作用：\\
提前申请空间。\\
减少扩容次数。

\noindent\rule{\linewidth}{0.4pt}

\section{15. shrink\_to\_fit()}

\subsection{完整写法}
\begin{lstlisting}[language=C++]
v.shrink_to_fit();
\end{lstlisting}
参数：\\
无。

作用：\\
尝试释放多余容量。

\noindent\rule{\linewidth}{0.4pt}

\section{16. at()}

\subsection{完整写法}
\begin{lstlisting}[language=C++]
v.at(index);
\end{lstlisting}
参数：\\
下标。

例如：
\begin{lstlisting}[language=C++]
v.at(2);
\end{lstlisting}
作用：\\
访问元素。

区别：
\begin{lstlisting}[language=C++]
v[2]
\end{lstlisting}
越界不检查。
\begin{lstlisting}[language=C++]
v.at(2)
\end{lstlisting}
越界抛异常。

\noindent\rule{\linewidth}{0.4pt}

\section{二、string速查表（重点）}

头文件：
\begin{lstlisting}[language=C++]
#include <string>
\end{lstlisting}
定义：
\begin{lstlisting}[language=C++]
string s;
\end{lstlisting}

\noindent\rule{\linewidth}{0.4pt}

\section{1. size()}
\begin{lstlisting}[language=C++]
s.size();
\end{lstlisting}
作用：\\
字符串长度。

\noindent\rule{\linewidth}{0.4pt}

\section{2. length()}
\begin{lstlisting}[language=C++]
s.length();
\end{lstlisting}
作用：\\
字符串长度。\\
和size完全一样。

\noindent\rule{\linewidth}{0.4pt}

\section{3. empty()}
\begin{lstlisting}[language=C++]
s.empty();
\end{lstlisting}
作用：\\
判断字符串为空。

\noindent\rule{\linewidth}{0.4pt}

\section{4. clear()}
\begin{lstlisting}[language=C++]
s.clear();
\end{lstlisting}
作用：\\
清空字符串。

\noindent\rule{\linewidth}{0.4pt}

\section{5. push\_back()}
\begin{lstlisting}[language=C++]
s.push_back(c);
\end{lstlisting}
参数：\\
字符。

例如：
\begin{lstlisting}[language=C++]
s.push_back('a');
\end{lstlisting}
作用：\\
尾部添加字符。

\noindent\rule{\linewidth}{0.4pt}

\section{6. pop\_back()}
\begin{lstlisting}[language=C++]
s.pop_back();
\end{lstlisting}
作用：\\
删除最后字符。

\noindent\rule{\linewidth}{0.4pt}

\section{7. append()}

\subsection{写法}
\begin{lstlisting}[language=C++]
s.append(str);
\end{lstlisting}
例如：
\begin{lstlisting}[language=C++]
s="hello";

s.append(" world");
\end{lstlisting}
结果：
\begin{lstlisting}
hello world
\end{lstlisting}
作用：\\
追加字符串。

\noindent\rule{\linewidth}{0.4pt}

\section{8. substr()}

\subsection{写法}
\begin{lstlisting}[language=C++]
s.substr(pos,len);
\end{lstlisting}
两个参数：\\
① 起始位置\\
② 长度

例如：
\begin{lstlisting}[language=C++]
string s="abcdef";


cout<<s.substr(1,3);
\end{lstlisting}
结果：
\begin{lstlisting}
bcd
\end{lstlisting}
作用：\\
截取子串。

\noindent\rule{\linewidth}{0.4pt}

\section{9. find()}

\subsection{写法}
\begin{lstlisting}[language=C++]
s.find(str);
\end{lstlisting}
参数：\\
要查找内容。

例如：
\begin{lstlisting}[language=C++]
s.find("abc");
\end{lstlisting}
返回：\\
第一次出现位置。

不存在：
\begin{lstlisting}[language=C++]
string::npos
\end{lstlisting}

\noindent\rule{\linewidth}{0.4pt}

\section{10. replace()}

\subsection{写法}
\begin{lstlisting}[language=C++]
s.replace(pos,len,str);
\end{lstlisting}
参数：\\
三个：\\
① 开始位置\\
② 替换长度\\
③ 新字符串

例如：
\begin{lstlisting}[language=C++]
s.replace(1,2,"xx");
\end{lstlisting}
作用：\\
替换字符串。

\noindent\rule{\linewidth}{0.4pt}

\section{11. erase()}

\subsection{删除位置}
\begin{lstlisting}[language=C++]
s.erase(pos,len);
\end{lstlisting}
例如：
\begin{lstlisting}[language=C++]
s.erase(1,3);
\end{lstlisting}
作用：\\
删除指定字符。

\noindent\rule{\linewidth}{0.4pt}

\section{12. insert()}

\subsection{写法}
\begin{lstlisting}[language=C++]
s.insert(pos,str);
\end{lstlisting}
例如：
\begin{lstlisting}[language=C++]
s.insert(2,"abc");
\end{lstlisting}
作用：\\
插入字符串。

\noindent\rule{\linewidth}{0.4pt}

\section{13. compare()}

\subsection{写法}
\begin{lstlisting}[language=C++]
s.compare(str);
\end{lstlisting}
作用：\\
比较字符串。

返回：
\begin{lstlisting}[language=C++]
0
\end{lstlisting}
相等。

\noindent\rule{\linewidth}{0.4pt}

\section{14. getline()}

\subsection{写法}
\begin{lstlisting}[language=C++]
getline(cin,s);
\end{lstlisting}
作用：\\
读取一整行。

区别：
\begin{lstlisting}[language=C++]
cin>>
\end{lstlisting}
遇空格停止。

getline：\\
包含空格。

\noindent\rule{\linewidth}{0.4pt}

\section{三、algorithm速查表}

头文件：
\begin{lstlisting}[language=C++]
#include <algorithm>
\end{lstlisting}

\noindent\rule{\linewidth}{0.4pt}

\section{1. sort()}

\subsection{写法}
\begin{lstlisting}[language=C++]
sort(begin,end);
\end{lstlisting}
参数：\\
两个迭代器。

例如：
\begin{lstlisting}[language=C++]
sort(v.begin(),v.end());
\end{lstlisting}
作用：\\
排序。

\noindent\rule{\linewidth}{0.4pt}

降序：
\begin{lstlisting}[language=C++]
sort(
v.begin(),
v.end(),
greater<int>()
);
\end{lstlisting}

\noindent\rule{\linewidth}{0.4pt}

\section{2. reverse()}

\subsection{写法}
\begin{lstlisting}[language=C++]
reverse(begin,end);
\end{lstlisting}
作用：\\
反转。

\noindent\rule{\linewidth}{0.4pt}

\section{3. max()}

\subsection{写法}
\begin{lstlisting}[language=C++]
max(a,b);
\end{lstlisting}
作用：\\
最大值。

\noindent\rule{\linewidth}{0.4pt}

\section{4. min()}
\begin{lstlisting}[language=C++]
min(a,b);
\end{lstlisting}
作用：\\
最小值。

\noindent\rule{\linewidth}{0.4pt}

\section{5. max\_element()}

\subsection{写法}
\begin{lstlisting}[language=C++]
max_element(begin,end);
\end{lstlisting}
返回：\\
迭代器。

使用：
\begin{lstlisting}[language=C++]
*max_element(v.begin(),v.end());
\end{lstlisting}
作用：\\
最大元素。

\noindent\rule{\linewidth}{0.4pt}

\section{6. min\_element()}
\begin{lstlisting}[language=C++]
*min_element(
v.begin(),
v.end()
);
\end{lstlisting}
作用：\\
最小元素。

\noindent\rule{\linewidth}{0.4pt}

\section{7. find()}

\subsection{写法}
\begin{lstlisting}[language=C++]
find(begin,end,x);
\end{lstlisting}
参数：\\
三个：\\
①开始\\
②结束\\
③查找值

作用：\\
寻找元素。

\noindent\rule{\linewidth}{0.4pt}

\section{8. count()}
\begin{lstlisting}[language=C++]
count(begin,end,x);
\end{lstlisting}
作用：\\
统计数量。

\noindent\rule{\linewidth}{0.4pt}

\section{9. fill()}
\begin{lstlisting}[language=C++]
fill(begin,end,x);
\end{lstlisting}
作用：\\
全部填充。

\noindent\rule{\linewidth}{0.4pt}

\section{10. replace()}
\begin{lstlisting}[language=C++]
replace(
begin,
end,
old,
new
);
\end{lstlisting}
作用：\\
替换元素。

\noindent\rule{\linewidth}{0.4pt}

\section{11. unique()}
\begin{lstlisting}[language=C++]
unique(begin,end);
\end{lstlisting}
作用：\\
删除连续重复。

常用：
\begin{lstlisting}[language=C++]
v.erase(
unique(v.begin(),v.end()),
v.end()
);
\end{lstlisting}

\noindent\rule{\linewidth}{0.4pt}

\section{12. lower\_bound()}
\begin{lstlisting}[language=C++]
lower_bound(
begin,
end,
x
);
\end{lstlisting}
要求：\\
有序。

作用：\\
第一个：
\begin{lstlisting}
>=x
\end{lstlisting}

\noindent\rule{\linewidth}{0.4pt}

\section{13. upper\_bound()}
\begin{lstlisting}[language=C++]
upper_bound(
begin,
end,
x
);
\end{lstlisting}
作用：\\
第一个：
\begin{lstlisting}
>x
\end{lstlisting}

\noindent\rule{\linewidth}{0.4pt}

\section{14. binary\_search()}
\begin{lstlisting}[language=C++]
binary_search(
begin,
end,
x
);
\end{lstlisting}
作用：\\
二分查找是否存在。

返回：\\
true/false。

\noindent\rule{\linewidth}{0.4pt}

\section{15. next\_permutation()}
\begin{lstlisting}[language=C++]
next_permutation(
begin,
end
);
\end{lstlisting}
作用：\\
下一个排列。

\noindent\rule{\linewidth}{0.4pt}

\section{四、numeric速查表}

头文件：
\begin{lstlisting}[language=C++]
#include <numeric>
\end{lstlisting}

\noindent\rule{\linewidth}{0.4pt}

\section{1. accumulate()}

\subsection{写法}
\begin{lstlisting}[language=C++]
accumulate(
begin,
end,
init
);
\end{lstlisting}
参数：\\
三个：\\
①开始\\
②结束\\
③初始值

例如：
\begin{lstlisting}[language=C++]
accumulate(
v.begin(),
v.end(),
0
);
\end{lstlisting}
作用：\\
求和。

\noindent\rule{\linewidth}{0.4pt}

\section{2. iota()}
\begin{lstlisting}[language=C++]
iota(
begin,
end,
start
);
\end{lstlisting}
作用：\\
连续赋值。

例如：
\begin{lstlisting}[language=C++]
iota(v.begin(),v.end(),1);
\end{lstlisting}
结果：
\begin{lstlisting}
1 2 3...
\end{lstlisting}

\noindent\rule{\linewidth}{0.4pt}

\section{3. partial\_sum()}
\begin{lstlisting}[language=C++]
partial_sum(
begin,
end,
out
);
\end{lstlisting}
作用：\\
前缀和。

\noindent\rule{\linewidth}{0.4pt}

\section{五、map速查表}

定义：
\begin{lstlisting}[language=C++]
map<int,int> mp;
\end{lstlisting}

\noindent\rule{\linewidth}{0.4pt}

\section{1. []}
\begin{lstlisting}[language=C++]
mp[key];
\end{lstlisting}
作用：\\
访问value。\\
不存在自动创建。

\noindent\rule{\linewidth}{0.4pt}

\section{2. insert()}
\begin{lstlisting}[language=C++]
mp.insert({key,value});
\end{lstlisting}
作用：\\
插入。

\noindent\rule{\linewidth}{0.4pt}

\section{3. erase()}
\begin{lstlisting}[language=C++]
mp.erase(key);
\end{lstlisting}
作用：\\
删除。

\noindent\rule{\linewidth}{0.4pt}

\section{4. find()}
\begin{lstlisting}[language=C++]
mp.find(key);
\end{lstlisting}
作用：\\
查找key。

\noindent\rule{\linewidth}{0.4pt}

\section{5. count()}
\begin{lstlisting}[language=C++]
mp.count(key);
\end{lstlisting}
作用：\\
判断key是否存在。

\noindent\rule{\linewidth}{0.4pt}

\section{6. 遍历}
\begin{lstlisting}[language=C++]
for(auto p:mp)
{
    cout<<p.first;
    cout<<p.second;
}
\end{lstlisting}

\noindent\rule{\linewidth}{0.4pt}

\section{六、stack速查表}

定义：
\begin{lstlisting}[language=C++]
stack<int> st;
\end{lstlisting}

\noindent\rule{\linewidth}{0.4pt}

push：
\begin{lstlisting}[language=C++]
st.push(x);
\end{lstlisting}
入栈。

\noindent\rule{\linewidth}{0.4pt}

pop：
\begin{lstlisting}[language=C++]
st.pop();
\end{lstlisting}
删除栈顶。

\noindent\rule{\linewidth}{0.4pt}

top：
\begin{lstlisting}[language=C++]
st.top();
\end{lstlisting}
访问栈顶。

\noindent\rule{\linewidth}{0.4pt}

empty：
\begin{lstlisting}[language=C++]
st.empty();
\end{lstlisting}
判断空。

\noindent\rule{\linewidth}{0.4pt}

\section{七、queue速查表}

定义：
\begin{lstlisting}[language=C++]
queue<int> q;
\end{lstlisting}

\noindent\rule{\linewidth}{0.4pt}

push：
\begin{lstlisting}[language=C++]
q.push(x);
\end{lstlisting}
入队。

\noindent\rule{\linewidth}{0.4pt}

pop：
\begin{lstlisting}[language=C++]
q.pop();
\end{lstlisting}
出队。

\noindent\rule{\linewidth}{0.4pt}

front：
\begin{lstlisting}[language=C++]
q.front();
\end{lstlisting}
队头。

\noindent\rule{\linewidth}{0.4pt}

back：
\begin{lstlisting}[language=C++]
q.back();
\end{lstlisting}
队尾。

\noindent\rule{\linewidth}{0.4pt}

\section{八、priority\_queue速查表}

定义：
\begin{lstlisting}[language=C++]
priority_queue<int> pq;
\end{lstlisting}

\noindent\rule{\linewidth}{0.4pt}

push：
\begin{lstlisting}[language=C++]
pq.push(x);
\end{lstlisting}

\noindent\rule{\linewidth}{0.4pt}

top：
\begin{lstlisting}[language=C++]
pq.top();
\end{lstlisting}
作用：\\
最大值。

\noindent\rule{\linewidth}{0.4pt}

pop：
\begin{lstlisting}[language=C++]
pq.pop();
\end{lstlisting}

\noindent\rule{\linewidth}{0.4pt}

小根堆：
\begin{lstlisting}[language=C++]
priority_queue<
int,
vector<int>,
greater<int>
> pq;
\end{lstlisting}

\noindent\rule{\linewidth}{0.4pt}

\section{九、set速查表}

定义：
\begin{lstlisting}[language=C++]
set<int> s;
\end{lstlisting}

\noindent\rule{\linewidth}{0.4pt}

插入：
\begin{lstlisting}[language=C++]
s.insert(x);
\end{lstlisting}
删除：
\begin{lstlisting}[language=C++]
s.erase(x);
\end{lstlisting}
查找：
\begin{lstlisting}[language=C++]
s.find(x);
\end{lstlisting}
数量：
\begin{lstlisting}[language=C++]
s.size();
\end{lstlisting}
判断存在：
\begin{lstlisting}[language=C++]
s.count(x);
\end{lstlisting}

\noindent\rule{\linewidth}{0.4pt}

\section{十、考试最常忘参数总结}

\begin{table}[h]
\centering
\begin{tabular}{|l|l|}
\hline
函数 & 参数 \\ \hline
sort & 开始，结束，比较函数 \\ \hline
reverse & 开始，结束 \\ \hline
find & 开始，结束，值 \\ \hline
count & 开始，结束，值 \\ \hline
fill & 开始，结束，值 \\ \hline
accumulate & 开始，结束，初始值 \\ \hline
lower\_bound & 开始，结束，值 \\ \hline
upper\_bound & 开始，结束，值 \\ \hline
substr & 位置，长度 \\ \hline
erase(vector) & 位置 或 区间 \\ \hline
insert(vector) & 位置，值 \\ \hline
replace(string) & 位置，长度，新字符串 \\ \hline
resize & 新大小 \\ \hline
reserve & 容量 \\ \hline
\end{tabular}
\end{table}

\noindent\rule{\linewidth}{0.4pt}

\section{第十章 C++ 常用头文件速查（cmath + 其他高频库）}

这一部分整理\textbf{上机考试除了 STL 之外最常用的函数库}。

重点：
\begin{itemize}
    \item \texttt{<cmath>} 数学函数（重点）
    \item \texttt{<iomanip>} 格式化输出
    \item \texttt{<cstdlib>} 常用工具函数
    \item \texttt{<cstring>} C字符串处理
    \item \texttt{<ctime>} 时间函数
    \item \texttt{<numeric>} 数学运算
    \item \texttt{<sstream>} 字符串流
    \item \texttt{<fstream>} 文件操作（了解）
    \item \texttt{<random>} 随机数
    \item \texttt{<limits>} 极值
\end{itemize}

\noindent\rule{\linewidth}{0.4pt}

\section{一、cmath（数学函数库）\textbf{【必背】}}

头文件：
\begin{lstlisting}[language=C++]
#include <cmath>
\end{lstlisting}
作用：\\
提供数学计算函数。

\noindent\rule{\linewidth}{0.4pt}

\section{1. abs()}

\subsection{名称}
\begin{lstlisting}[language=C++]
abs()
\end{lstlisting}

\noindent\rule{\linewidth}{0.4pt}

\subsection{用法}
\begin{lstlisting}[language=C++]
abs(x);
\end{lstlisting}
参数：\\
一个数字。

例如：
\begin{lstlisting}[language=C++]
int x=-10;

cout<<abs(x);
\end{lstlisting}
输出：
\begin{lstlisting}
10
\end{lstlisting}

\noindent\rule{\linewidth}{0.4pt}

\subsection{作用}
求绝对值。

注意：\\
整数：
\begin{lstlisting}[language=C++]
abs()
\end{lstlisting}
浮点：\\
推荐：
\begin{lstlisting}[language=C++]
fabs()
\end{lstlisting}

\noindent\rule{\linewidth}{0.4pt}

\section{2. fabs()}

\subsection{名称}
\begin{lstlisting}[language=C++]
fabs()
\end{lstlisting}

\noindent\rule{\linewidth}{0.4pt}

\subsection{用法}
\begin{lstlisting}[language=C++]
fabs(x);
\end{lstlisting}
例如：
\begin{lstlisting}[language=C++]
double x=-3.14;

cout<<fabs(x);
\end{lstlisting}
输出：
\begin{lstlisting}
3.14
\end{lstlisting}

\noindent\rule{\linewidth}{0.4pt}

\subsection{作用}
求浮点数绝对值。

\noindent\rule{\linewidth}{0.4pt}

\section{3. sqrt()}

\subsection{名称}
\begin{lstlisting}[language=C++]
sqrt()
\end{lstlisting}

\noindent\rule{\linewidth}{0.4pt}

\subsection{用法}
\begin{lstlisting}[language=C++]
sqrt(x);
\end{lstlisting}
例如：
\begin{lstlisting}[language=C++]
cout<<sqrt(9);
\end{lstlisting}
输出：
\begin{lstlisting}
3
\end{lstlisting}
浮点：
\begin{lstlisting}[language=C++]
cout<<sqrt(2.0);
\end{lstlisting}
输出：
\begin{lstlisting}
1.414...
\end{lstlisting}

\noindent\rule{\linewidth}{0.4pt}

\subsection{作用}
平方根。

\noindent\rule{\linewidth}{0.4pt}

\section{4. pow()}

\subsection{名称}
\begin{lstlisting}[language=C++]
pow()
\end{lstlisting}

\noindent\rule{\linewidth}{0.4pt}

\subsection{用法}
两个参数：
\begin{lstlisting}[language=C++]
pow(a,b);
\end{lstlisting}
表示：
\begin{lstlisting}
a的b次方
\end{lstlisting}
例如：
\begin{lstlisting}[language=C++]
cout<<pow(2,3);
\end{lstlisting}
输出：
\begin{lstlisting}
8
\end{lstlisting}

\noindent\rule{\linewidth}{0.4pt}

\subsection{作用}
幂运算。

常见：\\
平方：
\begin{lstlisting}[language=C++]
pow(x,2)
\end{lstlisting}
立方：
\begin{lstlisting}[language=C++]
pow(x,3)
\end{lstlisting}

\noindent\rule{\linewidth}{0.4pt}

\section{5. exp()}

\subsection{名称}
\begin{lstlisting}[language=C++]
exp()
\end{lstlisting}

\noindent\rule{\linewidth}{0.4pt}

\subsection{用法}
\begin{lstlisting}[language=C++]
exp(x);
\end{lstlisting}
表示：\\
$e^x$

例如：
\begin{lstlisting}[language=C++]
cout<<exp(1);
\end{lstlisting}
约：
\begin{lstlisting}
2.71828
\end{lstlisting}

\noindent\rule{\linewidth}{0.4pt}

\subsection{作用}
计算自然指数。

\noindent\rule{\linewidth}{0.4pt}

\section{6. log()}

\subsection{名称}
\begin{lstlisting}[language=C++]
log()
\end{lstlisting}

\noindent\rule{\linewidth}{0.4pt}

\subsection{用法}
\begin{lstlisting}[language=C++]
log(x);
\end{lstlisting}
表示：\\
$ln(x)$

例如：
\begin{lstlisting}[language=C++]
cout<<log(2);
\end{lstlisting}

\noindent\rule{\linewidth}{0.4pt}

\subsection{作用}
自然对数。

\noindent\rule{\linewidth}{0.4pt}

\section{7. log10()}

\subsection{名称}
\begin{lstlisting}[language=C++]
log10()
\end{lstlisting}

\noindent\rule{\linewidth}{0.4pt}

\subsection{用法}
\begin{lstlisting}[language=C++]
log10(x);
\end{lstlisting}
例如：
\begin{lstlisting}[language=C++]
cout<<log10(100);
\end{lstlisting}
输出：
\begin{lstlisting}
2
\end{lstlisting}

\noindent\rule{\linewidth}{0.4pt}

\subsection{作用}
十进制对数。

\noindent\rule{\linewidth}{0.4pt}

\section{8. sin()}

\subsection{名称}
\begin{lstlisting}[language=C++]
sin()
\end{lstlisting}

\noindent\rule{\linewidth}{0.4pt}

\subsection{用法}
\begin{lstlisting}[language=C++]
sin(x);
\end{lstlisting}
注意：\\
参数单位：\\
\textbf{弧度}\\
不是角度。

例如：
\begin{lstlisting}[language=C++]
sin(M_PI/2);
\end{lstlisting}
结果：
\begin{lstlisting}
1
\end{lstlisting}

\noindent\rule{\linewidth}{0.4pt}

\subsection{作用}
正弦。

\noindent\rule{\linewidth}{0.4pt}

\section{9. cos()}

\subsection{用法}
\begin{lstlisting}[language=C++]
cos(x);
\end{lstlisting}
作用：\\
余弦。

\noindent\rule{\linewidth}{0.4pt}

\section{10. tan()}

\subsection{用法}
\begin{lstlisting}[language=C++]
tan(x);
\end{lstlisting}
作用：\\
正切。

\noindent\rule{\linewidth}{0.4pt}

\section{11. asin()}

\subsection{名称}
\begin{lstlisting}[language=C++]
asin()
\end{lstlisting}
作用：\\
反正弦。

返回：\\
弧度。

\noindent\rule{\linewidth}{0.4pt}

\section{12. acos()}
\begin{lstlisting}[language=C++]
acos(x);
\end{lstlisting}
作用：\\
反余弦。

\noindent\rule{\linewidth}{0.4pt}

\section{13. atan()}
\begin{lstlisting}[language=C++]
atan(x);
\end{lstlisting}
作用：\\
反正切。

\noindent\rule{\linewidth}{0.4pt}

\section{14. atan2()}

\textbf{【重点】}

\subsection{用法}
\begin{lstlisting}[language=C++]
atan2(y,x);
\end{lstlisting}
返回：\\
点：\\
$(x,y)$\\
的角度。

例如：
\begin{lstlisting}[language=C++]
double angle=atan2(y,x);
\end{lstlisting}

\noindent\rule{\linewidth}{0.4pt}

\subsection{作用}
计算方向角。

常用于：
\begin{itemize}
    \item 几何
    \item 向量
\end{itemize}

\noindent\rule{\linewidth}{0.4pt}

\section{15. ceil()}

\subsection{名称}
\begin{lstlisting}[language=C++]
ceil()
\end{lstlisting}

\noindent\rule{\linewidth}{0.4pt}

\subsection{用法}
\begin{lstlisting}[language=C++]
ceil(x);
\end{lstlisting}
例如：
\begin{lstlisting}[language=C++]
ceil(3.2);
\end{lstlisting}
结果：
\begin{lstlisting}
4
\end{lstlisting}

\noindent\rule{\linewidth}{0.4pt}

\subsection{作用}
向上取整。

\noindent\rule{\linewidth}{0.4pt}

\section{16. floor()}

\subsection{名称}
\begin{lstlisting}[language=C++]
floor()
\end{lstlisting}

\noindent\rule{\linewidth}{0.4pt}

\subsection{用法}
\begin{lstlisting}[language=C++]
floor(x);
\end{lstlisting}
例如：
\begin{lstlisting}[language=C++]
floor(3.8);
\end{lstlisting}
结果：
\begin{lstlisting}
3
\end{lstlisting}

\noindent\rule{\linewidth}{0.4pt}

\subsection{作用}
向下取整。

\noindent\rule{\linewidth}{0.4pt}

\section{17. round()}

\subsection{名称}
\begin{lstlisting}[language=C++]
round()
\end{lstlisting}

\noindent\rule{\linewidth}{0.4pt}

\subsection{用法}
\begin{lstlisting}[language=C++]
round(x);
\end{lstlisting}
例如：
\begin{lstlisting}[language=C++]
round(3.5);
\end{lstlisting}
结果：
\begin{lstlisting}
4
\end{lstlisting}

\noindent\rule{\linewidth}{0.4pt}

\subsection{作用}
四舍五入。

\noindent\rule{\linewidth}{0.4pt}

\section{18. trunc()}

\subsection{名称}
\begin{lstlisting}[language=C++]
trunc()
\end{lstlisting}

\noindent\rule{\linewidth}{0.4pt}

\subsection{用法}
\begin{lstlisting}[language=C++]
trunc(x);
\end{lstlisting}
例如：
\begin{lstlisting}[language=C++]
trunc(3.8);
\end{lstlisting}
结果：
\begin{lstlisting}
3
\end{lstlisting}

\noindent\rule{\linewidth}{0.4pt}

\subsection{作用}
截断小数。

\noindent\rule{\linewidth}{0.4pt}

\section{19. fmod()}

\subsection{名称}
\begin{lstlisting}[language=C++]
fmod()
\end{lstlisting}

\noindent\rule{\linewidth}{0.4pt}

\subsection{用法}
\begin{lstlisting}[language=C++]
fmod(a,b);
\end{lstlisting}
例如：
\begin{lstlisting}[language=C++]
fmod(5.5,2);
\end{lstlisting}
结果：
\begin{lstlisting}
1.5
\end{lstlisting}

\noindent\rule{\linewidth}{0.4pt}

\subsection{作用}
浮点取余。

类似：\\
整数：
\begin{lstlisting}[language=C++]
a%b
\end{lstlisting}
浮点：
\begin{lstlisting}[language=C++]
fmod(a,b)
\end{lstlisting}

\noindent\rule{\linewidth}{0.4pt}

\section{20. hypot()}

\subsection{名称}
\begin{lstlisting}[language=C++]
hypot()
\end{lstlisting}

\noindent\rule{\linewidth}{0.4pt}

\subsection{用法}
\begin{lstlisting}[language=C++]
hypot(a,b);
\end{lstlisting}
计算：\\
$\sqrt{a^2+b^2}$

例如：
\begin{lstlisting}[language=C++]
hypot(3,4);
\end{lstlisting}
结果：
\begin{lstlisting}
5
\end{lstlisting}

\noindent\rule{\linewidth}{0.4pt}

\subsection{作用}
求直角三角形斜边。

\noindent\rule{\linewidth}{0.4pt}

\section{cmath常量}

\subsection{M\_PI}
圆周率：
\begin{lstlisting}[language=C++]
M_PI
\end{lstlisting}
例如：
\begin{lstlisting}[language=C++]
double r=2;

double area=M_PI*r*r;
\end{lstlisting}

\noindent\rule{\linewidth}{0.4pt}

\section{二、iomanip（格式化输出）\textbf{【重要】}}

头文件：
\begin{lstlisting}[language=C++]
#include <iomanip>
\end{lstlisting}

\noindent\rule{\linewidth}{0.4pt}

\section{1. setprecision()}

\subsection{名称}
\begin{lstlisting}[language=C++]
setprecision()
\end{lstlisting}

\noindent\rule{\linewidth}{0.4pt}

\subsection{用法}
\begin{lstlisting}[language=C++]
cout<<setprecision(n);
\end{lstlisting}
例如：
\begin{lstlisting}[language=C++]
cout<<setprecision(3)<<3.14159;
\end{lstlisting}
输出：
\begin{lstlisting}
3.14
\end{lstlisting}

\noindent\rule{\linewidth}{0.4pt}

默认：\\
控制有效数字。

\noindent\rule{\linewidth}{0.4pt}

配合：
\begin{lstlisting}[language=C++]
fixed
\end{lstlisting}

\noindent\rule{\linewidth}{0.4pt}

\section{2. fixed}

\subsection{用法}
\begin{lstlisting}[language=C++]
cout<<fixed;
\end{lstlisting}
作用：\\
固定小数位。

\noindent\rule{\linewidth}{0.4pt}

\section{3. setprecision + fixed}

\textbf{【必背】}

\begin{lstlisting}[language=C++]
cout
<<fixed
<<setprecision(2)
<<3.14159;
\end{lstlisting}
输出：
\begin{lstlisting}
3.14
\end{lstlisting}
考试最常用。

\noindent\rule{\linewidth}{0.4pt}

\section{4. setw()}

\subsection{名称}
\begin{lstlisting}[language=C++]
setw()
\end{lstlisting}

\noindent\rule{\linewidth}{0.4pt}

\subsection{用法}
\begin{lstlisting}[language=C++]
setw(n)
\end{lstlisting}
例如：
\begin{lstlisting}[language=C++]
cout<<setw(5)<<123;
\end{lstlisting}
输出：
\begin{lstlisting}
  123
\end{lstlisting}

\noindent\rule{\linewidth}{0.4pt}

\subsection{作用}
设置输出宽度。

\noindent\rule{\linewidth}{0.4pt}

\section{5. left}
\begin{lstlisting}[language=C++]
cout<<left;
\end{lstlisting}
作用：\\
左对齐。

\noindent\rule{\linewidth}{0.4pt}

\section{6. right}
\begin{lstlisting}[language=C++]
cout<<right;
\end{lstlisting}
作用：\\
右对齐。

\noindent\rule{\linewidth}{0.4pt}

\section{7. setfill()}

\subsection{用法}
\begin{lstlisting}[language=C++]
setfill('*')
\end{lstlisting}
例如：
\begin{lstlisting}[language=C++]
cout
<<setfill('0')
<<setw(4)
<<12;
\end{lstlisting}
输出：
\begin{lstlisting}
0012
\end{lstlisting}

\noindent\rule{\linewidth}{0.4pt}

\section{三、cstring（C风格字符串）}

头文件：
\begin{lstlisting}[language=C++]
#include <cstring>
\end{lstlisting}
处理：
\begin{lstlisting}[language=C++]
char数组
\end{lstlisting}

\noindent\rule{\linewidth}{0.4pt}

\section{1. strlen()}

\subsection{名称}
\begin{lstlisting}[language=C++]
strlen()
\end{lstlisting}

\subsection{用法}
\begin{lstlisting}[language=C++]
strlen(str);
\end{lstlisting}
例如：
\begin{lstlisting}[language=C++]
char s[]="hello";


cout<<strlen(s);
\end{lstlisting}
输出：
\begin{lstlisting}
5
\end{lstlisting}

\noindent\rule{\linewidth}{0.4pt}

\subsection{作用}
计算字符串长度。

\noindent\rule{\linewidth}{0.4pt}

\section{2. strcpy()}

\subsection{用法}
\begin{lstlisting}[language=C++]
strcpy(dest,src);
\end{lstlisting}
作用：\\
复制字符串。

\noindent\rule{\linewidth}{0.4pt}

\section{3. strcat()}

\subsection{用法}
\begin{lstlisting}[language=C++]
strcat(a,b);
\end{lstlisting}
作用：\\
连接字符串。

\noindent\rule{\linewidth}{0.4pt}

\section{4. strcmp()}

\subsection{用法}
\begin{lstlisting}[language=C++]
strcmp(a,b);
\end{lstlisting}
作用：\\
比较字符串。

返回：
\begin{lstlisting}
0 相等
<0 小于
>0 大于
\end{lstlisting}

\noindent\rule{\linewidth}{0.4pt}

\section{5. memset()}

\textbf{【必背】}

\subsection{用法}
\begin{lstlisting}[language=C++]
memset(数组,值,sizeof(数组));
\end{lstlisting}
例如：
\begin{lstlisting}[language=C++]
int a[100];


memset(a,0,sizeof(a));
\end{lstlisting}
作用：\\
快速初始化数组。

\noindent\rule{\linewidth}{0.4pt}

常用：

全部变：

0：
\begin{lstlisting}[language=C++]
memset(a,0,sizeof(a));
\end{lstlisting}

-1：
\begin{lstlisting}[language=C++]
memset(a,-1,sizeof(a));
\end{lstlisting}

\noindent\rule{\linewidth}{0.4pt}

注意：

不要：
\begin{lstlisting}[language=C++]
memset(a,1,sizeof(a));
\end{lstlisting}
整数不会全部变1。

\noindent\rule{\linewidth}{0.4pt}

\section{四、cstdlib（常用工具）}

头文件：
\begin{lstlisting}[language=C++]
#include <cstdlib>
\end{lstlisting}

\noindent\rule{\linewidth}{0.4pt}

\section{1. rand()}

\subsection{用法}
\begin{lstlisting}[language=C++]
rand();
\end{lstlisting}
作用：\\
生成随机数。

\noindent\rule{\linewidth}{0.4pt}

\section{2. srand()}

\subsection{用法}
\begin{lstlisting}[language=C++]
srand(time(0));
\end{lstlisting}
作用：\\
设置随机种子。

通常：
\begin{lstlisting}[language=C++]
srand(time(NULL));
\end{lstlisting}

\noindent\rule{\linewidth}{0.4pt}

\section{3. atoi()}

\subsection{用法}
\begin{lstlisting}[language=C++]
atoi(str);
\end{lstlisting}
例如：
\begin{lstlisting}[language=C++]
char s[]="123";


int x=atoi(s);
\end{lstlisting}
作用：\\
字符串转整数。

\noindent\rule{\linewidth}{0.4pt}

\section{4. exit()}

\subsection{用法}
\begin{lstlisting}[language=C++]
exit(0);
\end{lstlisting}
作用：\\
结束程序。

\noindent\rule{\linewidth}{0.4pt}

\section{五、sstream（字符串流）}

头文件：
\begin{lstlisting}[language=C++]
#include <sstream>
\end{lstlisting}

\noindent\rule{\linewidth}{0.4pt}

\section{1. stringstream}

\subsection{定义}
\begin{lstlisting}[language=C++]
stringstream ss;
\end{lstlisting}

\noindent\rule{\linewidth}{0.4pt}

\section{字符串转数字}
\begin{lstlisting}[language=C++]
string s="123";


stringstream ss(s);


int x;


ss>>x;
\end{lstlisting}
结果：
\begin{lstlisting}
x=123
\end{lstlisting}

\noindent\rule{\linewidth}{0.4pt}

\section{数字转字符串}
\begin{lstlisting}[language=C++]
stringstream ss;


ss<<123;


string s;


ss>>s;
\end{lstlisting}

\noindent\rule{\linewidth}{0.4pt}

\subsection{作用}
字符串和数字互转。

\noindent\rule{\linewidth}{0.4pt}

\section{六、limits（极值）}

头文件：
\begin{lstlisting}[language=C++]
#include <limits>
\end{lstlisting}

\noindent\rule{\linewidth}{0.4pt}

\section{1. INT\_MAX}

\subsection{用法}
\begin{lstlisting}[language=C++]
INT_MAX
\end{lstlisting}
作用：\\
int最大值。

约：
\begin{lstlisting}
2147483647
\end{lstlisting}

\noindent\rule{\linewidth}{0.4pt}

\section{2. INT\_MIN}
\begin{lstlisting}[language=C++]
INT_MIN
\end{lstlisting}
作用：\\
int最小值。

\noindent\rule{\linewidth}{0.4pt}

\section{3. numeric\_limits}

\subsection{用法}
\begin{lstlisting}[language=C++]
numeric_limits<int>::max();
\end{lstlisting}
作用：\\
获取类型最大值。

例如：
\begin{lstlisting}[language=C++]
long long INF=
numeric_limits<long long>::max();
\end{lstlisting}

\noindent\rule{\linewidth}{0.4pt}

\section{七、常用头文件总表（考试版）}

\begin{table}[h]
\centering
\begin{tabular}{|l|l|l|}
\hline
头文件 & 作用 & 常用 \\ \hline
cmath & 数学计算 & sqrt,pow,sin \\ \hline
iomanip & 格式输出 & fixed,setprecision \\ \hline
string & 字符串 & string \\ \hline
cstring & C字符串 & strlen,memset \\ \hline
sstream & 字符串转换 & stringstream \\ \hline
cstdlib & 工具函数 & rand,exit \\ \hline
ctime & 时间 & time \\ \hline
limits & 极值 & INT\_MAX \\ \hline
numeric & 数学算法 & accumulate,iota \\ \hline
algorithm & 算法 & sort,find \\ \hline
vector & 动态数组 & vector \\ \hline
queue & 队列 & queue \\ \hline
stack & 栈 & stack \\ \hline
map & 映射 & map \\ \hline
set & 集合 & set \\ \hline
\end{tabular}
\end{table}

\noindent\rule{\linewidth}{0.4pt}

\section{第十一章 C++ 常用语法速查（机考必备）}

这一章不是 STL，而是\textbf{写题时每天都会用到的 C++ 基础语法工具}。

重点：
\begin{itemize}
    \item 输入输出
    \item 类型转换
    \item 字符处理
    \item ASCII
    \item 位运算
    \item long long
    \item 常用技巧
    \item 防止错误
\end{itemize}

\noindent\rule{\linewidth}{0.4pt}

\section{一、输入输出（iostream）}

头文件：
\begin{lstlisting}[language=C++]
#include <iostream>
\end{lstlisting}

\noindent\rule{\linewidth}{0.4pt}

\section{1. cin 输入}

\subsection{用法}
\begin{lstlisting}[language=C++]
cin >> x;
\end{lstlisting}
例如：
\begin{lstlisting}[language=C++]
int a;

cin >> a;
\end{lstlisting}
作用：\\
读取变量。

\noindent\rule{\linewidth}{0.4pt}

\section{2. 多个变量输入}
\begin{lstlisting}[language=C++]
int a,b,c;

cin>>a>>b>>c;
\end{lstlisting}
输入：
\begin{lstlisting}
1 2 3
\end{lstlisting}
结果：
\begin{lstlisting}
a=1,b=2,c=3
\end{lstlisting}

\noindent\rule{\linewidth}{0.4pt}

\section{3. cout 输出}

\subsection{用法}
\begin{lstlisting}[language=C++]
cout << x;
\end{lstlisting}
例如：
\begin{lstlisting}[language=C++]
cout<<123;
\end{lstlisting}
输出：
\begin{lstlisting}
123
\end{lstlisting}

\noindent\rule{\linewidth}{0.4pt}

\section{4. 输出多个内容}
\begin{lstlisting}[language=C++]
cout<<"sum="<<sum;
\end{lstlisting}

\noindent\rule{\linewidth}{0.4pt}

\section{5. 换行}

\subsection{方法1}
\begin{lstlisting}[language=C++]
cout<<endl;
\end{lstlisting}
作用：\\
换行，并刷新缓冲区。

\noindent\rule{\linewidth}{0.4pt}

\subsection{方法2（推荐）}
\begin{lstlisting}[language=C++]
cout<<"\n";
\end{lstlisting}
速度更快。

考试推荐：
\begin{lstlisting}[language=C++]
cout<<"\n";
\end{lstlisting}

\noindent\rule{\linewidth}{0.4pt}

\section{6. getline()}

\textbf{【必背】}

\subsection{用法}
\begin{lstlisting}[language=C++]
string s;

getline(cin,s);
\end{lstlisting}
输入：
\begin{lstlisting}
hello world
\end{lstlisting}
读取：
\begin{lstlisting}
hello world
\end{lstlisting}

\noindent\rule{\linewidth}{0.4pt}

\subsection{注意}
如果前面用了：
\begin{lstlisting}[language=C++]
cin>>n;
\end{lstlisting}
后面：
\begin{lstlisting}[language=C++]
getline(cin,s);
\end{lstlisting}
可能读到空行。

解决：
\begin{lstlisting}[language=C++]
cin.ignore();
\end{lstlisting}
例如：
\begin{lstlisting}[language=C++]
int n;

cin>>n;

cin.ignore();

getline(cin,s);
\end{lstlisting}

\noindent\rule{\linewidth}{0.4pt}

\section{二、类型转换（非常重要）}

\noindent\rule{\linewidth}{0.4pt}

\section{1. int 转 double}

自动转换：
\begin{lstlisting}[language=C++]
int a=5;

double b=a;
\end{lstlisting}
结果：
\begin{lstlisting}
5.0
\end{lstlisting}

\noindent\rule{\linewidth}{0.4pt}

\section{2. 强制类型转换}

\subsection{写法1}
\begin{lstlisting}[language=C++]
(double)a
\end{lstlisting}
例如：
\begin{lstlisting}[language=C++]
int a=5,b=2;


cout<<(double)a/b;
\end{lstlisting}
结果：
\begin{lstlisting}
2.5
\end{lstlisting}

\noindent\rule{\linewidth}{0.4pt}

\subsection{写法2（C++推荐）}
\begin{lstlisting}[language=C++]
static_cast<double>(a)
\end{lstlisting}
例如：
\begin{lstlisting}[language=C++]
double x=
static_cast<double>(a)/b;
\end{lstlisting}

\noindent\rule{\linewidth}{0.4pt}

\section{3. string 转 int}

\subsection{stoi()}
头文件：
\begin{lstlisting}[language=C++]
#include <string>
\end{lstlisting}
用法：
\begin{lstlisting}[language=C++]
string s="123";


int x=stoi(s);
\end{lstlisting}
结果：
\begin{lstlisting}
123
\end{lstlisting}

\noindent\rule{\linewidth}{0.4pt}

\section{4. string 转 long long}
\begin{lstlisting}[language=C++]
long long x=stoll(s);
\end{lstlisting}

\noindent\rule{\linewidth}{0.4pt}

\section{5. string 转 double}
\begin{lstlisting}[language=C++]
double x=stod(s);
\end{lstlisting}

\noindent\rule{\linewidth}{0.4pt}

\section{6. int 转 string}

\subsection{to\_string()}
用法：
\begin{lstlisting}[language=C++]
string s=to_string(123);
\end{lstlisting}
结果：
\begin{lstlisting}
"123"
\end{lstlisting}

\noindent\rule{\linewidth}{0.4pt}

\section{三、字符处理（cctype）}

头文件：
\begin{lstlisting}[language=C++]
#include <cctype>
\end{lstlisting}

\noindent\rule{\linewidth}{0.4pt}

\section{1. isdigit()}

\subsection{名称}
判断数字字符

\subsection{用法}
\begin{lstlisting}[language=C++]
isdigit(c);
\end{lstlisting}
例如：
\begin{lstlisting}[language=C++]
char c='5';


if(isdigit(c))
{
    cout<<"数字";
}
\end{lstlisting}
作用：\\
判断：
\begin{lstlisting}
0~9
\end{lstlisting}

\noindent\rule{\linewidth}{0.4pt}

\section{2. isalpha()}

判断字母。
\begin{lstlisting}[language=C++]
isalpha(c);
\end{lstlisting}
例如：
\begin{lstlisting}[language=C++]
'a'
\end{lstlisting}
返回：\\
true。

\noindent\rule{\linewidth}{0.4pt}

\section{3. islower()}

判断小写。
\begin{lstlisting}[language=C++]
islower(c);
\end{lstlisting}

\noindent\rule{\linewidth}{0.4pt}

\section{4. isupper()}

判断大写。
\begin{lstlisting}[language=C++]
isupper(c);
\end{lstlisting}

\noindent\rule{\linewidth}{0.4pt}

\section{5. tolower()}

转换小写。
\begin{lstlisting}[language=C++]
tolower(c);
\end{lstlisting}
例如：
\begin{lstlisting}[language=C++]
char c='A';

c=tolower(c);
\end{lstlisting}
结果：
\begin{lstlisting}
a
\end{lstlisting}

\noindent\rule{\linewidth}{0.4pt}

\section{6. toupper()}

转换大写。
\begin{lstlisting}[language=C++]
toupper(c);
\end{lstlisting}
例如：
\begin{lstlisting}
a→A
\end{lstlisting}

\noindent\rule{\linewidth}{0.4pt}

\section{四、ASCII码（高频）}

字符本质：\\
数字。

\noindent\rule{\linewidth}{0.4pt}

\section{1. 字符转数字}

例如：
\begin{lstlisting}[language=C++]
char c='5';


int x=c-'0';
\end{lstlisting}
结果：
\begin{lstlisting}
5
\end{lstlisting}
原理：\\
ASCII：
\begin{lstlisting}
'5'=53

'0'=48

53-48=5
\end{lstlisting}

\noindent\rule{\linewidth}{0.4pt}

\section{2. 数字转字符}
\begin{lstlisting}[language=C++]
char c=x+'0';
\end{lstlisting}
例如：
\begin{lstlisting}[language=C++]
int x=7;

char c=x+'0';
\end{lstlisting}
结果：
\begin{lstlisting}
'7'
\end{lstlisting}

\noindent\rule{\linewidth}{0.4pt}

\section{3. 字母大小写转换}

大写：
\begin{lstlisting}[language=C++]
char c='A';

c=c+32;
\end{lstlisting}
得到：
\begin{lstlisting}
a
\end{lstlisting}

\noindent\rule{\linewidth}{0.4pt}

小写转大写：
\begin{lstlisting}[language=C++]
c=c-32;
\end{lstlisting}
不过推荐：
\begin{lstlisting}[language=C++]
toupper()
tolower()
\end{lstlisting}

\noindent\rule{\linewidth}{0.4pt}

\section{五、long long使用（非常重要）}

\noindent\rule{\linewidth}{0.4pt}

\section{1. int范围}
\begin{lstlisting}
-2^31 ~ 2^31-1
\end{lstlisting}
约：
\begin{lstlisting}
21亿
\end{lstlisting}

\noindent\rule{\linewidth}{0.4pt}

\section{2. long long范围}
约：
\begin{lstlisting}
9×10^18
\end{lstlisting}

\noindent\rule{\linewidth}{0.4pt}

\section{3. 定义}
\begin{lstlisting}[language=C++]
long long x;
\end{lstlisting}

\noindent\rule{\linewidth}{0.4pt}

\section{4. 常量}
错误：
\begin{lstlisting}[language=C++]
long long x=1000000000000;
\end{lstlisting}
可能溢出。

正确：
\begin{lstlisting}[language=C++]
long long x=1000000000000LL;
\end{lstlisting}

\noindent\rule{\linewidth}{0.4pt}

\section{六、三目运算符}

\subsection{名称}
条件运算符

\noindent\rule{\linewidth}{0.4pt}

\subsection{格式}
\begin{lstlisting}[language=C++]
条件 ? 真 : 假;
\end{lstlisting}

\noindent\rule{\linewidth}{0.4pt}

例如：
\begin{lstlisting}[language=C++]
int max=
a>b?a:b;
\end{lstlisting}
等价：
\begin{lstlisting}[language=C++]
if(a>b)
max=a;

else
max=b;
\end{lstlisting}

\noindent\rule{\linewidth}{0.4pt}

\section{七、位运算（了解）}

符号：
\begin{lstlisting}
&
|
^
~
<<
>>
\end{lstlisting}

\noindent\rule{\linewidth}{0.4pt}

\section{1. 按位与 \&}

例如：
\begin{lstlisting}[language=C++]
5&3
\end{lstlisting}
二进制：
\begin{lstlisting}
101
011
---
001
\end{lstlisting}
结果：
\begin{lstlisting}
1
\end{lstlisting}
作用：\\
判断奇偶。

\noindent\rule{\linewidth}{0.4pt}

\section{判断奇偶\textbf{【常用】}}

偶数：
\begin{lstlisting}[language=C++]
if((x&1)==0)
\end{lstlisting}
奇数：
\begin{lstlisting}[language=C++]
if(x&1)
\end{lstlisting}

\noindent\rule{\linewidth}{0.4pt}

\section{2. 按位或 |}
\begin{lstlisting}[language=C++]
a|b
\end{lstlisting}
作用：\\
合并二进制位。

\noindent\rule{\linewidth}{0.4pt}

\section{3. 异或 \textasciicircum{}}

\subsection{特点}
相同：
\begin{lstlisting}
0
\end{lstlisting}
不同：
\begin{lstlisting}
1
\end{lstlisting}

\noindent\rule{\linewidth}{0.4pt}

应用：\\
交换变量：
\begin{lstlisting}[language=C++]
a^=b;

b^=a;

a^=b;
\end{lstlisting}

\noindent\rule{\linewidth}{0.4pt}

\section{4. 左移 <<}
\begin{lstlisting}[language=C++]
x<<1
\end{lstlisting}
相当：
\begin{lstlisting}
x*2
\end{lstlisting}
例如：
\begin{lstlisting}[language=C++]
5<<1
\end{lstlisting}
结果：
\begin{lstlisting}
10
\end{lstlisting}

\noindent\rule{\linewidth}{0.4pt}

\section{5. 右移 >>}
\begin{lstlisting}[language=C++]
x>>1
\end{lstlisting}
相当：
\begin{lstlisting}
x/2
\end{lstlisting}

\noindent\rule{\linewidth}{0.4pt}

\section{八、常用数学技巧}

\noindent\rule{\linewidth}{0.4pt}

\section{1. 最大公约数 gcd}

头文件：
\begin{lstlisting}[language=C++]
#include <numeric>
\end{lstlisting}
用法：
\begin{lstlisting}[language=C++]
gcd(a,b);
\end{lstlisting}
例如：
\begin{lstlisting}[language=C++]
cout<<gcd(12,8);
\end{lstlisting}
输出：
\begin{lstlisting}
4
\end{lstlisting}

\noindent\rule{\linewidth}{0.4pt}

\section{2. 最小公倍数 lcm}

C++17：
\begin{lstlisting}[language=C++]
lcm(a,b);
\end{lstlisting}
例如：
\begin{lstlisting}[language=C++]
cout<<lcm(4,6);
\end{lstlisting}
输出：
\begin{lstlisting}
12
\end{lstlisting}

\noindent\rule{\linewidth}{0.4pt}

\section{3. 判断素数模板}
\begin{lstlisting}[language=C++]
bool isPrime(int n)
{
    if(n<2)
        return false;


    for(int i=2;i*i<=n;i++)
    {
        if(n%i==0)
            return false;
    }


    return true;
}
\end{lstlisting}

\noindent\rule{\linewidth}{0.4pt}

\section{九、swap交换}

\subsection{名称}
\begin{lstlisting}[language=C++]
swap()
\end{lstlisting}
头文件：
\begin{lstlisting}[language=C++]
#include <algorithm>
\end{lstlisting}

\noindent\rule{\linewidth}{0.4pt}

\subsection{用法}
\begin{lstlisting}[language=C++]
swap(a,b);
\end{lstlisting}
例如：
\begin{lstlisting}[language=C++]
int a=1,b=2;

swap(a,b);
\end{lstlisting}
结果：
\begin{lstlisting}
a=2,b=1
\end{lstlisting}

\noindent\rule{\linewidth}{0.4pt}

\section{十、auto自动类型推导}

\noindent\rule{\linewidth}{0.4pt}

\subsection{用法}
\begin{lstlisting}[language=C++]
auto x=10;
\end{lstlisting}
等价：
\begin{lstlisting}[language=C++]
int x=10;
\end{lstlisting}

\noindent\rule{\linewidth}{0.4pt}

vector遍历：
\begin{lstlisting}[language=C++]
for(auto x:v)
{
}
\end{lstlisting}
map：
\begin{lstlisting}[language=C++]
for(auto p:mp)
{
}
\end{lstlisting}

\noindent\rule{\linewidth}{0.4pt}

\section{十一、引用 \&（STL常用）}

\noindent\rule{\linewidth}{0.4pt}

\subsection{普通变量}
\begin{lstlisting}[language=C++]
int a=10;


int &b=a;
\end{lstlisting}
b就是a别名。

\noindent\rule{\linewidth}{0.4pt}

\subsection{修改vector元素}
错误：
\begin{lstlisting}[language=C++]
for(auto x:v)
{
    x++;
}
\end{lstlisting}
不会改变。

\noindent\rule{\linewidth}{0.4pt}

正确：
\begin{lstlisting}[language=C++]
for(auto &x:v)
{
    x++;
}
\end{lstlisting}

\noindent\rule{\linewidth}{0.4pt}

\section{十二、const（防止修改）}

\noindent\rule{\linewidth}{0.4pt}

定义：
\begin{lstlisting}[language=C++]
const int x=10;
\end{lstlisting}
不能：
\begin{lstlisting}[language=C++]
x=20;
\end{lstlisting}

\noindent\rule{\linewidth}{0.4pt}

常见：
\begin{lstlisting}[language=C++]
const vector<int>& v
\end{lstlisting}
表示：\\
引用vector，但不能修改。

\noindent\rule{\linewidth}{0.4pt}

\section{十三、宏定义（了解）}

\noindent\rule{\linewidth}{0.4pt}

\section{1. 常量}
\begin{lstlisting}[language=C++]
#define PI 3.14159
\end{lstlisting}

\noindent\rule{\linewidth}{0.4pt}

\section{2. 宏函数}
\begin{lstlisting}[language=C++]
#define MAX(a,b) ((a)>(b)?(a):(b))
\end{lstlisting}
现在C++更推荐：
\begin{lstlisting}[language=C++]
const
inline
\end{lstlisting}

\noindent\rule{\linewidth}{0.4pt}

\section{十四、结构体 struct}

考试经常用。

\noindent\rule{\linewidth}{0.4pt}

定义：
\begin{lstlisting}[language=C++]
struct Student
{
    int id;

    string name;

};
\end{lstlisting}

\noindent\rule{\linewidth}{0.4pt}

使用：
\begin{lstlisting}[language=C++]
Student s;


s.id=1;

s.name="Tom";
\end{lstlisting}

\noindent\rule{\linewidth}{0.4pt}

数组：
\begin{lstlisting}[language=C++]
Student a[100];
\end{lstlisting}

\noindent\rule{\linewidth}{0.4pt}

vector：
\begin{lstlisting}[language=C++]
vector<Student> v;
\end{lstlisting}

\noindent\rule{\linewidth}{0.4pt}

\section{十五、自定义排序（结构体）}

例如：
\begin{lstlisting}[language=C++]
struct Node
{
    int x;

    int y;
};


bool cmp(Node a,Node b)
{
    return a.x<b.x;
}
\end{lstlisting}
使用：
\begin{lstlisting}[language=C++]
sort(v.begin(),v.end(),cmp);
\end{lstlisting}

\noindent\rule{\linewidth}{0.4pt}

\section{十六、万能main模板（考试推荐）}

\begin{lstlisting}[language=C++]
#include <bits/stdc++.h>

using namespace std;


int main()
{
    ios::sync_with_stdio(false);
    cin.tie(nullptr);


    

    return 0;
}
\end{lstlisting}

\noindent\rule{\linewidth}{0.4pt}

\section{十七、考试最常用组合}

\subsection{1. 输入数组排序求和}
\begin{lstlisting}[language=C++]
vector<int> v(n);


for(auto &x:v)
    cin>>x;


sort(v.begin(),v.end());


int sum=
accumulate(v.begin(),v.end(),0);
\end{lstlisting}

\noindent\rule{\linewidth}{0.4pt}

\subsection{2. 统计字符次数}
\begin{lstlisting}[language=C++]
map<char,int> mp;


for(char c:s)
{
    mp[c]++;
}
\end{lstlisting}

\noindent\rule{\linewidth}{0.4pt}

\subsection{3. 删除重复元素}
\begin{lstlisting}[language=C++]
sort(v.begin(),v.end());


v.erase(
unique(v.begin(),v.end()),
v.end()
);
\end{lstlisting}

\noindent\rule{\linewidth}{0.4pt}

\subsection{4. 保留两位小数}
\begin{lstlisting}[language=C++]
cout
<<fixed
<<setprecision(2)
<<x;
\end{lstlisting}

\noindent\rule{\linewidth}{0.4pt}

\subsection{5. 最大最小值}
\begin{lstlisting}[language=C++]
int mx=*max_element(
v.begin(),
v.end()
);


int mn=*min_element(
v.begin(),
v.end()
);
\end{lstlisting}

\noindent\rule{\linewidth}{0.4pt}

\section{最终考试头文件优先级}

\textbf{【必背】}
\begin{lstlisting}[language=C++]
#include <bits/stdc++.h>
\end{lstlisting}
里面全部包含。

重点掌握：
\begin{lstlisting}
vector
string
algorithm
numeric
cmath
iomanip
map
set
queue
stack
sstream
cstring
\end{lstlisting}

\noindent\rule{\linewidth}{0.4pt}

\end{document}